\documentclass[11pt]{amsart}
\usepackage{amsmath,amssymb,amsthm,mathtools,mathrsfs}
\usepackage[margin=1.08in]{geometry}
\usepackage{booktabs}
\usepackage{array}
\usepackage{longtable}
\usepackage{tabularx}
\usepackage[hypertexnames=false]{hyperref}
\hypersetup{hidelinks}
\usepackage[final,expansion=false]{microtype}
\usepackage{url}
\setlength{\parskip}{4pt plus 1pt}
\emergencystretch=2.5em

\newtheorem{theorem}{Theorem}[section]
\newtheorem{proposition}[theorem]{Proposition}
\newtheorem{lemma}[theorem]{Lemma}
\newtheorem{corollary}[theorem]{Corollary}
\theoremstyle{definition}
\newtheorem{definition}[theorem]{Definition}
\newtheorem{problem}[theorem]{Problem}
\theoremstyle{remark}
\newtheorem{remark}[theorem]{Remark}

\newcommand{\C}{\mathbb C}
\newcommand{\R}{\mathbb R}
\newcommand{\N}{\mathbb N}
\newcommand{\Z}{\mathbb Z}
\newcommand{\Q}{\mathbb Q}
\newcommand{\Rel}{\operatorname{Re}}
\newcommand{\Iml}{\operatorname{Im}}
\newcommand{\supp}{\operatorname{supp}}
\newcommand{\TV}{\operatorname{TV}}
\newcommand{\Kp}{\mathcal K^{\mathrm{prim}}_{2T}}
\newcommand{\Kw}{\mathcal K_{2T}}
\newcommand{\PT}{\mathcal P_T}
\newcommand{\IT}{I_T}
\newcommand{\lamP}{\lambda^{\mathrm{prim}}_T}
\newcommand{\lamW}{\lambda^{\mathrm W}_T}
\newcommand{\lamA}{\lambda_{\mathrm A}(T)}
\newcommand{\sT}{\sigma_T}
\newcommand{\gh}[1]{\widehat g_{#1}}
\newcommand{\nuh}{\widehat\nu}
\newcommand{\wh}{\widehat w}
\newcommand{\Fh}{\widehat F}
\newcommand{\cF}{c_F}
\newcommand{\Lev}{\mathfrak E}
\newcommand{\Wcl}{\mathcal W_T}
\newcommand{\NT}{\mathcal N_T}
\newcommand{\Aconst}{\mathsf A_T}
\newcommand{\Ainf}{\mathcal A_\infty}
\newcommand{\PiT}{\Pi_T}
\newcommand{\mT}{\mathfrak m_T}
\newcommand{\laminf}{\lambda^{\infty}_{\min}}
\newcommand{\laminfP}{\lambda^{\infty,\mathrm P}_{\min}}
\newcolumntype{Y}{>{\raggedright\arraybackslash}X}

\begin{document}
\title[The sign of the radial intersection form]{The sign of the radial intersection form on the idele class group}
\author{David Alejandro Trejo Pizzo}
\date{}
\begin{abstract}
The radial trace on the idele class group of $\mathbb Q$ defines a Hermitian
form on functions supported in $(-T,T)$. We determine which sign of this form
corresponds to the intersection inequality on the kernel of the two Tate
functionals. If $\zeta(s)$ has no zero to the right of the critical line,
we prove $W_T(F)\ge0$ on the entire form domain for every $T>0$.
On the arithmetic side we reduce this all-window inequality to positive
semidefiniteness in every degree of a single sequence of explicit matrices.
On bounded windows we also obtain quantitative estimates: for
$T\le\tfrac12\log5$ the lower bounds on unit vectors are $10^{-18}$ for the
complete form and $2.5\cdot10^{-12}$ for its primitive restriction.
A certificate using finite-height zero data gives an explicit positive
lower bound up to $T=2.8022$.
\end{abstract}
\maketitle
\setcounter{tocdepth}{1}\tableofcontents
\section{Introduction}\label{sec:intro}
In Weil's intersection argument for a curve over a finite field, one removes
the two degree directions before applying the index inequality to a correspondence.
The separation matters: the two directions pair hyperbolically, so their
contribution to self-intersection has no fixed sign. For the correspondence
algebra attached to $\operatorname{Spec}\mathbb Z$, the constructions of \cite{P43,P44,Paper45} identify
the two degrees and compute the radial self-intersection. The fourth step,
which we call Row~D, asks whether the restriction to their common kernel
has the required sign at every support radius.

The diagonal contact of the arithmetic correspondence of degree $n$ has
weight $\Lambda(n)/\sqrt n$. The same coefficient occurs in the explicit
formula for the radial trace. The question here is what those contacts,
taken together with the archimedean place, imply for the self-intersection
of a class of bidegree zero.

We work on the carrier of the preceding papers: the semilocal idele class
groups for finite sets of places of $\mathbb Q$ containing infinity, with
the radial trace taken in their colimit. The involution, trace and Tate functionals
are constructed in \cite{P43}. The diagonal contact is computed in \cite{P44},
and the radial trace is expressed as a sum over the zeros with a polar
correction in \cite{Paper45}.
A compact support interval meets only finitely many prime-power contacts.
Enlarging the interval brings further primes into the same form.

There are two questions about the sign. At a fixed radius one can seek a
positive numerical lower bound for unit vectors. For arbitrary support the
intersection inequality requires only nonnegativity, but it requires this
for every radius. We treat the first question by bounds for the Fourier
multiplier, and the second by testing a sequence of finite matrices. The
conclusion returns to the zero sum and gives the precise logical relation
between its sign and the position of the zeros.

\section{The radial trace}\label{sec:trace}
Write $I_T=(-T,T)$ and extend functions by zero. Our two transforms are
\[
 \widehat F(\tau)=\int F(t)e^{i\tau t}\,dt,\qquad
 \Phi_F(s)=\int F(t)e^{(s-1/2)t}\,dt .
\]
Thus $\Phi_F(1)=A(F)$ and $\Phi_F(0)=B(F)$ are the Tate moments.
A class is primitive when both vanish. Put
$c_F(u)=\int F(t)\overline{F(t-u)}\,dt$.
For a smooth real test, the comparison theorem \cite[Theorem~C]{Paper45} gives
\begin{equation}\label{comparison}
 I_\infty(\Gamma_F,\Gamma_F)
 =2A(F)B(F)-\sum_\rho\Phi_F(\rho)\Phi_F(1-\rho).
\end{equation}
Zeros are counted with multiplicity. The sum converges absolutely: repeated
integration by parts gives rapid decay in the ordinate uniformly across
the critical strip, and the zero count up to height $H$ is $O(H\log H)$.
The primitive intersection inequality is $I_\infty\le0$. This sign
convention accounts for the minus sign in the next formula.


There is a suggestive way of seeing both degree directions before passing
to the Fourier transform. On real tests let $M(F)=(A(F),B(F))$ and give
$\mathbb R^2$ the symmetric pairing
\[
 H((a,b),(a',b'))=ab'+ba'.
\]
The adjoint on the radial variable is reflection. Since
$\Phi_{F^\dagger}(s)=\Phi_F(1-s)$ for real $F$, it interchanges the two
coordinates of $M$. Under convolution $\Phi_{F*G}(s)=\Phi_F(s)\Phi_G(s)$.
Thus the polar contribution on a product $F*G^\dagger$ is exactly
$H(M(F),M(G))$. To check that the two conditions are independent in
every window, take a smooth nonnegative even bump near the origin and
two small translates by distinct amounts $r_1,r_2$. Its moment vector
at $r$ is a positive constant times $(e^{r/2},e^{-r/2})$. The determinant
of the two vectors is proportional to $2\sinh((r_1-r_2)/2)$.
This identifies the primitive space as a codimension-two kernel even
before completion.

For $F,G$ real, polarization of \eqref{comparison} gives
\[
 I_\infty(\Gamma_F,\Gamma_G)
 = H(M(F),M(G))-
 \tfrac12\sum_\rho
 \{\Phi_F(\rho)\Phi_G(1-\rho)+\Phi_G(\rho)\Phi_F(1-\rho)\}.
\]
In particular a positive direction of $I_\infty$ in the common kernel is
exactly a negative direction for the zero sum there. This is where the
sign enters the intersection argument. Computing the bidegree and
identifying the local coefficients fix the pairing, but neither operation
determines its signature. On each window the rest of the argument studies
that signature with the arithmetic coefficients left in place.

For the arithmetic side of the comparison we use the following normalization.
Throughout, $T>0$, $\IT=(-T,T)$, and for $F\in L^2(\IT)$ (extended by zero) we write
$\Fh(\tau)=\int F(t)e^{i\tau t}\,dt$ and $\cF(u)=\int F(t)\overline{F(t-u)}\,dt$ for the
autocorrelation, so that $\cF$ is continuous, supported in $[-2T,2T]$, and
$\widehat{\cF}=|\Fh|^2$ on $\R$. Put
\begin{equation}\label{eq:symbol}
m(\tau)=\Rel\psi\!\left(\tfrac14+\tfrac{i\tau}{2}\right)-\log\pi,\qquad
\sT(\tau)=m(\tau)-\sum_{n\in\NT}v_n\cos(\tau\log n),\qquad v_n=\frac{2\Lambda(n)}{\sqrt n},
\end{equation}
where $\psi$ is the digamma function, $\NT=\{n\ge2:\ \log n<2T\}$ is the set of \emph{contacts} and
$\Aconst:=\sum_{n\in\NT}v_n$ is the total mass of the prime comb. The truncated Weil form is
\begin{equation}\label{eq:form}
\mathfrak a_T(F)=\frac1{2\pi}\int_\R\sT(\tau)\,|\Fh(\tau)|^2\,d\tau
\end{equation}
on the form domain $\mathcal D_T=\{F\in L^2(\IT):\int\log(2+|\tau|)|\Fh|^2<\infty\}$. The two \emph{Tate
moments} are $A(F)=\int F(t)e^{t/2}dt$ and $B(F)=\int F(t)e^{-t/2}dt$, the \emph{primitive subspace} is
$\PT=\{F\in L^2(\IT):A(F)=B(F)=0\}$, and the \emph{complete Weil form} on the window is
\begin{equation}\label{eq:W}
W_T(F)=\mathfrak a_T(F)+2\,\Rel\big(A(F)\overline{B(F)}\big).
\end{equation}
In particular the arithmetic side gives $I_\infty(\Gamma_F,\Gamma_F)=-\mathfrak a_T(F)$ for real smooth tests. By the explicit formula (\cite{Weil52}, \cite[Thm.~8.4]{Paper45}; Lemma~\ref{lem:roughEF} below for
tests that are not smooth), for real smooth tests $W_T(F)=\sum_\rho\Phi_F(\rho)\Phi_F(1-\rho)$ with
$\Phi_F(s)=\int F(t)e^{(s-\frac12)t}dt$, the sum over the nontrivial zeros of $\zeta$. We write
\[
\lamA=\inf\{\mathfrak a_T(F):F\in\PT\cap\mathcal D_T,\ \|F\|=1\},\qquad
\lamW=\inf\{W_T(F):F\in\mathcal D_T,\ \|F\|=1\};
\]
$\lamW\le\lamA$ because the two forms agree on $\PT$. The object of this paper is the following
statement, which we call \emph{window positivity}:
\begin{equation}\label{eq:target}
\lamA\ \ge\ 0\qquad\text{for every }T>0 ,
\end{equation}
together with its complete version $\lamW\ge0$ for every $T>0$. Since $T\mapsto\lamA$ is nonincreasing
(Lemma~\ref{lem:monoT}), \eqref{eq:target} is equivalent to $\lim_{T\to\infty}\lamA\ge0$; it is a
statement about every window at once, and every window carries a finite amount of arithmetic data, the
vector $(\Lambda(n))_{n<e^{2T}}$.
\subsection{Conventions}\label{subsec:conv}
$\widehat h(\tau)=\int_\R h(u)e^{i\tau u}du$; for even $h$ this is $\int h\cos(\tau u)$. For a finite even
measure $\nu$, $\nuh(\tau)=\int\cos(\tau u)\,d\nu(u)$, and for a probability measure $\mu$ on $\R$,
$\check\mu(u)=\int\cos(\tau u)\,d\mu(\tau)$. The letter $\tau$ always denotes a frequency, and $u$ a
variable dual to it. The right shift is $(V_aF)(t)=F(t-a)$, the reflection is $(JF)(t)=F(-t)$, and $Q_T$
denotes restriction to $\IT$ followed by extension by zero. We write $\mT(F)=\frac1{2\pi}\int m|\Fh|^2$
for the archimedean form, $\PiT(F)=-\sum_{n\in\NT}v_n\Rel\cF(\log n)$ for the prime form, so that
$\mathfrak a_T=\mT+\PiT$, and $\Ainf=\mT+2\Rel(A\bar B)$ for the complete archimedean form. The functions
\begin{equation}\label{eq:g}
g_0(u)=\cosh\tfrac u2\,\mathbf 1_{|u|<2T},\qquad g_1(u)=u\sinh\tfrac u2\,\mathbf 1_{|u|<2T},
\end{equation}
have transforms $\gh0,\gh1$ (Lemma~\ref{lem:gbounds}). For a parity label
$\mathrm{par}\in\{\mathrm{ev},\mathrm{odd}\}$, $\chi_{\mathrm{ev}}=\cosh\frac t2\,\mathbf 1_{\IT}$ and
$\chi_{\mathrm{odd}}=\sinh\frac t2\,\mathbf 1_{\IT}$. Letters not listed here are local to the proof or
the subsection in which they occur.

The Fourier form studied by Yoshida and Bombieri gives the variational setting for this question
\cite{Yoshida1992,Bombieri2000}. Connes and Consani obtain positivity in the
archimedean regime through the semilocal trace formula \cite{ConnesConsani2021}.
Here the coefficients of the prime translations are fixed by the diagonal
contacts, and both polar directions survive until the primitive restriction is made.
We first study the multiplier contributed by the real place. We then use
autocorrelations to construct dual lower bounds and to identify their relation
to Weil's explicit formula \cite{Weil52}. Finally we separate the polar directions
and reduce the primitive sign at a fixed radius to a finite Schur complement.
\section{The archimedean multiplier}\label{sec:levy}

\begin{lemma}[the L\'evy density]\label{lem:K}
For $u>0$ put $K(u)=\dfrac{2e^{3u/2}}{e^{2u}-1}=\dfrac{1}{2\sinh(u/2)}+\dfrac{1}{2\cosh(u/2)}
=2\sum_{k\ge0}e^{-(2k+\frac12)u}$. Then $K$ is positive, completely monotone,
$K(u)=u^{-1}+O(1)$ as $u\downarrow0$, $K(u)=2e^{-u/2}(1+O(e^{-u}))$ as $u\to\infty$, and
\begin{equation}\label{eq:m}
\begin{split}
m(\tau)&=\Rel\psi\!\left(\tfrac14+\tfrac{i\tau}2\right)-\log\pi
= C_0+\int_0^\infty(1-\cos\tau u)\,K(u)\,du,\\
C_0&=\psi(\tfrac14)-\log\pi=-5.3721834\ldots
\end{split}
\end{equation}
\end{lemma}
\begin{proof}
The three expressions for $K$ agree: $\frac{1}{2\sinh(u/2)}+\frac1{2\cosh(u/2)}
=\frac{\cosh(u/2)+\sinh(u/2)}{\sinh u}=\frac{2e^{u/2}}{e^u-e^{-u}}=\frac{2e^{3u/2}}{e^{2u}-1}$, and
$\frac{2e^{3u/2}}{e^{2u}-1}=2e^{-u/2}\sum_{k\ge0}e^{-2ku}$. Complete monotonicity is that of a
convergent series of decaying exponentials with positive coefficients. For \eqref{eq:m} use
$\psi(z)=-\gamma+\int_0^\infty\frac{e^{-t}-e^{-zt}}{1-e^{-t}}dt$ ($\Rel z>0$): with $z=\frac14+\frac{i\tau}2$,
\[
\Rel\psi(z)-\psi(\tfrac14)=\int_0^\infty\frac{e^{-t/4}\big(1-\cos(\tau t/2)\big)}{1-e^{-t}}\,dt
=\int_0^\infty\frac{2e^{-u/2}(1-\cos\tau u)}{1-e^{-2u}}\,du\qquad(t=2u),
\]
and $\frac{2e^{-u/2}}{1-e^{-2u}}=\frac{2e^{3u/2}}{e^{2u}-1}=K(u)$. The integral converges since
$1-\cos\tau u=O(u^2)$ at $0$ and $K=O(e^{-u/2})$ at $\infty$.
\end{proof}

\begin{theorem}[complete Bernstein structure]\label{thm:S}
Put $a_k=k+\frac14$ and $\varphi(x)=\sum_{k\ge0}\dfrac{x}{a_k(x+4a_k^2)}$ for $x\ge0$. Then
\emph{(i)} $\Rel\psi(\frac14+\frac{i\tau}2)-\psi(\frac14)=\varphi(\tau^2)$ and $\varphi$ is a
complete Bernstein function; \emph{(ii)} the L\'evy measure of the negative-definite function
$\tau\mapsto\varphi(\tau^2)$ on $\R$ is $\frac12K(|u|)du$, which is completely monotone on each
half-line, so the associated L\'evy process is a pure-jump subordinate Brownian motion
$B_{S_t}$ with a special subordinator $S$ of Laplace exponent $\varphi$; in particular it has no
Gaussian component; \emph{(iii)} $\int_{2T}^\infty K(u)\,du=\log\coth\frac T2+\pi-2\arctan e^{T}$.
\end{theorem}
\begin{proof}
(i) From $\psi(z)=-\gamma+\sum_{k\ge0}\big[\frac1{k+1}-\frac1{k+z}\big]$,
$\Rel\big[\frac1{a_k}-\frac1{a_k+i\tau/2}\big]=\frac1{a_k}-\frac{a_k}{a_k^2+\tau^2/4}
=\frac{\tau^2}{a_k(4a_k^2+\tau^2)}$, which is the stated series at $x=\tau^2$. Each summand
$x\mapsto a_k^{-1}\,x/(x+4a_k^2)$ is a complete Bernstein function (Stieltjes representation with
a point mass at $4a_k^2$), the summands are $O(k^{-3})$ locally uniformly, and the complete
Bernstein functions form a convex cone closed under pointwise limits
\cite[Ch.~6--7]{SSV2012}. (ii) By Lemma~\ref{lem:K}, $\varphi(\tau^2)=\int_\R(1-\cos\tau u)\,
\frac12K(|u|)du$, so the L\'evy measure is $\frac12K(|u|)du$; it is completely monotone on
$(0,\infty)$ and $\varphi(x)/x\to0$, so there is no Gaussian part; the representation as a
subordinate Brownian motion with a special subordinator is \cite[Thm.~13.10, Cor.~13.12]{SSV2012}.
(iii) $\int\frac{du}{2\sinh(u/2)}=\log\tanh\frac u4$ and $\int\frac{du}{2\cosh(u/2)}=2\arctan e^{u/2}$.
\end{proof}

\begin{theorem}[the form as a L\'evy energy]\label{thm:A}
For every $F\in L^2(\IT)$, with $\Lev(\Phi):=\int_0^\infty(\Phi(0)-\Rel\Phi(u))K(u)\,du\in[0,\infty]$
(for complex $F$ the autocorrelation is Hermitian and only its even part $\Rel\cF$ enters; the
prime term below is likewise $\Rel\cF(\log n)$, and from \S\ref{sec:duality} on all measures are even),
\begin{equation}\label{eq:Aform}
\mathfrak a_T(F)=C_0\|F\|^2+\Lev(\cF)-\sum_{n\in\NT}v_n\cF(\log n),
\end{equation}
both sides being finite exactly when $\int\log(2+|\tau|)|\Fh|^2<\infty$. Consequently, with
$\Kw=\{\cF:F\in L^2(\IT),\|F\|=1\}$ and the primitive class
\begin{equation}\label{eq:Kprim}
\Kp=\Big\{\Phi\in\Kw:\ \int\Phi(u)\cosh\tfrac u2\,du=0,\ \int\Phi(u)\,u\sinh\tfrac u2\,du=0\Big\},
\end{equation}
one has
\begin{align*}
\lamP&:=\inf_{\Phi\in\Kp}\Big[C_0+\Lev(\Phi)-\sum_{n\in\NT}v_n\Phi(\log n)\Big]=C_0-\Lambda_T,\\
\Lambda_T&:=\sup_{\Phi\in\Kp}\Big[\sum_{n\in\NT}v_n\Phi(\log n)-\Lev(\Phi)\Big],
\end{align*}
and $\lamP\le\lamA$.
\end{theorem}
\begin{proof}
By \eqref{eq:m} and Tonelli (both $1-\cos\tau u$ and $K$ are nonnegative),
$\frac1{2\pi}\int m|\Fh|^2=C_0\|F\|^2+\int_0^\infty K(u)\Big[\frac1{2\pi}\int(1-\cos\tau u)|\Fh|^2d\tau\Big]du$,
and $\frac1{2\pi}\int\cos(\tau u)|\Fh(\tau)|^2d\tau=\cF(u)$ by Plancherel; likewise for the
prime terms. Finiteness: $\int_0^\infty(1-\cos\tau u)K(u)du=\log|\tau|+O(1)$ as $|\tau|\to\infty$
by Lemma~\ref{lem:K}, since $\int_0^1(1-\cos\tau u)\frac{du}u=\log|\tau|+O(1)$. For the last
statements: the two integrals in \eqref{eq:Kprim} are $\Rel(A\bar B)$ and $\Rel(A'\bar B-A\bar B')$
with $A'=\int tF(t)e^{t/2}dt$, $B'=\int tF(t)e^{-t/2}dt$ (expand $\cosh\frac{t-s}2$ and
$(t-s)\sinh\frac{t-s}2$ in exponentials), so $\{\cF:F\in\PT,\|F\|=1\}\subseteq\Kp$, whence
$\lamP\le\lamA$. In terms of $\widehat\Phi$, an entire function of exponential type $2T$ with
$\widehat\Phi(i/2)=\overline{\widehat\Phi(-i/2)}$, the two conditions say that
$\Rel\widehat\Phi(i/2)=0$ and $\Rel\big[i(\widehat\Phi'(i/2)-\widehat\Phi'(-i/2))\big]=0$; when
$F\in\PT$ one has the stronger $\widehat\Phi(\pm i/2)=\widehat\Phi'(\pm i/2)=0$, i.e.\ $\widehat\Phi$
vanishes to order $\ge2$ at $\pm i/2$ (Remark~\ref{rem:Kprim}).
\end{proof}

\begin{remark}\label{rem:Kprim}
$\Kp$ is slightly larger than the set of primitive autocorrelations: for real $F$ its two
conditions read $AB=0$ and $A'B-AB'=0$, i.e.\ $F\in\PT$, or $A=A'=0$, or $B=B'=0$. Every
certificate below bounds $\lamP$, hence $\lamA$; the equality $\lamP=\lamA$ is not needed
anywhere and is not asserted. Since $W_T(F)=\mathfrak a_T(F)$ whenever $A\bar B$ is purely
imaginary, also $\lamW\le\lamP$. Whether $\lamP=\lamW$, i.e.\ whether the optimal multipliers of the
primitive problem are the polar ones, is characterized in Theorem~\ref{thm:mult}.
\end{remark}

\begin{corollary}[monotonicity; upper bound]\label{cor:mono}
With the contact set held fixed, $T\mapsto\Lambda_T$ is nondecreasing and $T\mapsto\lamP$
nonincreasing. Moreover
$\lamA\le C_0+\Aconst+\varphi\big((3\pi/2T)^2\big)$, which equals $1.291$ at $T=\frac12\log3$ and
$2.327$ at $T=\log 2$.
\end{corollary}
\begin{proof}
Monotonicity: $\Kp$ grows with $T$ and $\Lev$ does not depend on $T$. Upper bound: the form
$\Lev(\cF)$ is the Dirichlet form of the L\'evy process of Theorem~\ref{thm:S} killed on
leaving $\IT$; its eigenvalues are bounded by those of the subordinate killed Brownian motion,
$\lambda_n\le\varphi((n\pi/2T)^2)$ \cite[Prop.~2.3]{SongVondracek2003}, the Brownian motion being
normalized with characteristic exponent $\xi^2$; $\PT$ has codimension
$2$, so the bottom of $\Lev$ on $\PT$ is at most $\lambda_3$; and $-\sum v_n\cF(\log n)\le\Aconst\|F\|^2$.
\end{proof}

The following elementary bounds on $m$ and on the prime comb are used throughout.

\begin{lemma}[elementary bounds]\label{lem:mbounds}
\emph{(i)} $m$ is increasing on $[0,\infty)$. \emph{(ii)} For every $\tau>0$,
\[
\log\frac{\tau}{2\pi}-\frac1\tau-\frac1{2\tau^2}\ \le\ m(\tau)\ \le\ \log\frac{\tau}{2\pi}+\frac1\tau+\frac12\log\Big(1+\frac1{4\tau^2}\Big),
\]
and the lower bound, which we denote by $m_{\mathrm{lo}}(\tau)$, is increasing in $\tau$. \emph{(iii)}
\[
-\sum_{n\in\NT}v_n\cos(\tau\log n)\ge-\Aconst,\qquad\text{with}\quad
\sup_\tau\sum_{n\in\NT}v_n\cos(\tau\log n)=\Aconst.
\]
\end{lemma}
\begin{proof}
(i) In $\Rel\psi(\frac14+iy)=-\gamma+\sum_{n\ge0}\big[\frac1{n+1}-\frac{n+\frac14}{(n+\frac14)^2+y^2}\big]$
every summand increases with $|y|$. (ii) Use $\psi(z)=\log z-\frac1{2z}-\int_0^\infty g(t)e^{-zt}dt$
for $\Rel z>0$, $g(t)=\frac1{e^t-1}-\frac1t+\frac12$ \cite[5.9.13]{DLMF}, at $z=\frac14+iy$, $y=\tau/2$.
Here $g(0^+)=0$, $g'(t)=t^{-2}-\frac14\sinh^{-2}(t/2)\in(0,\infty)$ and $g(\infty)=\frac12$, so one
integration by parts gives $|\int_0^\infty g e^{-zt}dt|=|z|^{-1}|\int_0^\infty g'e^{-zt}dt|\le\frac1{2|z|}\le\frac1{2y}$.
Also $\Rel\log z=\frac12\log(\frac1{16}+y^2)\in[\log y,\log y+\frac12\log(1+\frac1{16y^2})]$ and
$\Rel\frac1{2z}=\frac{1/8}{1/16+y^2}\in[0,\frac1{8y^2}]$. Substituting $y=\tau/2$ and $-\log\pi$
gives (ii); the derivative of the lower bound is $\tau^{-1}+\tau^{-2}+\tau^{-3}>0$. (iii) The
first is trivial. For the second, the numbers $\log p$, $p$ prime, are linearly independent over
$\Q$, so by Kronecker's theorem the vector $(\tau\log p)_{p<e^{2T}}$ is dense modulo $2\pi$; at
$\tau$ where all $\cos(\tau\log p)>1-\varepsilon$ also $\cos(\tau k\log p)>1-k^2\varepsilon$, so the
sum over all $n\in\NT$ exceeds $\Aconst-O(\varepsilon)$.
\end{proof}

The hypothesis that the contact set is fixed can be removed from Corollary~\ref{cor:mono}, for
all three bottoms at once.

\begin{lemma}[monotonicity across contact thresholds]\label{lem:monoT}
The functions $T\mapsto\lamA$, $T\mapsto\lamP$ and $T\mapsto\lamW$ are nonincreasing on $(0,\infty)$,
including at the points $T=\frac12\log n$ where a contact enters.
\end{lemma}
\begin{proof}
Let $T<T'$ and extend $F\in L^2(\IT)$ by zero to $I_{T'}$. The norm, the Tate moments $A(F),B(F)$ and
the autocorrelation $\cF$ are unchanged, and $\cF$ vanishes on $|u|\ge2T$. A contact
$n\in\mathcal N_{T'}\setminus\NT$ has $\log n\ge2T$, so $\cF(\log n)=0$, and by \eqref{eq:Aform}
$\mathfrak a_{T'}(F)=\mathfrak a_T(F)$. Hence $\PT\subset\mathcal P_{T'}$ isometrically with the same
form values, which gives the statement for $\lamA$; the same computation with $W_T=\mathfrak a_T+2\Rel(A\bar B)$
gives it for $\lamW$, and for $\lamP$ the class $\Kp$ grows with $T$, $\Lev$ is $T$-free and the new
contacts contribute $\Phi(\log n)=0$.
\end{proof}

For the complete form, continuity is \cite[Thm.~5]{Bombieri2000}:
Bombieri proves that the even and odd minima $\mu_\pm$ of the complete functional
on $L^2([M^{-1},M])$ are continuous and decreasing. Here $W_T$ splits into its values on the even and
odd parts of $F$: the symbol is even, and $\Rel(A\bar B)=|C|^2-|S|^2$ with $C=\int F\cosh\frac t2$, which
sees only the even part, and $S=\int F\sinh\frac t2$, which sees only the odd part. So, after the change of
variables $x=e^t$, $\lamW=\min(\mu_+,\mu_-)$ at $M=e^T$. A direct proof in the present normalization follows.

\begin{lemma}[continuity of the complete bottom]\label{lem:contT}
The function $T\mapsto\lamW$ is continuous on $(0,\infty)$.
\end{lemma}
\begin{proof}
Fix $T$ and let $r\in[\frac12,2]$. The dilation $U_rF(t)=r^{-1/2}F(t/r)$ is unitary from $L^2(\IT)$ onto
$L^2(I_{rT})$, with $\widehat{U_rF}(\tau)=r^{1/2}\Fh(r\tau)$ and $c_{U_rF}(u)=\cF(u/r)$. It maps the form
domain onto the form domain, since $\log(2+|s|/r)\le\log(2+|s|)+\log2$. Hence
$\lambda^{\mathrm W}_{rT}=\inf\{W_{rT}(U_rF):F\text{ in the form domain},\ \|F\|=1\}$. By \eqref{eq:form}
and \eqref{eq:W},
\[
W_{rT}(U_rF)=\frac1{2\pi}\int_\R m(s/r)|\Fh(s)|^2ds-\sum_{2\le n<e^{4T}}v_n\Rel\cF(\tfrac{\log n}r)
+2\Rel\big(A(U_rF)\overline{B(U_rF)}\big),
\]
where the comb may be summed over $n<e^{4T}$ for both windows, because $\cF$ vanishes on $|u|\ge2T$. This
also takes care of contacts that enter or leave. We compare the three terms with those of $W_T(F)$.
\emph{(i) Archimedean.} $m$ is even and $m'(x)=-\frac12\Iml\psi'(\frac14+\frac{ix}2)$ with
$\psi'(z)=O(1/|z|)$ on $\Rel z=\frac14$. So $M_1:=\sup_x|xm'(x)|<\infty$ and
$|m(s/r)-m(s)|\le M_1|\log r|$, and the difference of the archimedean terms is at most
$M_1|\log r|\,\|F\|^2$.
\emph{(ii) Comb.} $|\Rel\cF(u)-\Rel\cF(u')|\le\frac1{2\pi}\int\min(2,|\tau||u-u'|)|\Fh|^2d\tau$. Put
$\delta=4T|1/r-1|$ and $\|F\|_Q^2:=\frac1{2\pi}\int\log(2+|\tau|)|\Fh|^2d\tau$. Splitting at
$|\tau|=\delta^{-1/2}$, the difference of the comb terms is at most
$S_T\big(\delta^{1/2}\|F\|^2+2\|F\|_Q^2/\log(2+\delta^{-1/2})\big)$, where $S_T=\sum_{n<e^{4T}}v_n$.
\emph{(iii) Polar.} $A(U_rF)=r^{1/2}\int F(s)e^{rs/2}ds$, so
$|A(U_rF)-A(F)|\le\|F\|\,\|r^{1/2}e^{rs/2}-e^{s/2}\|_{L^2(\IT)}$, and likewise for $B$. Both moments are
bounded by $C_T\|F\|$ uniformly in $r\in[\frac12,2]$. So the difference of the polar terms is at most
$\delta_P(r)\|F\|^2$ with $\delta_P(r)\to0$ as $r\to1$.
By Lemma~\ref{lem:mbounds}(ii) there is $C_1$ with $\log(2+|\tau|)\le m(\tau)+C_1$. By
Lemma~\ref{lem:mbounds}(iii) and $|2A\bar B|\le4\sinh T\,\|F\|^2$ (Cauchy--Schwarz), this gives
$\|F\|_Q^2\le W_T(F)+C'\|F\|^2$ with $C'=C_1+\Aconst+4\sinh T$. Collecting (i)--(iii), there is
$\eta(r)\to0$ as $r\to1$ with
\[
|W_{rT}(U_rF)-W_T(F)|\le\eta(r)\big(W_T(F)+C'\|F\|^2\big)
\]
for all $F$ in the form domain, and the right side is nonnegative. Take infima over unit $F$. This gives
$(1-\eta)\lamW-\eta C'\le\lambda^{\mathrm W}_{rT}\le(1+\eta)\lamW+\eta C'$, so
$|\lambda^{\mathrm W}_{rT}-\lamW|\le\eta(r)(|\lamW|+C')\to0$.
\end{proof}

The same proof with the comb deleted gives continuity of the bottom $\laminf$ of $\Ainf$
(\S\ref{sec:archstr}). It does not apply directly to $\lamA$, because $U_r$ does not preserve $\PT$;
the correction of the two Tate moments is the content of Lemma~\ref{lem:contA} below.

\begin{lemma}[compactness]\label{lem:compact}
For every $V\in\R$ the set $\{F\in\mathcal D_T:\|F\|=1,\ \mathfrak a_T(F)\le V\}$ is relatively compact in
$L^2(\IT)$. Consequently the form $\mathfrak a_T$ restricted to $\PT\cap\mathcal D_T$ is the form of a
self-adjoint operator $H_T$ on $\PT$ with compact resolvent, $\lamA=\min\operatorname{spec}H_T$ is attained,
and the bottoms of $\mathfrak a_T$ and $W_T$ on the unit sphere of $L^2(\IT)$, and on any closed subspace, are
attained.
\end{lemma}
\begin{proof}
By Lemma~\ref{lem:mbounds}, $\sT\ge m-\Aconst$ and $m-C_0\ge0$, so for $\|F\|=1$ and
$\mathfrak a_T(F)\le V$ one has $\frac1{2\pi}\int(m-C_0)|\Fh|^2\le V+\Aconst-C_0$, hence
$\frac1{2\pi}\int_{|\tau|>R}|\Fh|^2\le(V+\Aconst-C_0)/(m(R)-C_0)\to0$ as $R\to\infty$, uniformly. For each
$R$ the map $F\mapsto\Fh\mathbf 1_{[-R,R]}$ is Hilbert--Schmidt from $L^2(\IT)$ to $L^2(\R)$ (its kernel
$e^{i\tau t}$ is square integrable on $[-R,R]\times\IT$), hence compact. A diagonal argument gives a
subsequence along which $\Fh_n$ converges in $L^2([-R,R])$ for every $R$, and with the uniform tail bound
$\Fh_n$ is Cauchy in $L^2(\R)$. The form $\mathfrak a_T$ is closed and bounded below by $C_0-\Aconst$, since its
form norm is equivalent to $\int\log(2+|\tau|)|\Fh|^2$ (Lemma~\ref{lem:mbounds}(ii)); $\PT$ is closed, and
the polar term $2\Rel(A\bar B)$ is continuous on $L^2(\IT)$. Lower semicontinuity of closed forms gives
attainment, and relative compactness of the sublevel sets is compactness of the resolvent.
\end{proof}

On $\PT$ the operator $H_T$ commutes with $J$, since $\sT$ is even and $A(JF)=B(F)$; so its eigenspaces
split into even and odd parts.

\begin{lemma}[continuity of the primitive bottom]\label{lem:contA}
The function $T\mapsto\lamA$ is continuous on $(0,\infty)$.
\end{lemma}

\begin{proof}
Fix $T$ and consider windows $T'\in[T/2,2T]$ and factors $r\in[\frac12,2]$ with $rT'\in[T/2,2T]$; all
constants depend only on $T$. Put $M_0=\mathsf A_{2T}-C_0$, so that $\mathfrak a_S+M_0\|\cdot\|^2\ge0$
for every $S\le2T$ (Lemma~\ref{lem:mbounds}).

\emph{Step 1.} Steps (i) and (ii) of the proof of Lemma~\ref{lem:contT} do not involve the polar term.
Taking $\delta=8T|1/r-1|$ and summing the comb over $n<e^{8T}$, which covers every window in
$[T/2,2T]$ and every contact that enters or leaves, they give $\eta(r)\to0$ as $r\to1$ with
$|\mathfrak a_{rT'}(U_rF)-\mathfrak a_{T'}(F)|\le\eta(r)\big(\mathfrak a_{T'}(F)+C\|F\|^2\big)$ for all
$F$ in the form domain of $\mathfrak a_{T'}$, uniformly in $T'$.

\emph{Step 2.} If $F\in\mathcal P_{T'}$, then $A(U_rF)=r^{1/2}\int F(s)\big(e^{rs/2}-e^{s/2}\big)ds$,
so $|A(U_rF)|\le\varepsilon(r)\|F\|$ with $\varepsilon(r)=r^{1/2}\|e^{rs/2}-e^{s/2}\|_{L^2(I_{2T})}\to0$,
and likewise for $B$.

\emph{Step 3.} Fix real $\phi_1,\phi_2\in C_c^\infty(I_{T/4})$ for which the matrix
$\big(A(\phi_j),B(\phi_j)\big)_{j}$ is invertible. By zero extension (proof of Lemma~\ref{lem:monoT}),
$\mathfrak a_S(\phi_i,\phi_j)$ does not depend on $S\ge T/4$. For $u=U_rF$ with $F\in\mathcal P_{T'}$,
$\|F\|=1$, let $G=c_1\phi_1+c_2\phi_2$ have the same Tate moments as $u$. Then
$v:=u-G\in\mathcal P_{rT'}$, $\|G\|\le C_2\varepsilon(r)$ and $\mathfrak a_{rT'}(G)\le C_3\varepsilon(r)^2$.

\emph{Step 4.} By the Cauchy--Schwarz inequality for the nonnegative form
$\mathfrak a_{rT'}+M_0\|\cdot\|^2$,
$|\mathfrak a_{rT'}(u,G)|\le(\mathfrak a_{rT'}(u)+M_0)^{1/2}(\mathfrak a_{rT'}(G)+M_0\|G\|^2)^{1/2}+M_0\|G\|$.
Hence, if $\mathfrak a_{T'}(F)\le V$ for a fixed $V$, then
$\mathfrak a_{rT'}(v)\le\mathfrak a_{T'}(F)+\omega(r)$ with $\omega(r)\to0$ uniformly in such $F$. Since
also $\mathfrak a_{rT'}(v)\ge-M_0\|v\|^2$ and $\big|\|v\|-1\big|\le C_2\varepsilon(r)$, the Rayleigh
quotient gives $\lambda_{\mathrm A}(rT')\le\mathfrak a_{T'}(F)+\omega'(r)$ with
$\omega'(r)\to0$.

\emph{Step 5.} By Lemma~\ref{lem:monoT}, $\lamA$ is bounded on $[T/2,2T]$ between $C_0-\mathsf A_{2T}$
and $\lambda_{\mathrm A}(T/2)$; take $V$ one more than the latter. Applying Step 4 to
near-minimizers at $T$ gives $\limsup_{r\to1}\lambda_{\mathrm A}(rT)\le\lamA$, and applying
it to near-minimizers at $rT$ with the factor $1/r$ gives the reverse inequality for the $\liminf$.
\end{proof}


\section{Duality for the window form}\label{sec:duality}

\subsection{Primal and dual}
Fix $T$ and the weights $v=(v_n)_{n\in\NT}$ (any nonnegative weights may be used; $v_n=2\Lambda(n)/\sqrt n$
is the case of $\zeta$).

\begin{definition}\label{def:PD}
\textbf{(P)} $\ \lamP=\inf\big\{\int_\R\sT\,d\mu:\ \mu$ an even probability measure on $\R$,
$\check\mu(u)=0$ for $|u|\ge2T$, $\int\gh0\,d\mu=\int\gh1\,d\mu=0\big\}$.

\textbf{(D)} $\ D_T=\sup\big\{\inf_{\tau\in\R}\big[\sT(\tau)+\nuh(\tau)+\beta_0\gh0(\tau)+\beta_1\gh1(\tau)\big]:\
\nu$ a finite even signed Borel measure with $\supp\nu\subseteq\{|u|\ge2T\}$, $\beta\in\R^2\big\}$.
A triple $(\nu,\beta_0,\beta_1)$ with $f:=\sT+\nuh+\beta\cdot\widehat g\ge c$ on $\R$ is a
\emph{certificate of value $c$}.
\end{definition}

\begin{proposition}[the primal is the eigenvalue problem; weak duality]\label{prop:weak}
The value of \textup{(P)} is the $\lamP$ of Theorem~\ref{thm:A}. Every certificate of value $c$
satisfies $c\le\lamP\le\lamA$.
\end{proposition}
\begin{proof}
If $F$ is a unit vector, $\mu_F=|\Fh|^2d\tau/2\pi$ is feasible ($\check\mu_F=\cF$) with
$\int\sT d\mu_F=\mathfrak a_T(F)$, and $\int\gh j\,d\mu_F=\int g_j\cF$. Conversely if $\mu$ is
feasible, $\Phi=\check\mu$ is continuous, positive definite, $\Phi(0)=1$, supported in $[-2T,2T]$,
so by the Boas--Kac theorem \cite{BoasKac1945} $\Phi=\cF$ for some $F\in L^2(\IT)$, $\|F\|=1$, and
$\Phi\in\Kp$; thus the feasible set of (P) is exactly $\Kp$ read through $\mu=|\Fh|^2d\tau/2\pi$.
Weak duality: for feasible $\mu$, $\int\nuh\,d\mu=\int\check\mu\,d\nu=0$ because $\check\mu=\cF$
vanishes identically on the closed set $\{|u|\ge2T\}$ ($\cF(\pm2T)=0$ by continuity, so an atom
of $\nu$ at $\pm2T$ is harmless), and $\int\gh j\,d\mu=0$; hence
$\int\sT d\mu=\int f\,d\mu\ge c$.
\end{proof}

\begin{lemma}[bounds on the multipliers]\label{lem:gbounds}
$\gh0(\tau)=2\int_0^{2T}\cosh\frac u2\cos\tau u\,du$ and $\gh1(\tau)=2\int_0^{2T}u\sinh\frac u2\cos\tau u\,du$
are even, real, bounded, tend to $0$ at infinity, and for $\tau>0$
\[
|\gh0(\tau)|\le k_0(\tau):=\frac{2\cosh T}{\tau}+\frac{\sinh T}{\tau^2},\qquad
|\gh1(\tau)|\le k_1(\tau):=\frac{8T\sinh T}{\tau},
\]
both bounds decreasing in $\tau$. Moreover $\gh0(\tau)=\frac{2\cosh T\sin2T\tau}{\tau}+O(\tau^{-2})$,
and for any $g=y_0g_0+y_1g_1$, $\widehat g(\tau)=\frac{2g(2T^-)\sin2T\tau}{\tau}+\frac{2g'(2T^-)\cos2T\tau}{\tau^2}+O(\tau^{-3})$.
\end{lemma}
\begin{proof}
Integrate by parts: $\gh0=\frac{2\cosh T\sin2T\tau}\tau-\frac1\tau\int_0^{2T}\sinh\frac u2\sin\tau u\,du$,
and once more, $\frac1\tau\int_0^{2T}\sinh\frac u2\sin\tau u\,du=-\frac{\sinh T\cos2T\tau}{\tau^2}+\frac{\gh0(\tau)}{4\tau^2}$;
solving, $|\gh0|(1+\frac1{4\tau^2})\le\frac{2\cosh T}\tau+\frac{\sinh T}{\tau^2}$. For $\gh1$, one
integration by parts gives $\gh1=\frac{4T\sinh T\sin2T\tau}{\tau}-\frac2\tau\int_0^{2T}\big(\sinh\frac u2+\frac u2\cosh\frac u2\big)\sin\tau u\,du$,
and the integrand is the derivative of $u\sinh\frac u2$, whose integral over $[0,2T]$ is $2T\sinh T$.
The asymptotic expansion is two integrations by parts for a function real-analytic on $[0,2T]$
with $g'(0)=0$.
\end{proof}

\subsection{Strong duality}

\begin{theorem}[strong duality; finitely atomic dual sequences]\label{thm:D1}
$D_T=\lamP$. Moreover, for every $\varepsilon>0$ there are finitely many atoms
$u_1,\dots,u_N\ge2T$, reals $a_1,\dots,a_N$ and $\beta\in\R^2$ such that
$\sT(\tau)+\sum_j2a_j\cos(u_j\tau)+\beta_0\gh0(\tau)+\beta_1\gh1(\tau)>\lamP-\varepsilon$ for all $\tau\in\R$.
In particular, for every $T$ with $\lamP>0$, $\lamA>0$ is provable by a finite certificate.
\end{theorem}
\begin{proof}
\emph{(i) Finite relaxations.} For a finite set $U=\{u_1<\dots<u_N\}\subset[2T,\infty)$ let
$P(U)=\inf\{\int\sT d\mu:\ \mu\ \text{even probability},\ \check\mu(u_i)=0\ \forall i,\ \int\gh j d\mu=0\}$
and $D(U)$ the supremum in (D) restricted to $\nu$ supported on $\pm U$. Then $P(U)\le\lamP$
(the feasible set of (P) is contained in that of $P(U)$) and $D(U)\le D_T\le\lamP$.

\emph{(ii) $P(U)=D(U)$.} Write $f_0\equiv1$, $f_i(\tau)=\cos(u_i\tau)$ ($1\le i\le N$),
$f_{N+1}=\gh0$, $f_{N+2}=\gh1$, $b=(1,0,\dots,0)\in\R^{N+3}$, and let
$C=\{\sum_kt_kf(\tau_k):\ t_k\ge0\}$ be the convex cone generated by the moment curve
$\tau\mapsto f(\tau)=(f_0,\dots,f_{N+2})(\tau)$. The value function
$w(b')=\inf\{\int\sT d\mu:\ \mu\ge0,\ \int f\,d\mu=b'\}$ is convex, $>-\infty$ everywhere
($\sT\ge-\Aconst-|C_0|$ and $\int f_0d\mu=b'_0$ bounds the mass) and finite on $C$ (finite atomic
$\mu$ have finite value). Its Legendre biconjugate at $b$ is the Lagrangian dual
$w^{**}(b)=\sup\{y_0:\ \sT-\sum_ky_kf_k\ge0\text{ on }\R\}$, which is $D(U)$ after
$a_i=-y_i/2$, $\beta_j=-y_{N+1+j}$. A proper convex function agrees with its biconjugate at every
interior point of its domain \cite[Thm.~7.4, Cor.~12.1.1]{Rockafellar}; so it suffices to show
$b\in\operatorname{int}C$.

Suppose not. Since $C$ is a convex cone containing $0$, there is $y\ne0$ with
$\langle y,z\rangle\ge\langle y,b\rangle$ for all $z\in C$; taking $z=0$ gives $y_0\le0$, and
taking $z=tf(\tau)$, $t\to\infty$, gives
$G(\tau):=y_0+\sum_iy_i\cos(u_i\tau)+y_{N+1}\gh0(\tau)+y_{N+2}\gh1(\tau)\ge0$ for all $\tau$.
The Ces\`aro mean $X^{-1}\int_0^XG\to y_0$ (each $\cos(u_i\tau)$ has mean $0$ because $u_i>0$,
and $\gh j\to0$ by Lemma~\ref{lem:gbounds}); as $G\ge0$ this forces $y_0=0$, and then
$X^{-1}\int_0^X|G|\to0$. Testing against the bounded functions $\cos(u_i\tau)$ gives
$y_i/2=\lim X^{-1}\int_0^XG\cos(u_i\tau)\,d\tau=0$ for each $i$. There remains
$\widehat g\ge0$ with $g=y_{N+1}g_0+y_{N+2}g_1$. By the expansion in Lemma~\ref{lem:gbounds},
$\widehat g\ge0$ forces $g(2T^-)=0$ and then $g'(2T^-)=0$; but
$g(2T^-)=y_{N+1}\cosh T+2T\,y_{N+2}\sinh T$ and $g'(2T^-)=\frac12y_{N+1}\sinh T+y_{N+2}(\sinh T+T\cosh T)$,
whose determinant is $\sinh T\cosh T+T>0$. So $y=0$, a contradiction; hence $b\in\operatorname{int}C$.

\emph{(iii) The limit.} Choose finite sets $U_1\subset U_2\subset\cdots$ with $\bigcup U_n$ dense in
$[2T,\infty)$; $P(U_n)$ is nondecreasing and $\le\lamP$. Pick feasible $\mu_n$ with
$\int\sT d\mu_n\le P(U_n)+1/n$. Since $\sT\to+\infty$ and is bounded below, $(\mu_n)$ is tight;
let $\mu$ be a weak-$*$ limit point, a probability measure. For $u\in U_N$ and $n\ge N$,
$\check\mu_n(u)=0$ and $\cos(u\,\cdot)$ is bounded continuous, so $\check\mu(u)=0$; $\check\mu$
is continuous, so $\check\mu\equiv0$ on $[2T,\infty)$; likewise $\int\gh jd\mu=0$. Thus $\mu$ is
feasible for (P), and by lower semicontinuity of $\mu\mapsto\int\sT d\mu$ under weak convergence
($\sT$ continuous and bounded below), $\lamP\le\int\sT d\mu\le\liminf\int\sT d\mu_n=\lim P(U_n)\le\lamP$.
Hence $D_T\ge\sup_nD(U_n)=\sup_nP(U_n)=\lamP\ge D_T$, and the atomic statement follows.
\end{proof}

\begin{remark}\label{rem:D1complete}
Theorem~\ref{thm:D1} says that, at a fixed window, $\lamP\ge0$ holds if and only if for every
$\varepsilon>0$ there are an atomic $\nu$ on $\{|u|\ge2T\}$ and $\beta\in\R^2$ with
$\sT+\nuh+\beta\cdot\widehat g\ge-\varepsilon$ on $\R$; and $\lamP\ge0$ implies $\lamA\ge0$. The same
proof, with the two $\widehat g$-constraints deleted, gives strong duality for the complete form:
$\lamW=\sup\{\inf_\tau[\sT+2\gh0+\nuh]\}$ over the same class of $\nu$ (Proposition~\ref{prop:P}
explains the constant $2$).
\end{remark}

\subsection{Non-attainment}

\begin{theorem}[non-attainment]\label{thm:O}
Let $\nu$ be any finite even signed measure with $\supp\nu\subseteq\{|u|\ge2T\}$ and let
$h_0,\dots,h_J$ be finitely many continuous functions on $[-2T,2T]$ that annihilate $\Kp$
(for instance $g_0,g_1$). Then for every $\beta\in\R^{J+1}$
\[
\inf_{\tau\in\R}\Big[\sT(\tau)+\nuh(\tau)+\sum_j\beta_j\widehat h_j(\tau)\Big]<\lamP\quad\text{strictly.}
\]
The statement holds for every finite weight family, including $v=0$.
\end{theorem}
\begin{proof}
A minimizer $\Phi_\star=c_{F_\star}\in\Kp$ of the primal exists: the sublevel sets of
$\mathfrak a_T$ on the unit sphere of $L^2(\IT)$ are compact (Lemma~\ref{lem:compact}),
and, in the language of measures, the feasible set of (P) is tight and weak-$*$ closed and $\mu\mapsto\int\sT d\mu$ is
weak-$*$ lower semicontinuous, exactly as in step (iii) of the proof of Theorem~\ref{thm:D1}. Put $\varrho=\sT+\nuh+\sum\beta_j\widehat h_j-\lamP$, a
continuous function ($\nu$ finite). By the annihilation properties,
$\frac1{2\pi}\int\varrho|\Fh_\star|^2=\mathfrak a_T(F_\star)-\lamP=0$. If $\varrho\ge0$ then $\varrho=0$ a.e.\ on
$\{\Fh_\star\ne0\}$, whose complement is the discrete zero set of a nonzero entire function of
exponential type; so $\varrho\equiv0$ on $\R$, and by Fourier inversion, as distributions in $u$,
\[
\check\sT+\nu+\sum_j\beta_jh_j=\lamP\,\delta_0 .
\]
Restrict to the open set $0<|u|<2T$, where $\nu=0$: by \eqref{eq:m}, $\check m=-\frac12K(|u|)$
there, and the prime part of $\check\sT$ is $-\sum_{n\in\NT}\frac{v_n}2(\delta_{\log n}+\delta_{-\log n})$;
hence $\sum_j\beta_jh_j(u)=\frac12K(|u|)+\sum_{n\in\NT}\frac{v_n}2(\delta_{\log n}+\delta_{-\log n})$
on $0<|u|<2T$. The left side is a bounded continuous function; the right side is unbounded as
$u\to0$ ($K(u)\sim1/u$) and carries atoms as soon as $\NT\neq\emptyset$. Contradiction.
\end{proof}

\begin{remark}\label{rem:OvsB}
Theorem~\ref{thm:O} forbids attainment, not certificates: whenever $\lamP>0$ any dual value in
$(0,\lamP)$ suffices, and Theorem~\ref{thm:D1} says such values are reached by finite atomic
$\nu$. Every maximizing sequence of the dual must therefore degenerate (its total variation is
unbounded; Corollary~\ref{cor:T2} below quantifies the cost). On the primal side the infimum is
attained \cite[Thm.~3]{Bombieri2000}. The proof shows the reason for the asymmetry: the
singularity $K(u)\sim1/u$ of the archimedean L\'evy density at $u=0$ and the prime atoms lie inside
the window, where the dual measure is not allowed to have support.
\end{remark}

\subsection{The Euler--Lagrange equation}

\begin{theorem}[finite-interval Wiener--Hopf equation]\label{thm:W}
Let $\mathsf S_T$ denote the convolution operator on $\R$ with symbol $\sT$,
$\widehat{\mathsf S_TF}=\sT\Fh$, acting on $L^2(\IT)$ extended by zero and taking values in tempered
distributions. A unit vector $F\in\PT\cap\mathcal D_T$ is a critical point of $\mathfrak a_T$ on the unit
sphere of $\PT$ with critical value $\lambda$ if and only if there are $\alpha,\beta\in\C$ (real for real
$F$) with
\[
(\mathsf S_T-\lambda)F=\alpha\cosh\tfrac t2+\beta\sinh\tfrac t2\qquad\text{on }\IT,
\]
in the sense of distributions on $\IT$. The symbol $\sT-\lambda$ is real, even, tends to
$+\infty$, and changes sign for every $\lambda>C_0-\Aconst$; in that range it is not elliptic, and no
Wiener--Hopf index is defined for it.
\end{theorem}
\begin{proof}
For $G\in C_c^\infty(\IT)$, $\langle(\mathsf S_T-\lambda)F,G\rangle=\frac1{2\pi}\int(\sT-\lambda)\Fh\,\overline{\widehat G}$
is the polarization of $\mathfrak a_T-\lambda\|\cdot\|^2$. Criticality on the unit sphere of $\PT$
with value $\lambda$ means that this vanishes for all $G$ in $\PT\cap\mathcal D_T$; since
$C_c^\infty(\IT)\cap\PT$ is dense in $\PT\cap\mathcal D_T$ for the form norm (mollify and correct the two
Tate moments by a fixed pair of $C_c^\infty$ functions), it is equivalent to require it for
$G\in C_c^\infty(\IT)$ with $\int Ge^{t/2}=\int Ge^{-t/2}=0$. This is a subspace of codimension $2$
of $C_c^\infty(\IT)$, and a distribution on $\IT$ that annihilates the common kernel of two continuous
linear functionals is a linear combination of them; the functionals are integration against
$e^{\pm t/2}$, i.e.\ against $\cosh\frac t2\pm\sinh\frac t2$. Conversely, if $(\mathsf S_T-\lambda)F$
is such a combination on $\IT$, it annihilates every $G$ in that subspace, so $F$ is critical. Finally
$\sT(0)=C_0-\Aconst$ and $\sT\to+\infty$ (Lemma~\ref{lem:mbounds}).
\end{proof}

\subsection{The polar multiplier}

\begin{lemma}[moment interior]\label{lem:W}
For $F\in L^2(\IT)$ real put $Q_0(F)=\int\cF\cosh\frac u2=A B$ and
$Q_1(F)=\int\cF\,u\sinh\frac u2=A'B-AB'$. The image of the set of unit vectors under
$F\mapsto(Q_0,Q_1)$ contains a neighborhood of $0$ in $\R^2$, realised by $C_c^\infty(\IT)$
functions; equivalently, $0$ is an interior point of the (convex) moment image
$\mathfrak m(\mu)=(\int\gh0d\mu,\int\gh1d\mu)$ of the feasible set of (P) without the two
constraints, and the interior is realised by measures of finite value.
\end{lemma}
\begin{proof}
The formulas for $Q_0,Q_1$ follow from expanding $\cosh\frac{t-s}2$ and $(t-s)\sinh\frac{t-s}2$
in exponentials. The functions $e^{t/2},e^{-t/2},te^{t/2},te^{-t/2}$ are linearly independent
on $\IT$, and a nontrivial combination is real-analytic, so it cannot annihilate
$C_c^\infty(\IT)$: the linear map $F\mapsto(A,B,A',B')$ is onto $\R^4$ from $C_c^\infty(\IT)$.
If $0$ were not interior, some $(t_0,t_1)\ne0$ would satisfy $t_0Q_0(F)+t_1Q_1(F)\ge0$ for all
$F$ (normalization only rescales $Q_j$ by a positive factor). Take $(A,B,A',B')=(1,0,0,\mp1)$:
$Q_0=0$, $Q_1=\pm1$, so $t_1=0$; take $(1,\pm1,0,0)$: $Q_1=0$, $Q_0=\pm1$, so $t_0=0$. A finite
convex combination of the four witnesses covers a ball, and each has finite value because
$\Fh$ is Schwartz. For the image of the unit sphere itself: a function with
$(A,B,A',B')=(1,x,0,-y)$ has $(Q_0,Q_1)=(x,y)$; adding an element of the kernel of the moment map
does not change $(Q_0,Q_1)$ and makes the norm any number above the minimal one, and
$(Q_0,Q_1)(F/\|F\|)=(Q_0,Q_1)(F)/\|F\|^2$. Since the minimal norm is bounded on $x^2+y^2=1$, every
point $t(x,y)$ with $0<t\le t_0$, $t_0>0$ uniform, is attained on the unit sphere, and $(0,0)$ is
attained on $\PT$.
\end{proof}

\begin{proposition}[the polar multiplier]\label{prop:P}
For every $F\in L^2(\IT)$ in the form domain,
$\frac1{2\pi}\int\gh0|\Fh|^2=\Rel\big(A(F)\overline{B(F)}\big)$ and hence
\[
\frac1{2\pi}\int\big(\sT+2\gh0\big)|\Fh|^2\,d\tau=W_T(F);
\]
if moreover $F$ is real and $\sum_\rho|\Phi_F(\rho)\Phi_F(1-\rho)|<\infty$ (in particular if
$F\in C_c^\infty(\IT)$ is real), then $W_T(F)=\sum_\rho\Phi_F(\rho)\Phi_F(1-\rho)$
(Lemma~\ref{lem:roughEF} below).
Consequently a certificate with $(\beta_0,\beta_1)=(2,0)$ and value $c$ proves $W_T\ge c$ on all
of $L^2(\IT)$, and $\beta_0=2$ is the unique value of $\beta_0$ with this property.
\end{proposition}
\begin{proof}
$\frac1{2\pi}\int\gh0|\Fh|^2=\int g_0\cF=\iint\cosh\frac{t-s}2F(t)\overline{F(s)}$ and
$\cosh\frac{t-s}2=\frac12e^{t/2}e^{-s/2}+\frac12e^{-t/2}e^{s/2}$ gives $\frac12(A\bar B+B\bar A)$.
The identification of $W_T$ with the zero sum is the explicit formula for the autocorrelation
$\cF$, \cite{Weil52,Bombieri2000}, \cite[Thm.~8.4]{Paper45}, whose polar term is
$2\Phi_F(0)\Phi_F(1)=2\Rel(A\bar B)$ for the real forms considered here. Uniqueness: for
$\beta_0\ne2$, $\frac1{2\pi}\int(\sT+\beta_0\gh0)|\Fh|^2=W_T(F)+(\beta_0-2)\Rel(A\bar B)$, and
$\Rel(A\bar B)$ takes both signs on the unit sphere of $L^2(\IT)$ (Lemma~\ref{lem:W}, witnesses
$(1,\pm1,0,0)$), so no lower bound of the left side bounds $W_T$; the same argument with
$\Rel(A'\bar B-A\bar B')$ forces $\beta_1=0$.
\end{proof}

\cite[Thm.~8.4]{Paper45} is stated for real $F\in C_c^\infty$. The first equality of
Proposition~\ref{prop:P} holds on the whole form domain and is the only one the certificates use;
wherever the zero sum is used for a rough test (Theorem~\ref{thm:margin}), it enters through the
following lemma.

\begin{lemma}[the explicit formula for rough tests]\label{lem:roughEF}
Let $F\in L^2(\IT)$ be real, with \emph{(a)} $\int\log(2+|\tau|)|\Fh|^2<\infty$ and \emph{(b)}
$\sum_\rho|\Phi_F(\rho)\Phi_F(1-\rho)|<\infty$. Then
\[
\sum_\rho\Phi_F(\rho)\Phi_F(1-\rho)=W_T(F)=2\Phi_F(0)\Phi_F(1)+\mathfrak a_T(F),
\]
the sum being absolutely convergent. In particular $\mathfrak a_T(F)=\sum_\rho\Phi_F(\rho)\Phi_F(1-\rho)$
for $F\in\PT$, and a zero $\rho=\frac12+i\gamma$ on the critical line contributes $|\Fh(\gamma)|^2$.
\end{lemma}
\begin{proof}
For real $F$, $\Phi_F(1)=A(F)$ and $\Phi_F(0)=B(F)$, so $2\Phi_F(0)\Phi_F(1)=2\Rel(A\bar B)$. Let
$\eta_\varepsilon\ge0$, $\eta_\varepsilon\in C_c^\infty(-\varepsilon,\varepsilon)$, $\int\eta_\varepsilon=1$,
and $F_\varepsilon=F*\eta_\varepsilon\in C_c^\infty$, supported in $[-T-\varepsilon,T+\varepsilon]$.
Apply \cite[Thm.~8.4]{Paper45} to $F_\varepsilon$ at the window $T+\varepsilon$. Zero side:
$\Phi_{F_\varepsilon}=\Phi_F\Phi_{\eta_\varepsilon}$, $|\Phi_{\eta_\varepsilon}(s)|\le e^{\varepsilon|\Rel s-\frac12|}\le e^{\varepsilon/2}$
in the critical strip and $\Phi_{\eta_\varepsilon}\to1$ pointwise, so by (b) and dominated convergence
the zero sum (which for $F_\varepsilon$ converges absolutely, so that the symmetric order is
immaterial) tends to $\sum_\rho\Phi_F(\rho)\Phi_F(1-\rho)$. Polar term:
$\Phi_{F_\varepsilon}(0)\Phi_{F_\varepsilon}(1)\to\Phi_F(0)\Phi_F(1)$. Archimedean term:
$|\Fh_\varepsilon|^2=|\Fh|^2|\widehat{\eta_\varepsilon}|^2\le|\Fh|^2$ and $|m(\tau)|\le|C_0|+\log(1+|\tau|)$
(Lemma~\ref{lem:mbounds}(i),(ii); $m\le0$ for $|\tau|\le\tau_0$), so by (a) and dominated convergence
$\frac1{2\pi}\int m|\Fh_\varepsilon|^2\to\frac1{2\pi}\int m|\Fh|^2$. Prime term: $c_{F_\varepsilon}\to\cF$
uniformly, the sum runs over the finitely many $n<e^{2T+2\varepsilon}$, and for
$e^{2T}\le n<e^{2T+2\varepsilon}$ one has $\cF(\log n)=0$; so it tends to $-\sum_{n\in\NT}v_n\cF(\log n)$.
\end{proof}


\section{Atomic certificates}\label{sec:tails}

A certificate is an inequality $f\ge c$ on all of $\R$. The lemmas of this section reduce it to a
bounded interval; they rest on Lemmas~\ref{lem:mbounds} and \ref{lem:gbounds}.

\subsection{Concentration of the near-optimal feasible measures}\label{subsec:conc}

The program (P) has a minimizer, but no dual optimum is attained (Theorem~\ref{thm:O}); a reduction of
the dual value to a bounded interval cannot be routed through an extremal certificate, and has to be
uniform over the near-optimal feasible set. Write
\[
\mathcal F_\varepsilon=\Big\{\mu\ \text{feasible for (P)}:\ \int\sT\,d\mu\le\lamP+\varepsilon\Big\},
\qquad B_*=\Aconst-C_0,\qquad G(\varepsilon)=\lamP+\varepsilon+B_*,
\]
which is non-empty for every $\varepsilon>0$.

\begin{theorem}[concentration of near-optimal feasible measures]\label{thm:P3}
Let $\varepsilon>0$ and $\eta\in(0,1)$.
\emph{(i)} Every $\mu\in\mathcal F_\varepsilon$ satisfies $\int(m-C_0)\,d\mu\le G(\varepsilon)$ and,
for every $R>0$, $\mu(\{|\tau|\ge R\})\le G(\varepsilon)/(m(R)-C_0)$.
\emph{(ii)} Consequently $\mu(\{|\tau|\ge R\})\le\eta$ for every $\mu\in\mathcal F_\varepsilon$ and
every $R\ge\max\{2,R(\varepsilon,\eta)\}$, where
\[
R(\varepsilon,\eta):=2\pi\exp\Big(1+C_0+\frac{G(\varepsilon)}\eta\Big).
\]
\emph{(iii)} Let $(\nu,\beta)$ be a certificate of value $c$, put
$E_R:=\Aconst+\|\nu\|_{\TV}+|\beta_0|k_0(R)+|\beta_1|k_1(R)-c$, and let
$R\ge\max\{2,R(\varepsilon,\eta)\}$. Then, for every $\mu\in\mathcal F_\varepsilon$,
\[
0\le\lamP-c\le\int_{[-R,R]}(f-c)\,d\mu+\varepsilon_{\mathrm{tail}}(\mu,R),\qquad
\varepsilon_{\mathrm{tail}}(\mu,R)\le\int_{|\tau|>R}(m-C_0)\,d\mu+(E_R+C_0)^+\,\eta ,
\]
the first tail integral being $\le G(\varepsilon)$ always, and $\to0$ as $R\to\infty$ for each fixed
$\mu$.
\end{theorem}
\begin{proof}
(i) $m-C_0\ge0$ by Lemma~\ref{lem:mbounds}(i), and $\int\sT d\mu\ge\int m\,d\mu-\Aconst$ by
Lemma~\ref{lem:mbounds}(iii); hence $\int(m-C_0)d\mu\le\int\sT d\mu+\Aconst-C_0\le G(\varepsilon)$.
Chebyshev's inequality on $\{|\tau|\ge R\}$, where $m-C_0\ge m(R)-C_0$, gives the second bound. (ii) By
Lemma~\ref{lem:mbounds}(ii), $m(R)\ge\log(R/2\pi)-1$ for $R\ge2$, so $m(R)-C_0\ge G(\varepsilon)/\eta$
as soon as $\log(R/2\pi)\ge1+C_0+G(\varepsilon)/\eta$. (iii) By Proposition~\ref{prop:weak},
$\int f\,d\mu=\int\sT d\mu$ for feasible $\mu$, so $\int(f-c)d\mu=\int\sT d\mu-c\in[\lamP-c,\lamP+\varepsilon-c]$;
both pieces of the split at $R$ are nonnegative because $f\ge c$. On $|\tau|>R$,
$f-c\le m+E_R=(m-C_0)+(E_R+C_0)\le(m-C_0)+(E_R+C_0)^+$, by Lemma~\ref{lem:mbounds}(iii),
$|\nuh|\le\|\nu\|_{\TV}$ and Lemma~\ref{lem:gbounds} ($k_0,k_1$ decreasing); apply (i) to the first
summand and (ii) to the second. For fixed $\mu$, $\int(m-C_0)d\mu<\infty$, and the tail integral tends to
$0$ by dominated convergence.
\end{proof}

The positive part in (iii) is needed: $C_0<0$, so $E_R+C_0$ can be negative for certificates of
small total variation, although it is positive for the certificate of Theorem~\ref{thm:C}
($\|\nu\|_{\TV}=90.75$).

\begin{corollary}[verification on a bounded interval]\label{cor:AP}
Suppose $(\nu,\beta)$ satisfies $f\ge c$ \emph{only} on $[-R,R]$ with
$R\ge\max\{2,R(\varepsilon,\eta)\}$, and put
\[
\mathsf V:=\Aconst+\|\nu\|_{\TV}+|\beta_0|\,\gh0(0)+|\beta_1|\,\gh1(0)-C_0 ,
\]
so that $f\ge-\mathsf V$ on all of $\R$. Then $c\le(\lamP+\varepsilon+\mathsf V\eta)/(1-\eta)$ for
every $\varepsilon>0$.
\end{corollary}
\begin{proof}
Since $g_0,g_1\ge0$ on $[0,2T]$, $|\gh j|\le\gh j(0)$, and $m\ge C_0$; this gives $f\ge-\mathsf V$.
For $c<0$ there is nothing to prove. For $c\ge0$ take $\mu\in\mathcal F_\varepsilon$: then
$\lamP+\varepsilon\ge\int\sT d\mu=\int f\,d\mu\ge c\,\mu([-R,R])-\mathsf V\mu(|\tau|>R)\ge c(1-\eta)-\mathsf V\eta$
by Theorem~\ref{thm:P3}(ii).
\end{proof}

\begin{remark}\label{rem:P3exp}
The exponential dependence of $R(\varepsilon,\eta)$ on $1/\eta$ is forced and is not an artifact of
Chebyshev's inequality. Theorem~\ref{thm:Vp}(ii) below constructs, for every $\delta>0$, feasible measures
$\mu'_{\eta,R}$ with mass $\ge\eta(1-o(1))$ beyond $R-\sqrt R$ and value
$\le\lamP+\delta+\eta(\log R+C_1)+o(\eta)$; such a $\mu'$ lies in $\mathcal F_\varepsilon$ as soon as
$\eta(\log R+C_1)+o(\eta)\le\varepsilon-\delta$. So the radius at which mass $\eta$ is forced to have
left satisfies $\log R\asymp1/\eta$, exactly the shape of (ii); we do not know the optimal constant,
which lies between $\varepsilon$ and $G(\varepsilon)$. The radius $R(\varepsilon,\eta)$ has the same
shape as the envelope threshold of Proposition~\ref{prop:env}, with $B_*/\eta$ in place of $\Aconst$.
\end{remark}

\subsection{Lattice atoms: the graded periodic tail}

\begin{lemma}[graded periodic tail]\label{lem:F1}
Let $d=2T/M$, $u_j=jd$ for $M\le j\le J$ (so $u_M=2T$), $\nu=\sum_ja_j(\delta_{u_j}+\delta_{-u_j})$,
$P=2\pi/d$, so that $\nuh(\tau)=\sum_j2a_j\cos(u_j\tau)$ is $P$-periodic. Let
$f=\sT+\nuh+\beta_0\gh0+\beta_1\gh1$. Fix $\tau_1>0$, $N\ge1$, and for $x\in[\tau_1,\tau_1+P)$ put
\begin{align*}
Y_k(x)&=m(x+kP)-\sum_{n\in\NT}v_n\cos\big((x+kP)\log n\big)-|\beta_0|k_0(x+kP)-|\beta_1|k_1(x+kP)\ \ (k<N),\\
Y_N(x)&=m(x+NP)-\Aconst-|\beta_0|k_0(x+NP)-|\beta_1|k_1(x+NP).
\end{align*}
If $\nuh(x)+Y_k(x)\ge c$ for all $x\in[\tau_1,\tau_1+P)$ and all $0\le k\le N$, then $f(\tau)\ge c$
for all $\tau\ge\tau_1$.
\end{lemma}
\begin{proof}
Every $\tau\ge\tau_1$ is uniquely $x+kP$ with $x\in[\tau_1,\tau_1+P)$, $k\ge0$, and
$\nuh(\tau)=\nuh(x)$ since $u_jP=2\pi j$. For $k<N$ the $k$-th row is $f(x+kP)$ with the exact
prime term and the two multiplier terms replaced by the lower bounds of Lemma~\ref{lem:gbounds}.
For $k\ge N$: $m(x+kP)\ge m(x+NP)$, $k_i(x+kP)\le k_i(x+NP)$, and the prime term is $\ge-\Aconst$.
\end{proof}

The prime term is never required to be periodic: it is evaluated exactly on the first $N$
copies of one period of $\nuh$ and replaced by its infimum beyond. The lemma therefore holds
for any finite set of contacts; the incommensurability of $\log2$ and $\log3$ plays no role.
It is a sufficient criterion only (the rows are conservative in $\beta$ by
$|\beta_0|k_0+|\beta_1|k_1$).

\subsection{Exponential columns: the monotone tail}

\begin{proposition}[exponential columns]\label{prop:T0}
For $s>0$ and an integer $p\ge0$ let $d\nu_{s,p}(u)=\frac12(|u|-2T)^pe^{-s(|u|-2T)}\mathbf 1_{|u|\ge2T}du$,
a positive finite measure of mass $p!/s^{p+1}$ supported outside the window. Then
\[
\nuh_{s,p}(\tau)=\Rel\Big[e^{2iT\tau}\frac{p!}{(s-i\tau)^{p+1}}\Big],\qquad
|\nuh_{s,p}(\tau)|\le\operatorname{cap}_{s,p}(\tau):=\frac{p!}{(s^2+\tau^2)^{(p+1)/2}},
\]
and $\operatorname{cap}_{s,p}$ is decreasing in $\tau\ge0$.
\end{proposition}
\begin{proof}
$\nuh_{s,p}(\tau)=\int_{2T}^\infty(u-2T)^pe^{-s(u-2T)}\cos\tau u\,du=\Rel\big[e^{2iT\tau}\int_0^\infty t^pe^{-(s-i\tau)t}dt\big]$.
\end{proof}

\begin{lemma}[monotone tail]\label{lem:T}
Let $\nu=\sum_{k=1}^Kc_k\nu_{s_k,p_k}+\sum_{j=1}^Ja_j(\delta_{u_j}+\delta_{-u_j})$ with $u_j\ge2T$, and
$f=\sT+\nuh+\beta_0\gh0+\beta_1\gh1$. Define for $\tau>0$
\[
L(\tau)=\log\frac\tau{2\pi}-\frac1\tau-\frac1{2\tau^2}-\Aconst-\sum_k|c_k|\operatorname{cap}_{s_k,p_k}(\tau)
-2\sum_j|a_j|-|\beta_0|k_0(\tau)-|\beta_1|k_1(\tau).
\]
Then $f\ge L$ on $(0,\infty)$, $L$ is nondecreasing, and $L(\tau_1)\ge c$ for one $\tau_1$ implies
$f\ge c$ on $[\tau_1,\infty)$.
\end{lemma}
\begin{proof}
Lemma~\ref{lem:mbounds}(ii),(iii), Proposition~\ref{prop:T0}, the trivial bound on the atoms and
Lemma~\ref{lem:gbounds}; every term of $L$ is nondecreasing.
\end{proof}

\begin{corollary}[cost of the tail]\label{cor:T2}
For an atom-free $\nu$ with $A:=\sum_k|c_k|p_k!$ and $A':=|\beta_0|(2\cosh T+\sinh T)+|\beta_1|8T\sinh T$,
the tail inequality $L(\tau_1)\ge c$ holds for every $\tau_1\ge\max\{A+A'+2,\ 2\pi e^{c+\Aconst+1}\}$,
so the tail costs linearly in the $\ell^1$ size of the certificate. For atoms it costs
exponentially, $\tau_1\approx2\pi e^{\Aconst+2\sum|a_j|}$, which is why Lemma~\ref{lem:F1} and not
Lemma~\ref{lem:T} is used for lattice certificates. Neither cost is uniform over the dual class:
by Theorem~\ref{thm:O}, $\|\nu\|\to\infty$ along every maximizing sequence.
\end{corollary}
\begin{proof}
For $\tau\ge1$, $\operatorname{cap}_{s,p}(\tau)\le p!/\tau$ and $|\beta_i|k_i\le A'/\tau$, so
$L(\tau)\ge\log\frac\tau{2\pi}-\Aconst-\frac{A+A'+2}{\tau}$.
\end{proof}

\begin{lemma}[hybrid tail]\label{lem:hybrid}
Let $\nu=\nu_{\mathrm{lat}}+\sum_\ell c_\ell\nu_{s_\ell,p_\ell}$ with $\nu_{\mathrm{lat}}$ as in Lemma~\ref{lem:F1}.
For $x\in[\tau_1,\tau_1+P)$ define $H_k(x)=\nuh_{\mathrm{lat}}(x)+Y_k(x)+\sum_\ell c_\ell\nuh_{s_\ell,p_\ell}(x+kP)$
for $0\le k<N$, and $H_N(x)=\nuh_{\mathrm{lat}}(x)+Y_N(x)-\sum_\ell|c_\ell|\operatorname{cap}_{s_\ell,p_\ell}(x+NP)$.
If $H_k\ge c$ on $[\tau_1,\tau_1+P)$ for all $0\le k\le N$ then $f\ge c$ on $[\tau_1,\infty)$.
\end{lemma}
\begin{proof}
As in Lemma~\ref{lem:F1}; on the rows $k<N$ the column term is evaluated exactly, and for
$k\ge N$ it is bounded by $-\sum_\ell|c_\ell|\operatorname{cap}(x+kP)\ge-\sum_\ell|c_\ell|\operatorname{cap}(x+NP)$
since $\operatorname{cap}$ decreases.
\end{proof}

\subsection{Exhaustion by exponential columns}

\begin{theorem}[exponential columns exhaust the dual]\label{thm:E}
Let $\mathcal E$ be the real span of $\{\nu_{s,p}:s>0,p\ge0\}$. Then
$\sup\{\inf_\tau[\sT+\nuh+\beta\cdot\widehat g]:\ \nu\in\mathcal E,\beta\in\R^2\}=\lamP$, and the same
holds with $\mathcal E$ replaced by the absolutely continuous measures with bounded density and
compact support in $\{|u|\ge2T\}$.
\end{theorem}
\begin{proof}
``$\le$'' is weak duality. For ``$\ge$'' fix $\varepsilon>0$ and, by Theorem~\ref{thm:D1}, an atomic
$\nu_0=\sum_{j\le J}a_j(\delta_{u_j}+\delta_{-u_j})$ and $\beta$ with
$c_0:=\inf[\sT+\nuh_0+\beta\cdot\widehat g]>\lamP-\varepsilon/2$. By Lemma~\ref{lem:T} there is
$\tau_1<\infty$ with $\sT+\nuh_0+\beta\cdot\widehat g\ge c_0-\varepsilon/8$ on $[\tau_1,\infty)$,
and the lemma uses only $\sum|a_j|$, so the same bound holds for every signed measure of the
same total variation supported in $\{|u|\ge2T\}$. Replace each atom $a_j\delta_{u_j}$ by
$a_j$ times the uniform probability density on $(u_j,u_j+\eta)$: since
$|\cos u\tau-\cos u_j\tau|\le|\tau|\eta$, the transform changes by at most $2\sum|a_j|\tau_1\eta$
on $[0,\tau_1]$, which is $<\varepsilon/8$ for small $\eta$; the new $\nu_1$ has bounded density
and compact support and value $\ge c_0-\varepsilon/4$. Finally $\{t^pe^{-t}\}_{p\ge0}$ is total in
$L^1(0,\infty)$ (if $g\in L^\infty$ annihilates them, its Laplace transform is holomorphic on
$\Rel z>0$ with all derivatives vanishing at $z=1$, hence $g=0$), so there is $\nu_2\in\mathcal E$
with $\|\nu_2-\nu_1\|_{\TV}<\varepsilon/8$, and $\|\nuh_2-\nuh_1\|_\infty\le\|\nu_2-\nu_1\|_{\TV}$ on all of $\R$.
\end{proof}


\begin{proposition}[the canonical column; band-limitation]\label{prop:M}
Let $d\nu_K=\frac12K(|u|)\mathbf 1_{|u|\ge2T}du$, a positive measure of mass
$c_K=\int_{2T}^\infty K$ (Theorem~\ref{thm:S}(iii)). Then for all $\tau$
\[
\sT(\tau)+\nuh_K(\tau)=\big(C_0+c_K\big)+\int_0^{2T}(1-\cos\tau u)K(u)\,du-\sum_{n\in\NT}v_n\cos(\tau\log n),
\]
a function whose distributional Fourier transform is supported in $[-2T,2T]$ (it still grows
like $\log|\tau|$: the singularity of $K$ at $u=0$ lies inside the window). Moreover
$d\nu_K=2e^{-T}\big(1+O(e^{-2|u|})\big)\,d\nu_{\frac12,0}$:
the archimedean L\'evy density selects the scale $s=\frac12$ among the exponential columns.
\end{proposition}
\begin{proof}
By \eqref{eq:m}, split $\int_0^\infty(1-\cos\tau u)K$ at $2T$ and use
$\nuh_K(\tau)=\int_{2T}^\infty K(u)\cos\tau u\,du$; the tail contributes $c_K-\nuh_K(\tau)$. For the support
statement let $\psi$ be a Schwartz function whose Fourier transform vanishes on a neighbourhood of
$[-2T,2T]$. Then $\int\psi=0$, $\int\psi(\tau)\cos(\tau\log n)d\tau=0$ for $n\in\NT$, and, by Fubini
(justified by $|1-\cos\tau u|K(u)\le\min(2,\tau^2u^2)K(u)$),
$\int\psi(\tau)\int_0^{2T}(1-\cos\tau u)K(u)\,du\,d\tau=\int_0^{2T}K(u)\big(\widehat\psi(0)-\Rel\widehat\psi(u)\big)du=0$.
Finally, by Lemma~\ref{lem:K}, $\frac12K(u)=e^{-u/2}(1+O(e^{-2u}))$, while the density of $\nu_{\frac12,0}$
is $\frac12e^{T}e^{-u/2}$ for $u\ge2T$.
\end{proof}

By Theorem~\ref{thm:O}, $(\nu_K,0,0)$ is not a certificate of value $\lamP$. The identity says that the
outside-window L\'evy tail removes from the symbol every frequency outside the window; apart from the
$\log|\tau|$ growth generated by $K$ near $u=0$, what remains is carried by frequencies in $[0,2T]$.
Section~\ref{sec:Z} identifies this remainder.

\subsection{An explicit certificate at \texorpdfstring{$T=\tfrac12\log3$}{T = (1/2) log 3}}\label{sec:cert}

At $T=\frac12\log3$ the window is $[-\log3,\log3]$, the contact set is $\NT=\{2\}$, and
$v_2=2\log2/\sqrt2=0.98025814\ldots=\Aconst$. Positivity of the complete form here is known by
\cite[Thm.~1.2]{Zhu2026}, and positivity of both forms also follows from Theorems~\ref{thm:log5A} and
\ref{thm:log5W} below. The point of the following theorem is the method: a proof by an explicit dual
feasible point, with a margin about $9.4\%$ of the numerical value $\lambda_{\mathrm A}(\frac12\log3)\approx9.448\cdot10^{-4}$.

\begin{theorem}[explicit certificate]\label{thm:C}
Let $T=\frac12\log3$, $d=2T/24=\frac{\log3}{24}$,
$u_j=jd$ for $24\le j\le460$ (so $u_{24}=\log3$ and $u_{460}=21.0567\ldots$), let
$a_{24},\dots,a_{460}$ be the $437$ rational numbers given as exact decimal strings in the
certificate data accompanying this paper, and let
$\beta_0=2.00256170925026\ldots$, $\beta_1=-0.00965443404485\ldots$. Put
$\nu=\sum_ja_j(\delta_{u_j}+\delta_{-u_j})$, so $\|\nu\|_{\TV}=2\sum|a_j|=90.75$, $\nuh(0)=1.7337$.
Then
\[
\sT(\tau)+\nuh(\tau)+\beta_0\gh0(\tau)+\beta_1\gh1(\tau)\ \ge\ c:=8.8577\cdot10^{-5}\qquad\text{for every }\tau\in\R,
\]
and therefore $\lamA\ge\lamP\ge8.8577\cdot10^{-5}>0$. The proof includes a computer verification
in interval arithmetic.
\end{theorem}

\begin{proof}
The function $f$ is even, so it suffices to take $\tau\ge0$. Split $[0,\infty)$ as $[0,\tau_1]$
with $\tau_1=30$, where $f$ itself is bounded below, and $[\tau_1,\infty)$, where
Lemma~\ref{lem:F1} with $N=2$ and $P=2\pi/d=137.2608\ldots$ reduces the inequality to the three
rows $\nuh+Y_k\ge c$, $k=0,1,2$, on the single period $[30,167.2608]$.

Each of the four functions ($f$ on $[0,\tau_1]$ and the three rows) is covered by the exhaustive
second-order procedure of Proposition~\ref{prop:TAY}. The second-derivative
bound is $M_2=6264.92$ for $f$, computed from the exact decimals; the three rows admit the smaller
constant $6262.99$, and the larger one is used for all four. The initial cell width is $0.02$, the
minimal width $w_{\min}=10^{-9}$; there is no cap on the number of cells and no cell is discarded.
The cover uses $2\,378\,593$ evaluations, leaves no unresolved interval, and consists of $20$
segments tiling the four intervals exactly.

The target in the cover is $8.8578\cdot10^{-5}$ minus a floating-point allowance
$\varepsilon_{\mathrm{fp}}=10^{-9}$. This allowance exceeds the a priori error of the double-precision
evaluator: argument rounding in the $437$ cosines contributes at most
$2\sum_j|a_j|\cdot u_{460}\cdot(\tau_1+NP)\cdot2^{-53}<10^{-10}$, the Taylor remainder at the minimal
width at most $M_2w_{\min}^2/8<10^{-15}$, and the summation of the $445$ terms together with the
evaluation of $\psi$ less than $10^{-12}$.

The constants $H_0,H_1$, $\psi'(\frac14)$, $k_0,k_1$ and the values of $m$ at the four minimizers
are enclosed in interval arithmetic at $40$ digits, $m$ through the recurrence
$\Rel\psi(z)=\Rel\psi(z+L)-\sum_{j<L}\Rel(z+j)^{-1}$ and the asymptotic series with Stieltjes'
remainder bound. Writing $h_k=\nuh+Y_k$ for the rows, the four minima $f(17.645)$, $h_0(37.815)$, $h_1$, $h_2(58.817)$
are strictly positive intervals with lower ends $\ge8.8578\cdot10^{-5}$; the smallest is
$h_0(37.8153)=8.85786\cdot10^{-5}$. The certificate data and the verification statistics accompany
the paper, and an independent replay from the decimal strings reproduces all
figures.
\end{proof}
\subsection{Verification on the bounded part}

\begin{proposition}[second-order cover]\label{prop:TAY}
Write $f=m+q+r$ with $q=-\sum_{n\in\NT}v_n\cos(\tau\log n)+\nuh$ ($\nu$ atomic) and
$r=\beta_0\gh0+\beta_1\gh1$, and put
$M_2=\sum_{n\in\NT}v_n\log^2n+2\sum_j|a_j|u_j^2+|\beta_0|H_0+|\beta_1|H_1$ with
$H_0=\int_{-2T}^{2T}u^2\cosh\frac u2\,du$, $H_1=\int_{-2T}^{2T}|u|^3\sinh\frac{|u|}2\,du$. Then
$\sup|q''+r''|\le M_2$, and on every cell $[x_0-\frac w2,x_0+\frac w2]\subset[0,\infty)$
\[
\min_{\text{cell}}f\ \ge\ m\!\big(x_0-\tfrac w2\big)+(q+r)(x_0)-|(q+r)'(x_0)|\,\tfrac w2-M_2\tfrac{w^2}8 .
\]
The same bound applies to each row $\nuh+Y_k$ of Lemmas~\ref{lem:F1}, \ref{lem:hybrid} with
$m(x+kP)$ and $-|\beta_i|k_i(x+kP)$ evaluated at the left endpoint. A cell is certified when the
right side exceeds the target and is bisected otherwise; near a contact of $f$ the admissible width
is $\approx\sqrt{8c/M_2}$, so the number of evaluations grows like $c^{-1/2}$ rather than the
$c^{-1}$ of a Lipschitz cover.
\end{proposition}
\begin{proof}
Taylor's formula with integral remainder for $q+r\in C^2$, and monotonicity of $m$ and of
$-|\beta_i|k_i$.
\end{proof}



\subsection{Mass against value}

\begin{theorem}[logarithmic exchange]\label{thm:Vp}
Let $B_*=\Aconst-C_0$ ($=7.621$ at $T\in(\frac12\log3,\log2]$).
\emph{(i)} For every feasible $\mu$ of (P), with $\varepsilon:=\mu(\{|\tau|\ge R\})$,
$\int\sT d\mu\ge-B_*+\varepsilon\,(m(R)-m(0))$, and $m(R)-m(0)=\log R+O(1)$.
\emph{(ii)} For every $\delta>0$ there is $C_1<\infty$ such that for all large $R$ and all
$\varepsilon\in[0,1]$ there is a feasible $\mu'_{\varepsilon,R}$ with
$\mu'_{\varepsilon,R}(\{|\tau|\ge R-\sqrt R\})\ge\varepsilon(1-o(1))$ and
$\int\sT d\mu'_{\varepsilon,R}\le\lamP+\delta+\varepsilon(\log R+C_1)+o(\varepsilon)$.
\emph{(iii)} Consequently no inequality $\int\sT d\mu\ge-B_*+\varepsilon\,\omega(R)$ with
$\omega(R)/\log R\to\infty$ holds on the feasible set: the exchange rate between escaped mass and
value is $\Theta(\log R)$, and converting the depth $B_*$ into a nonnegative value through escaped
mass alone requires $R\gtrsim e^{B_*/\varepsilon}$.
\end{theorem}
\begin{proof}
(i) $\sT\ge-B_*$ everywhere and $\sT\ge m(R)-\Aconst$ on $|\tau|\ge R$; split the integral.
(ii) Pick $F_0\in C_c^\infty(\IT)$ with $c_{F_0}\in\Kp$ and $\mathfrak a_T(F_0)\le\lamP+\delta$
(such $F_0$ are dense in the primitive class), $\mu_0=\mu_{F_0}$, which has a rapidly decaying
density. Set $\mu_{\varepsilon,R}=(1-\varepsilon)\mu_0+\frac\varepsilon2(\mu_0(\cdot-R)+\mu_0(\cdot+R))$;
then $\check\mu_{\varepsilon,R}=[(1-\varepsilon)+\varepsilon\cos Ru]\check\mu_0$ is still supported in
$[-2T,2T]$, so $\mu_{\varepsilon,R}$ is feasible for the relaxation without the two moment
constraints; its moment defect is $\varepsilon\,d(R)$ with $d(R)=(\int g_0\check\mu_0\cos Ru,\int g_1\check\mu_0\cos Ru)\to0$
by Riemann--Lebesgue. By Lemma~\ref{lem:W} there are $r_0>0$, $V^*<\infty$ and, for each $w$
with $|w|\le r_0$, a relaxed-feasible $\mu_w$ of value $\le V^*$ and moment $w$. Put
$\lambda=\varepsilon|d(R)|/(r_0+\varepsilon|d(R)|)$ and choose $\mu_w$ with
$w=-\frac{1-\lambda}\lambda\varepsilon d(R)$ (of norm $r_0$); then
$\mu'_{\varepsilon,R}=(1-\lambda)\mu_{\varepsilon,R}+\lambda\mu_w$ has moment $0$, is feasible for
(P), and its value exceeds that of $\mu_{\varepsilon,R}$ by at most
$\lambda(B_*+V^*)=o(\varepsilon)$. The value of $\mu_{\varepsilon,R}$ is
$(1-\varepsilon)\mathfrak a_T(F_0)+\varepsilon\int\sT(\tau\pm R)d\mu_0$, and by
Lemma~\ref{lem:mbounds}(ii) $\int m(\tau+R)d\mu_0\le\log R+\int|\tau|d\mu_0/R+O(1)$ while the prime
part is bounded by $\Aconst$. (iii) Apply (i) to $\mu'_{\varepsilon,R}$ and let $\varepsilon=1$.
\end{proof}

\subsection{The envelope threshold}

\begin{proposition}[envelope threshold]\label{prop:env}
\emph{(i)} Every certificate $(\nu,\beta)$ satisfies $f(\tau)\ge\log\frac\tau{2\pi}-\frac1\tau-\frac1{2\tau^2}-\Aconst-\|\nu\|_{\TV}-|\beta_0|k_0(\tau)-|\beta_1|k_1(\tau)$,
so $f\ge0$ holds automatically for $\tau\ge\tau_{\mathrm{env}}:=2\pi e^{\Aconst+\|\nu\|_{\TV}+o(1)}$.
\emph{(ii)} No bound of the form $f\ge m(\tau)-E(\tau)$ with $E$ depending on $\tau$ only through
an upper bound of the prime comb can replace $\Aconst$ by a smaller constant, because
$\sup_\tau\sum_{n\in\NT}v_n\cos(\tau\log n)=\Aconst$ (Lemma~\ref{lem:mbounds}(iii)).
\emph{(iii)} $\Aconst=\sum_{n<e^{2T}}2\Lambda(n)/\sqrt n=(4+o(1))e^{T}$ by the prime number theorem,
so $\tau_{\mathrm{env}}(T)$ grows doubly exponentially in $T$. The multiplier-free part alone gives
the threshold $2\pi e^{\Aconst}$, equal to $16.7$ and $59.5$ at $T=\frac12\log3$ and $\log2$ (the last
sign change of $\sT$ is at $10.53$ and $46.13$); for a certificate of total variation $\|\nu\|_{\TV}$
the threshold is $2\pi e^{\Aconst+\|\nu\|_{\TV}}$, about $10^{40}$ for the certificate of
\S\ref{sec:cert}, which is why its tail is handled by the periodic Lemma~\ref{lem:F1} and not by this
bound.
\end{proposition}
\begin{proof}
(i) Lemma~\ref{lem:mbounds} and $|\nuh|\le\|\nu\|_{\TV}$. (ii),(iii) as stated.
\end{proof}

Thus a pointwise argument can stop at $\tau_{\mathrm{env}}(T)$ and at no earlier point, and
$2\pi e^{\Aconst}$, doubly exponential in $T$, is the shortest interval such an argument can
reach; this is the dual counterpart of \cite[Thm.~1.4]{Zhu2026}. We do not know whether every
certificate must do work on all of $[0,2\pi e^{\Aconst}]$; the periodic tail lemmas show that a
lattice certificate is checked on $[0,\tau_1+(N+1)P]$, with $\tau_1+(N+1)P\approx442$ at $T=\frac12\log3$.

\section{The explicit formula on flat-top weights}\label{sec:Z}

So far every certificate has been a list of numbers produced by an optimizer. This section
identifies these numbers. Up to an arbitrarily small uniform correction, every certificate of the
complete Weil form is the Weil explicit formula evaluated on a family of modulated test functions,
and its value at $\tau$ is a smoothing of the zero measure at $\tau$.

\subsection{Flat-top weights}
\begin{definition}\label{def:W}
$\Wcl$ is the set of even measurable $w:\R\to\R$ with $w\equiv1$ on $[-2T,2T]$ and
$|w(u)|\le Ce^{-(\frac12+\delta)|u|}$ for $|u|\ge2T$, for some $C,\delta>0$. $\Wcl^{\mathrm{reg}}\subset\Wcl$
consists of the $w\in C^\infty(\R)$ with compact support. For $w\in\Wcl$ define the even measure
\begin{equation}\label{eq:nuw}
d\nu_w=\Big[\tfrac12\big(1-w(u)\big)K(|u|)+2w(u)\cosh\tfrac u2\Big]\mathbf 1_{|u|\ge2T}\,du
-\sum_{n\notin\NT}\tfrac{v_n}2\,w(\log n)\,\big(\delta_{\log n}+\delta_{-\log n}\big),
\end{equation}
and $f_w:=\sT+\nuh_w+2\gh0$.
\end{definition}
The measure $\nu_w$ is finite ($K$ is integrable on $[2T,\infty)$, $|w|\cosh\frac u2$ is
integrable by the decay condition, and $\sum_nv_n|w(\log n)|\le2C\sum\Lambda(n)n^{-1-\delta}$) and
supported in $\{|u|\ge2T\}$. Hence $(\nu_w,2,0)$ is a dual feasible point, and by
Propositions~\ref{prop:weak} and \ref{prop:P}, $W_T\ge\inf f_w$ on $L^2(\IT)$.

\begin{theorem}\label{thm:Z}
\emph{(i)} For every $w\in\Wcl$ and $\tau\in\R$,
\begin{equation}\label{eq:fw}
\begin{split}
f_w(\tau)=C_0+\int_0^\infty\big(1-w(u)\cos\tau u\big)K(u)\,du&-\sum_{n\ge2}v_n\,w(\log n)\cos(\tau\log n)\\
&+2\int_\R w(u)\cosh\tfrac u2\cos(\tau u)\,du,
\end{split}
\end{equation}
which is the right-hand side of the Weil explicit formula for the even test function
$h_\tau(u)=w(u)\cos(\tau u)$.
\emph{(ii)} For $w\in\Wcl^{\mathrm{reg}}$, with $\wh(t)=\int w(u)e^{itu}du$ (an even entire function
of exponential type),
\begin{equation}\label{eq:Z}
f_w(\tau)=\sum_\rho\wh(\tau-\gamma_\rho),\qquad\rho=\tfrac12+i\gamma_\rho,
\end{equation}
the sum running over the non-trivial zeros of $\zeta$ with multiplicity, and converging absolutely.
A zero on the critical line contributes $\wh(\tau-\gamma_\rho)$ with $\gamma_\rho$ real, so on
the zeros of the critical line $f_w$ is the smoothing of the zero-counting measure by the kernel $\wh$.
\end{theorem}
\begin{proof}
(i) By \eqref{eq:m}, $m(\tau)+\int_{2T}^\infty(1-w)K\cos\tau u\,du=C_0+\int_0^{2T}(1-\cos\tau u)K
+\int_{2T}^\infty K-\int_{2T}^\infty wK\cos\tau u=C_0+\int_0^\infty(1-w\cos\tau u)K$, using $w=1$ on
$[0,2T]$. The prime terms inside and outside the window combine into $-\sum_{n\ge2}v_nw(\log n)\cos(\tau\log n)$,
again because $w(\log n)=1$ for $n\in\NT$, and $2\gh0+2\int_{|u|\ge2T}w\cosh\frac u2\cos\tau u=2\int_\R w\cosh\frac u2\cos\tau u$.
(ii) For even $h\in C_c^\infty(\R)$ the Weil explicit formula \cite{Weil52}, \cite[Thm.~5.12]{IK}, in the
normalization of this paper, reads
\[
\sum_\rho\widehat h(\gamma_\rho)=2\int_\R h(u)\cosh\tfrac u2\,du-\sum_{n\ge2}v_n\,h(\log n)+\frac1{2\pi}\int_\R m(t)\,\widehat h(t)\,dt .
\]
Apply it to $h_\tau=w\cos(\tau\,\cdot)$, whose transform is $\frac12[\wh(t-\tau)+\wh(t+\tau)]$. The
zero sum is $\sum_\rho\wh(\tau-\gamma_\rho)$, since $\wh$ is even and the multiset $\{\gamma_\rho\}$ is
invariant under $\gamma\mapsto-\gamma$ ($\rho\mapsto1-\rho$). The polar and prime terms are the
last two terms of \eqref{eq:fw}. For the archimedean term insert \eqref{eq:m}: by Fourier
inversion $\frac1{2\pi}\int\wh(t\mp\tau)\,dt=w(0)=1$ and
$\frac1{2\pi}\int\wh(t\mp\tau)\cos(tu)\,dt=w(u)\cos(\tau u)$, and Fubini is justified by
$\int|\wh(t-\tau)|\int_0^\infty\min(2,t^2u^2)K(u)\,du\,dt<\infty$; the result is
$C_0+\int_0^\infty(1-w(u)\cos\tau u)K(u)\,du$. Absolute convergence of the zero sum holds because $\wh$ is
Schwartz on horizontal lines of the strip $|\Iml z|\le\frac12$, uniformly, and $N(t)\ll t\log t$.
\end{proof}


\begin{corollary}[certificates are statements about zeros]\label{cor:zeros}
If $w\in\Wcl^{\mathrm{reg}}$ and the elementary function \eqref{eq:fw} satisfies $f_w\ge c$ on $\R$,
then $\sum_\rho\wh(\tau-\gamma_\rho)\ge c$ for every real $\tau$, the sum including any zeros off the
critical line through $\wh$ at complex arguments. By Theorem~\ref{thm:finZ} below, this inequality
cannot be verified on any finite list of zeros: every truncation of the sum to finitely many
ordinates takes negative values.
\end{corollary}
\begin{proof}
Immediate from Theorem~\ref{thm:Z}(ii), the zero sum being absolutely convergent.
\end{proof}

\begin{remark}\label{rem:Zcheck}
Numerically, both sides of \eqref{eq:fw}--\eqref{eq:Z} agree for $C^2$ flat-top weights at eight
windows in $[\frac12\log3,1.20]$, the zero side computed from the first $10^5$ ordinates of Odlyzko's
table and the prime side from the elementary right-hand side of \eqref{eq:fw}, to between
$2.6\cdot10^{-7}$ and $2.8\cdot10^{-6}$ (double precision, not certified). At $T=\frac12\log3$ the
truncations of the zero side to the first $10$, $10^2$ and $10^3$ ordinates leave residuals
$4.0\cdot10^{-4}$, $6.0\cdot10^{-6}$ and $1.2\cdot10^{-7}$, of the orders that Lemma~\ref{lem:smooth}(iii)
predicts.
\end{remark}

\begin{theorem}[exhaustion by flat-top weights]\label{thm:Zexh}
Let $T\ge\frac12\log2$. \emph{(i)} $\sup_{w\in\Wcl^{\mathrm{reg}}}\inf_{\tau\in\R}f_w(\tau)=\lamW$.
\emph{(ii)} $\sup\{\inf_\tau[f_w+(\beta_0-2)\gh0+\beta_1\gh1]:\ w\in\Wcl^{\mathrm{reg}},\beta\in\R^2\}=\lamP$.
\end{theorem}
Thus every certificate of the complete Weil form is a flat-top explicit-formula certificate up to
an arbitrarily small uniform correction, and every certificate of the primitive form is one plus a
correction in the two primitivity multipliers.
\begin{proof}
``$\le$'' is weak duality in both cases. For ``$\ge$'' in (i): the proof of Theorem~\ref{thm:D1}
with the two $\widehat g$-constraints removed gives $\lamW=\sup\inf[\sT+2\gh0+\nuh]$ over finite
atomic $\nu$ outside the window, and the argument of Theorem~\ref{thm:E} (with $\beta=(2,0)$
fixed) replaces the atoms by a bounded even density $h$ with compact support in $\{|u|\ge2T\}$,
losing at most $\varepsilon/2$. It remains to realize $h\,du$ as $\nu_w$ up to total variation
$\varepsilon/2$, with $w\in\Wcl^{\mathrm{reg}}$; the transform then moves by at most $\varepsilon/2$
uniformly. On $u\ge2T$ the denominator $d_*(u)=2\cosh\frac u2-\frac12K(u)$ is positive and increasing
($d_*(0.3)>0$ and $d_*'>0$ on $(0,\infty)$), because $2T\ge\log2>0.3$; put
$w_h(u)=\big(h(u)-\tfrac12K(u)\big)/d_*(u)$ for $u\ge2T$, extended evenly and by $1$ on the window.
Then $|w_h|\ll e^{-|u|}$ and the absolutely continuous part of $\nu_{w_h}$ is exactly $h\,du$, while
its atomic part is $-\sum_{n\notin\NT}\frac{v_n}2w_h(\log n)(\delta_{\pm\log n})$, of total variation
$\ll\sum_n\Lambda(n)n^{-3/2}<\infty$. Now modify, in this order: (a) replace $h$ on $[2T,2T+\eta]$
by a smooth function matching $d_*(u)+\frac12K(u)$ to infinite order at $2T$ (so that $w_h\equiv1$
smoothly across the window edge); (b) mollify $h$ and cut $w_h$ off smoothly beyond a large $U_1$
(where $w_h\ll e^{-u}$; the a.c.\ part of $\nu_w$ becomes $\frac12K$ there, of total variation
$\ll e^{-U_1/2}$, and all atoms with $\log n>U_1$ disappear because $w=0$ there); (c) for the finitely
many $n\notin\NT$ with $\log n\le U_1$, replace $h$ by $\frac12K$ on disjoint intervals of small total
length around $\log n$ (so that $w_h$ vanishes there and the remaining atoms disappear). Each step
changes $\nu_{w_h}$ by an arbitrarily small total variation, and the result is in
$\Wcl^{\mathrm{reg}}$. (ii) is the same argument applied to the certificates of Theorem~\ref{thm:E}
with their own $\beta$.
\end{proof}

\subsection{Rigidity, smoothness, and finite zero sets}

\begin{lemma}[rigidity of flat-top weights]\label{lem:Rrig}
Let $w\in L^1(\R)$ be even and continuous with $w\equiv1$ on $[-2T,2T]$ for some $T>0$. Then
$\wh\ge0$ on $\R$ is impossible. Equivalently, if an even continuous $w\in L^1$ with $\wh\ge0$
satisfies $w\equiv1$ on an interval around the origin, then $w\equiv1$ on $\R$, which contradicts
integrability.
\end{lemma}
\begin{proof}
Suppose $w\in L^1$ is continuous with $\wh\ge0$. Then $\wh\in L^1$ and the inversion formula
$w(u)=\frac1{2\pi}\int\wh(t)e^{-itu}dt$ holds pointwise: this is the classical summability argument
(Gauss--Weierstrass or Fej\'er means converge to $w$ uniformly on compacta, and monotone convergence
applied to the nonnegative $\wh$ against the approximate identity gives $\int\wh=2\pi w(0)<\infty$;
see \cite[Ch.~VI]{Katznelson}). Hence $\wh\,dt/2\pi$ is a finite positive measure whose Fourier
transform is $w$, so by Bochner's theorem $w$ is positive definite, with $w(0)=1$. For a positive
definite $\phi$ the nonnegativity of the $3\times3$ Gram matrix at the points $0,-u,a$ gives the
standard inequality
\[
|\phi(u+a)-\phi(u)|^2\ \le\ 2\phi(0)\big(\phi(0)-\Rel\phi(a)\big)\qquad(u,a\in\R).
\]
Take any $a$ with $0<|a|\le2T$: then $w(a)=1=w(0)$, the right side vanishes, and $w(u+a)=w(u)$ for
every $u$. So $w$ is periodic with every period in $(0,2T]$; being continuous, it is constant, equal
to $w(0)=1$. A constant function is not in $L^1(\R)$.
\end{proof}

\begin{corollary}[positivity is never termwise]\label{cor:nontermwise}
For every even continuous $w\in\Wcl$ the kernel $\wh$ takes negative values. Consequently the
inequality $f_w=\sum_\rho\wh(\tau-\gamma_\rho)\ge0$ for $w\in\Wcl^{\mathrm{reg}}$ always expresses
cancellation across the zeros, never the termwise positivity of a kernel: there is no flat-top $w$
with $\wh\ge0$.
\end{corollary}
\begin{proof}
Lemma~\ref{lem:Rrig} applied to $w$, which lies in $L^1$ by the decay hypothesis of
Definition~\ref{def:W}, is even and equals $1$ on the window. A merely measurable $w$ with $\wh\ge0$
agrees almost everywhere with a continuous function, so nothing is lost.
\end{proof}

The following elementary lemma extends Theorem~\ref{thm:Z}(ii), stated for $w\in C_c^\infty$, to the weights $w\in C^k_c$ that arise in finite computations, and bounds the contribution of far zeros.

\begin{lemma}[smoothness and far zeros]\label{lem:smooth}
Let $w\in\Wcl$ be supported in $[-U,U]$, and let $k\ge0$ be an integer such that $w^{(k)}$ is
continuous on $\R$ (in particular across the window edge) and $w^{(k+1)}$ is of bounded variation,
with total variation $V=V(w^{(k+1)})$. Then:
\emph{(i)} for every $z\in\C$, $|\wh(z)|\le V\,e^{U|\Iml z|}|z|^{-k-2}$; in particular
$\wh(t)=O(|t|^{-k-2})$ on $\R$.
\emph{(ii)} The zero sum $\sum_\rho\wh(\tau-\gamma_\rho)$ converges absolutely and
Theorem~\ref{thm:Z}(ii) holds verbatim for such $w$; thus \eqref{eq:Z} is valid on
$\Wcl\cap C_c^k$ for every $k\ge0$, not only on $\Wcl^{\mathrm{reg}}$.
\emph{(iii)} Uniformly for $|\tau|\le H/2$,
\[
\sum_{|\gamma_\rho|>H}\big|\wh(\tau-\gamma_\rho)\big|\ \ll_{U}\ V\,H^{-k-1}\log H .
\]
\emph{(iv)} A zero $\rho=\frac12+i\gamma_\rho$ off the critical line, necessarily with
$|\Iml\gamma_\rho|\le\frac12$, contributes to the sum at most $e^{U/2}$ times the bound that
\emph{(i)} gives for a zero on the line with the same real ordinate.
\end{lemma}
\begin{proof}
(i) Integrating by parts $k+2$ times, the boundary terms vanishing because $w^{(j)}$ is continuous
and compactly supported for $j\le k$,
$\wh(z)=\int w(u)e^{izu}du=(iz)^{-(k+2)}\int e^{izu}\,dw^{(k+1)}(u)$, and $|e^{izu}|\le e^{U|\Iml z|}$
on $[-U,U]$. (ii) By (i) with $k\ge0$ we have $\wh(t)\ll V(1+|t|)^{-2}$ on $\R$ and
$|\wh(t+iy)|\ll Ve^{U/2}(1+|t|)^{-2}$ for $|y|\le\frac12$, while the number of zeros with
$|\gamma_\rho|\in[t,t+1]$ is $O(\log(2+t))$ by $N(t)\ll t\log t$; hence the sum converges absolutely,
and every step of the proof of Theorem~\ref{thm:Z}(ii) (the explicit formula for even $h$ of bounded
variation with compact support, and the Fubini justification, which used only
$\int|\wh(t-\tau)|\int_0^\infty\min(2,t^2u^2)K(u)du\,dt<\infty$) applies unchanged.
(iii) With $|\tau|\le H/2$ and $|\gamma_\rho|>H$ one has $|\tau-\gamma_\rho|\ge|\gamma_\rho|/2$,
so by (i) and partial summation against $N(t)=\frac t{2\pi}\log\frac t{2\pi e}+O(\log t)$,
$\sum_{|\gamma_\rho|>H}|\wh(\tau-\gamma_\rho)|\ll Ve^{U/2}\int_H^\infty t^{-k-2}\log t\,dt
\ll_UV H^{-k-1}\log H$. (iv) is the factor $e^{U|\Iml z|}$ of (i) at $|\Iml z|\le\frac12$.
\end{proof}

For $k=-1$, that is, a weight with a jump at the edge such as the box, the same computation gives
only $\wh(t)=O(|t|^{-1})$; the zero sum then converges conditionally, and (iii) becomes a statement
about oscillation rather than an absolute bound. For this reason the flat-top weights used in
practice are continuous across $u=\pm2T$.

\begin{lemma}[finite zero sets and positive definiteness]\label{lem:fdict}
Let $T>0$, let $Z\subset\R\setminus\{0\}$ be a finite symmetric multiset, and let $w:\R\to\R$ be
even, continuous and integrable with $w\equiv1$ on $[-2T,2T]$. Put
$g_Z(u)=\sum_{\gamma\in Z}e^{i\gamma u}=\sum_{\gamma\in Z}\cos(\gamma u)$ and $\phi:=w\,g_Z$. Then
$\phi$ is continuous and integrable, $\widehat\phi(\tau)=\sum_{\gamma\in Z}\wh(\tau-\gamma)$,
$\phi(0)=|Z|$, $\phi\equiv g_Z$ on $[-2T,2T]$, and
\[
\sum_{\gamma\in Z}\wh(\cdot-\gamma)\ \ge\ 0\ \text{ on }\R
\qquad\Longleftrightarrow\qquad
\phi\ \text{ is positive definite.}
\]
\end{lemma}
\begin{proof}
Continuity, $|\phi|\le|Z|\,|w|\in L^1$ and the values on the window are immediate from $w\equiv1$
there; $\widehat\phi=\wh*\mu_Z$ with $\mu_Z=\sum_{\gamma\in Z}\delta_\gamma$ finite is the
convolution theorem. ($\Leftarrow$) If $\phi$ is positive definite, its spectral measure is
nonnegative, and $\phi\in L^1\cap C$ forces that measure to be $\widehat\phi\,d\tau/2\pi$, so
$\widehat\phi\ge0$. ($\Rightarrow$) If $\widehat\phi\ge0$, the summability argument of
Lemma~\ref{lem:Rrig} (Gauss--Weierstrass means, then monotone convergence against the nonnegative
$\widehat\phi$) gives $\int\widehat\phi=2\pi\phi(0)<\infty$, so $\widehat\phi\,d\tau/2\pi$ is a
finite positive measure with Fourier transform $\phi$, and $\phi$ is positive definite by Bochner.
\end{proof}

\begin{theorem}[finite zero sets are determinate]\label{thm:finDet}
Let $T>0$, let $Z\subset\R\setminus\{0\}$ be a finite symmetric multiset with $|Z|=\ell\ge2$
(multiplicities allowed), and put $g_Z(u)=\sum_{\gamma\in Z}\cos(\gamma u)$. Let $\nu$ be a
finite positive symmetric measure on $\R$ with $\nuh=g_Z$ on $(-2T,2T)$. Then
$\nu=\mu_Z=\sum_{\gamma\in Z}\delta_\gamma$.
\end{theorem}
No decay hypothesis on $\nu$ and no product structure is assumed: the extension problem for a
finite zero set is determinate.
\begin{proof}
Since $\nu$ is finite, $\phi:=\nuh$ is continuous, bounded, real, even and positive definite.

\emph{Step 1 (a lattice inside the window).} Let $w_1,\dots,w_r$ be the distinct elements of
$Z$, with multiplicities $m_1,\dots,m_r\ge1$, so $\sum_jm_j=\ell$ and $r\le\ell$. Put $z_j=e^{iw_ja}$ for
a step $a>0$, so that $g_Z(ka)=\sum_{j\le r}m_jz_j^{\,k}$ for all $k\in\Z$. Since $Z$ is symmetric
with multiplicity, the list $(z_j,m_j)$ is closed under conjugation. Each condition $z_j=z_{j'}$
($j\ne j'$) excludes a discrete set of $a$, so let $E\subset(0,2T/\ell)$ be that exceptional discrete
set and fix $a\in(0,2T/\ell)\setminus E$, so that $z_1,\dots,z_r$ are pairwise distinct. As $r\le\ell$ we
have $ra\le\ell a<2T$, hence $\phi(ka)=g_Z(ka)$ for $|k|\le r$.

\emph{Step 2 (a rank-$r$ Toeplitz block).} Put $c_k:=\phi(ka)\in\R$ ($\phi$ is real and even) and
$A:=[c_{i-j}]_{0\le i,j\le r}$. By Step 1, $A=V^*DV$ with $V_{j,i}=z_j^{\,i}$ and
$D=\operatorname{diag}(m_1,\dots,m_r)$, $m_j>0$; $V$ has rank $r$ because the $z_j$ are distinct, so
$\operatorname{rank}A=r$ and
\[
\ker A=\Big\{v\in\C^{r+1}:\ \textstyle\sum_{i=0}^{r}v_iz_j^{\,i}=0\ (j\le r)\Big\}=\C\cdot v,
\]
$v$ being the coefficient vector of $P(z)=\prod_{j\le r}(z-z_j)$. The root set is closed under
conjugation, so $v\in\R^{r+1}$, $v_r=1$ and $|v_0|=\prod_j|z_j|=1\ne0$.

\emph{Step 3 (the recursion forced by positive definiteness).} Fix $m>r$. The Gram matrix of $\phi$
at $\{0,a,\dots,ra,ma\}$ is positive semidefinite; in block form it is
$M=\begin{psmallmatrix}A&b\\b^*&c_0\end{psmallmatrix}$ with $b_i=\phi(ia-ma)=c_{m-i}$. For
$v\in\ker A$ and $x=(v,t)$ one has $x^*Mx=2\Rel(\bar t\,b^*v)+|t|^2c_0\ge0$ for every $t\in\C$;
choosing the phase of $t$ adversely and letting $|t|\to0$ gives $b^*v=0$, i.e.
\begin{equation}\label{eq:frec}
\sum_{i=0}^{r}v_i\,c_{m-i}=0 .
\end{equation}
The sequence $k\mapsto g_Z(ka)$ satisfies the same relation: $\sum_iv_iz^{m-i}=z^{m-r}z^rP(1/z)$
vanishes at $z=z_j$ because $1/z_j=\bar z_j$ is again a root of $P$; multiplying by $m_j$ and
summing over $j$ gives $\sum_iv_ig_Z((m-i)a)=0$.

\emph{Step 4 (induction).} By Step 1, $c_k=g_Z(ka)$ for $0\le k\le r$. If $c_k=g_Z(ka)$ for all
$0\le k\le m-1$ with $m>r$, then in \eqref{eq:frec} every term with $i\ge1$ equals the
corresponding term for $g_Z$; subtracting the two relations and using $v_0\ne0$ gives
$c_m=g_Z(ma)$. Hence $\phi(ma)=g_Z(ma)$ for every integer $m\ge0$ and every
$a\in(0,2T/\ell)\setminus E$.

\emph{Step 5 (a variable step).} Fix an integer $m\ge1$. By Step 4 the identity
$\phi(ma)=g_Z(ma)$ holds for every $a\in(0,2T/\ell)$ outside the discrete set $E$; both sides are
continuous in $a$, so it holds for every $a\in(0,2T/\ell)$. Since
$\bigcup_{m\ge1}m\cdot(0,2T/\ell)=(0,\infty)$ and both functions are even, $\phi=g_Z$ on $\R$. The
spectral measure of a continuous positive definite function is unique (uniqueness of
Fourier--Stieltjes transforms), so $\nu=\mu_Z$.
\end{proof}

\begin{theorem}[no finite zero set carries a flat-top certificate]\label{thm:finZ}
Let $T>0$, let $Z\subset\R\setminus\{0\}$ be a finite symmetric multiset with $|Z|=\ell\ge2$
(multiplicities allowed), and let $w:\R\to\R$ be even, continuous and integrable with $w\equiv1$ on
$[-2T,2T]$ and $w(u)\to0$ as $|u|\to\infty$. Then
\[
f_{w,Z}(\tau):=\sum_{\gamma\in Z}\wh(\tau-\gamma)
\]
takes negative values on $\R$. In particular, for every continuous $w\in\Wcl$ and every
$H\ge\gamma_1=14.13\ldots$ such that all zeros with $|\Iml\rho|\le H$ lie on the critical line (by
\cite{PlattTrudgian2021}, every $H\le3\cdot10^{12}$), the truncation
$\sum_{|\gamma_\rho|\le H}\wh(\tau-\gamma_\rho)$ of a flat-top certificate is not nonnegative.
\end{theorem}
\begin{proof}
Suppose $f_{w,Z}\ge0$. By Lemma~\ref{lem:fdict} the function $\phi=w\,g_Z$ is continuous and
positive definite with $\phi(0)=\ell$, and $\widehat\phi=f_{w,Z}\in L^1$ with $\int f_{w,Z}=2\pi\ell$, so
$\nu_w:=f_{w,Z}\,d\tau/2\pi$ is a finite positive symmetric measure with $\nuh_w=\phi$, which equals
$g_Z$ on $(-2T,2T)$ because $w\equiv1$ there. By Theorem~\ref{thm:finDet}, $\nu_w=\mu_Z$, i.e.
$w\,g_Z=g_Z$ on $\R$. But along a lattice $\{ma\}$ as in Step~1 of that proof, the $z_j$ are
distinct points of the unit circle, so the Ces\`aro averages of $(z_j\bar z_{j'})^m$ vanish for
$j\ne j'$ (a geometric sum) and
\[
\lim_{M\to\infty}\frac1M\sum_{m=1}^{M}|g_Z(ma)|^2=\sum_{j\le r}m_j^2\ \ge\ r\ \ge\ 1 ,
\]
whence $\limsup_m|g_Z(ma)|\ge1$, while $|w(ma)g_Z(ma)|\le\ell|w(ma)|\to0$, a contradiction. Hence
$f_{w,Z}\ge0$ is impossible; being continuous with $\int_\R f_{w,Z}=2\pi\ell>0$, it takes negative
values.
\end{proof}


The proof of Theorem~\ref{thm:finZ} does not locate the negative values: it contradicts positive
definiteness along a lattice in the variable $u$ and says nothing about $\tau$. Under additional
smoothness hypotheses, integration by parts locates them.

\begin{proposition}[the location of the negativity]\label{prop:loc}
Let $w\in\Wcl$ be supported in $[-U,U]$, let $k\ge0$, assume that $w^{(k)}$ is continuous on $\R$ and
that $w^{(k+1)}$ is of bounded variation on $\R$, absolutely continuous on $\R\setminus\{\pm U\}$, with
a jump $J\ne0$ at $u=U$, i.e.\ $J=w^{(k+1)}(U^+)-w^{(k+1)}(U^-)$. Let $Z\subset\R\setminus\{0\}$ be a
finite symmetric multiset with $S:=\sum_{\gamma\in Z}\cos(U\gamma)\ne0$. Put $c_k(x)=(-1)^{k/2+1}\cos x$
for even $k$ and $c_k(x)=(-1)^{(k-1)/2}\sin x$ for odd $k$. Then, as $t\to+\infty$,
\[
\wh(t)=\frac{2J\,c_k(Ut)}{t^{k+2}}+o\big(t^{-k-2}\big),\qquad
f_{w,Z}(\tau)=\frac{2JS\,c_k(U\tau)}{\tau^{k+2}}+o\big(\tau^{-k-2}\big),
\]
so $f_{w,Z}(\tau)<0$ for arbitrarily large $\tau$, at the rate $\tau^{-k-2}$ and the frequency $U$.
\end{proposition}
\begin{proof}
Integrating by parts $k+2$ times as in Lemma~\ref{lem:smooth}(i), the boundary terms vanishing because
$w^{(j)}$ is continuous and compactly supported for $j\le k$,
$\wh(t)=(-1/it)^{k+2}\int e^{itu}\,dw^{(k+1)}(u)$. The measure $dw^{(k+1)}$ is an $L^1$ density plus
atoms at $\pm U$; the density contributes $o(1)$ to the integral by the Riemann--Lebesgue lemma. Since
$w^{(k+1)}$ has parity $(-1)^{k+1}$, its jump at $-U$ is $(-1)^kJ$, so the atoms contribute
$J(e^{itU}+(-1)^ke^{-itU})$, that is $2J\cos(Ut)$ for even $k$ and $2iJ\sin(Ut)$ for odd $k$. Multiplying
by $(-1/it)^{k+2}=(-1)^k i^{-(k+2)}t^{-(k+2)}$ gives the first formula. For the second, sum over the
finitely many $\gamma\in Z$, use $(\tau-\gamma)^{-k-2}=\tau^{-k-2}(1+O(1/\tau))$, and
$\sum_\gamma c_k(U(\tau-\gamma))=S\,c_k(U\tau)$, because $Z$ is symmetric, so that
$\sum_\gamma\sin(U\gamma)=0$. Since $JS\ne0$ and $c_k(U\tau)$ equals $\pm1$ on two arithmetic
progressions of step $2\pi/U$, the leading term is $\le-2|JS|\tau^{-k-2}$ on one of them.
\end{proof}

Theorem~\ref{thm:finZ} gives the negativity under minimal hypotheses; Proposition~\ref{prop:loc} gives
its location and size under stronger ones. Neither statement contains the other.

\begin{corollary}[finite sets of zeros do not suffice]\label{cor:infinite}
Let $T>0$. Every flat-top certificate $f_w\ge0$, $w\in\Wcl^{\mathrm{reg}}$, is a statement about the complete,
infinite zero multiset of $\zeta$: no truncation of $\sum_\rho\wh(\tau-\gamma_\rho)$ to finitely many
zeros is nonnegative, for any window, any admissible weight and any truncation height $H\ge\gamma_1$ up
to which the zeros are known to lie on the critical line. A finite
table of zeros can evaluate $f_w$ to any accuracy on a bounded range (Lemma~\ref{lem:smooth}(iii)),
but cannot establish $f_w\ge0$. Quantitatively, if $f_w\ge c>0$ for some $w\in\Wcl\cap C_c^k$
supported in $[-U,U]$ with $V=V(w^{(k+1)})$, and $H$ is large enough that
$C_UVH^{-k-1}\log H\le c/2$, then the truncation $f_{w,H}$ is $\ge c/2$ on
$|\tau|\le H/2$ while being negative somewhere; hence all of its negativity lies at
$|\tau|\gtrsim H/2$.
\end{corollary}
\begin{proof}
The first assertion is Theorem~\ref{thm:finZ} applied to the multiset
$Z=\{\gamma_\rho:|\gamma_\rho|\le H\}$, which is finite, symmetric, free of $0$, real by the choice of
$H$, and has at least two elements since $H\ge\gamma_1$. The
evaluation statement and the quantitative one follow from Lemma~\ref{lem:smooth}(iii), which bounds
$|f_w-f_{w,H}|$ by $C_UVH^{-k-1}\log H$ uniformly for $|\tau|\le H/2$.
\end{proof}


In the opposite direction, a finite zero set does admit a dual certificate, but only a trivial one.

\begin{proposition}[finite zero sets]\label{prop:finuniv}
Let $Z\subset\R\setminus\{0\}$ be a finite symmetric multiset, $\mu_Z=\sum_{\gamma\in Z}\delta_\gamma$,
$g_Z=\widehat{\mu_Z}$, and $W_Z(h):=\sum_{\gamma\in Z}\widehat h(\gamma)$ for even Schwartz $h$. Then:
\emph{(i)} $W_Z(F\star\widetilde F)=\sum_{\gamma\in Z}|\Fh(\gamma)|^2\ge0$ for every Schwartz $F$, where
$\widetilde F(u)=\overline{F(-u)}$;
\emph{(ii)} $\mu_Z$ is the unique symmetric positive tempered measure $\mu$ with
$W_Z(h)=\int\widehat h\,d\mu$ for all even Schwartz $h$;
\emph{(iii)} $g_Z$ is continuous and positive definite, with spectral measure $\mu_Z$.
\end{proposition}
\begin{proof}
(i) $\widehat{F\star\widetilde F}=|\Fh|^2$. (ii) $W_Z(h)=\int\widehat h\,d\mu_Z$ by definition; if $\mu$ is
another such measure, then $\int\widehat h\,d\mu=\int\widehat h\,d\mu_Z$ for every even Schwartz $h$, and
both measures are symmetric, so they agree on all Schwartz functions (the odd part of a test function
integrates to $0$ against both); hence $\mu=\mu_Z$. (iii) is Bochner's theorem for the finite positive
measure $\mu_Z$, and uniqueness of Fourier--Stieltjes transforms.
\end{proof}

For a finite zero set the dual certificate is therefore the spectrum itself: positivity of $W_Z$ is
immediate and the representing measure contains no arithmetic information, while by
Theorem~\ref{thm:finZ} the one-sided smoothing $\sum_{\gamma\in Z}\wh(\tau-\gamma)\ge0$ is never
available for a flat-top weight.

The mechanism is the following. A finite $Z$ makes the Toeplitz forms $[g_Z((i-j)a)]$ of finite rank;
beyond that rank, positive semidefiniteness of any extension propagates the linear recursion of $P$
indefinitely, so the extension is unique and does not decay (Theorem~\ref{thm:finDet}), whereas a
flat-top weight must decay. For the zero set of $\zeta$ no such $P$ exists and the rank is never finite.

\subsection{The decay tail of a flat-top certificate}

Proposition~\ref{prop:env} bounds the tail of any certificate at $2\pi e^{\Aconst+\|\nu\|_{\TV}}$,
and Lemma~\ref{lem:T} at a threshold exponential in the atomic mass. Both ignore the decay of
$\nuh_w$. For a $C^2$ flat-top weight this decay is available. The following theorem extends
Lemma~\ref{lem:T}: the single term $2\sum_j|a_j|$ is split into an atomic part, which does not decay,
and an absolutely continuous part, which decays like $\tau^{-2}$. It is also the smooth analogue of
Lemma~\ref{lem:F1}, with one essential difference recorded in Remark~\ref{rem:Dgrade}.

\begin{theorem}[decay tail bound for flat-top weights]\label{thm:Dtail}
Let $T>0$ and let $w\in\Wcl\cap C^2(\R)$ (so $w$ is even, $w\equiv1$ on $[-2T,2T]$, and $w'$, $w''$
match across $u=\pm2T$; in particular $w'(2T)=w''(2T)=0$), with $w''$ of bounded variation, and
assume either \emph{(a)} $\supp w\subset[-U,U]$, or \emph{(b)} $|w^{(j)}(u)|\le Ce^{-s|u|}$ for
$j\le2$ with $s>\frac12$. Put $\Omega(u)=w(|u|)\cosh\frac u2$ and $R=(1-w)K$ on $[2T,\infty)$, and
\[
D_c(w)=4\int_0^\infty\Big|\big(w(u)\cosh\tfrac u2\big)''\Big|\,du,\qquad
D_K(w)=\int_{2T}^\infty\big|\big((1-w)K\big)''(u)\big|\,du,
\]
\[
D(w)=D_c(w)+D_K(w),\qquad E_{\mathrm{at}}(w)=\sum_{\log n>2T}v_n\,|w(\log n)| ,
\]
all four finite. Then, with $\Psi_w(\tau)=\int_{2T}^\infty(1-w)K\cos\tau u\,du$,
\begin{equation}\label{eq:Dsplit}
\nuh_w(\tau)+2\gh0(\tau)=2\widehat\Omega(\tau)+\Psi_w(\tau)-\sum_{\log n>2T}v_n\,w(\log n)\cos(\tau\log n),
\end{equation}
and
\[
|2\widehat\Omega(\tau)|\le D_c(w)\,\tau^{-2},\qquad |\Psi_w(\tau)|\le D_K(w)\,\tau^{-2}.
\]
Consequently, for $f=f_w+(\beta_0-2)\gh0+\beta_1\gh1$ and every $\tau>0$,
\[
f(\tau)\ \ge\ L_w(\tau):=\log\frac\tau{2\pi}-\frac1\tau-\frac1{2\tau^2}-\Aconst-E_{\mathrm{at}}(w)
-|\beta_0-2|k_0(\tau)-|\beta_1|k_1(\tau)-D(w)\,\tau^{-2},
\]
$L_w$ is nondecreasing on $(0,\infty)$, and positivity of $f$ on $\R$ reduces to the compact
interval $[0,\tau_{\mathrm{stop}}(w)]$, $\tau_{\mathrm{stop}}(w)=\min\{\tau:L_w(\tau)\ge0\}$.
\end{theorem}
\begin{proof}
\emph{The identity \eqref{eq:Dsplit}, and why the $\cosh$ part must be regrouped with $2\gh0$.}
By \eqref{eq:nuw} the absolutely continuous part of $\nu_w$ contributes
$\Psi_w(\tau)+4\int_{2T}^\infty w\cosh\frac u2\cos\tau u\,du$ and its atomic part contributes the
last sum of \eqref{eq:Dsplit}. Since $\gh0(\tau)=2\int_0^{2T}\cosh\frac u2\cos\tau u\,du$ and
$w\equiv1$ on the window, $4\int_{2T}^\infty w\cosh\frac u2\cos\tau u\,du+2\gh0(\tau)
=4\int_0^\infty w\cosh\frac u2\cos\tau u\,du=2\widehat\Omega(\tau)$. The regrouping is necessary: the
density of $\nu_w$ at $u=2T^+$ equals $2\cosh T\ne0$, so $\nu_w$ has a jump at the window edge and
one integration by parts leaves the irreducible term $-4\cosh T\sin(2T\tau)/\tau$. Thus $\nuh_w$
alone is only $O(\tau^{-1})$, and no bound $|\nuh_w|\le D\tau^{-2}$ can hold. By
Lemma~\ref{lem:gbounds}, $2\gh0(\tau)=4\cosh T\sin(2T\tau)/\tau+O(\tau^{-2})$ cancels that term
exactly, with the same coefficient and the opposite sign; this is visible at once in
$2\widehat\Omega$, since $\Omega$ has no edge.

\emph{The two bounds.} $\Omega$ is even and, by the hypotheses, globally $C^2$ on $\R$ with $\Omega''$ of
bounded variation: on the window $\Omega=\cosh\frac u2$, and $w,w',w''$ match across $\pm2T$. In case
(a) $\Omega$ is compactly supported; in case (b) $|\Omega^{(j)}|\le C'e^{-(s-\frac12)|u|}$ for $j\le2$ is
integrable, which is exactly where $s>\frac12$ is needed. Two integrations by parts, all boundary
terms vanishing, give $\widehat\Omega(\tau)=-\tau^{-2}\int_\R e^{i\tau u}\,d\Omega'(u)$, so
$|2\widehat\Omega|\le2V(\Omega')\tau^{-2}=4\int_0^\infty|\Omega''|\,du\cdot\tau^{-2}=D_c(w)\tau^{-2}$ by evenness.
For $R$, extend by $0$ to $(-2T,2T)$ and evenly: $R(2T)=(1-w(2T))K(2T)=0$ and
$R'(2T)=-w'(2T)K(2T)+(1-w(2T))K'(2T)=0$ because $w\equiv1$ on the window and $w\in C^1$, so the
extension is $C^1$ on $\R$ with $R'$ of bounded variation, $V(R')=2\int_{2T}^\infty|R''|$ (in case
(a), $R=K$ beyond $U$, and $K,K',K''$ are integrable at infinity; in case (b),
$|R''|\le(|w''|+2|w'|+|1-w|)\,(K+|K'|+|K''|)$ with $K+|K'|+|K''|\le Ce^{-u/2}$ on $[2T,\infty)$). As $\Psi_w=\frac12\widehat R$,
two integrations by parts give $|\Psi_w|\le\frac12V(R')\tau^{-2}=D_K(w)\tau^{-2}$.

\emph{The atoms carry no decay}, and are bounded only by their mass:
\[
\Big|\sum_{\log n>2T}v_nw(\log n)\cos(\tau\log n)\Big|\le E_{\mathrm{at}}(w),
\]
the sum being finite in case
(a) and convergent in case (b) because $\sum_n\Lambda(n)n^{-\frac12-s}<\infty$ for $s>\frac12$,
the same threshold as above.

\emph{The lower bound.} $\sT(\tau)\ge\log\frac\tau{2\pi}-\frac1\tau-\frac1{2\tau^2}-\Aconst$ by
Lemma~\ref{lem:mbounds}(ii),(iii), the two multiplier terms are bounded by Lemma~\ref{lem:gbounds},
and the three pieces of \eqref{eq:Dsplit} by the above. Each term of $L_w$ is nondecreasing, hence so
is $L_w$, and $L_w(\tau_{\mathrm{stop}})\ge0$ gives $f\ge0$ on $[\tau_{\mathrm{stop}},\infty)$.
\end{proof}

\begin{remark}\label{rem:Dgrade}
The exponent $\tau^{-2}$ is the one Lemma~\ref{lem:smooth}(i) gives at $k=0$ after the two extra
derivatives, and it cannot be improved without $w\in C^3$: integrating by parts once more, a jump of
$w'''$ at the support endpoint contributes a genuine $\tau^{-2}$ to $\widehat\Omega$. The bound on
the atoms is sharp in the worst case by Lemma~\ref{lem:mbounds}(iii). This is where the method
differs from Lemma~\ref{lem:F1}: a lattice $\nuh$ is $P$-periodic, so the comb can be evaluated
exactly on $N$ periods and replaced by its infimum beyond, whereas a flat-top $\nuh_w$ is not
periodic and its comb can only be bounded by its mass. Theorem~\ref{thm:Dtail} handles the absolutely
continuous part by decay. We do not know how to combine the two techniques; the remaining
exponential dependence on the atomic mass comes from this gap.
\end{remark}

\begin{corollary}[comb-vanishing weights]\label{cor:comb}
Assume in addition $\supp w\subset[-U,U]$ and $w(\log n)=0$ for every prime power $n$ with
$2T<\log n\le U$ (a finite set of linear conditions). Then $E_{\mathrm{at}}(w)=0$, and for
$D(w)\le D_0$,
\[
\tau_{\mathrm{stop}}(w)\ \le\ \max\Big\{2\pi e^{\Aconst+1+|\beta_0-2|k_0(1)+|\beta_1|k_1(1)},\ \sqrt{D_0}+2\Big\},
\]
since for $\tau$ beyond the first entry one has $\log\frac\tau{2\pi}-\Aconst-|\beta_0-2|k_0-|\beta_1|k_1\ge1$
and then $L_w(\tau)\ge0$ as soon as $\tau^2\ge D_0+O(\tau)$. In particular
$\tau_{\mathrm{stop}}=O(D(w)^{1/2})$ as $D(w)\to\infty$: the verification interval is polynomial,
not exponential, in the size of the weight, compared with $2\pi e^{\Aconst+E_{\mathrm{at}}(w)}$ in
general and $2\pi e^{\Aconst+\|\nu\|_{\TV}}$ in Proposition~\ref{prop:env}.
\end{corollary}


\begin{corollary}[bounded profiles]\label{cor:bprof}
Let $w$ satisfy the hypotheses of Theorem~\ref{thm:Dtail}(a), with $0\le w\le1$ and $\supp w\subset[-U,U]$,
and put $\kappa_0(T)=K(2T)$, $\kappa_1(T)=|K'(2T)|$. Then
\[
D_c(w)\le4\cosh\tfrac U2\Big(\|w''\|_1+\|w'\|_1+\tfrac14\|w\|_1\Big),\qquad
D_K(w)\le\kappa_0(T)\|w''\|_1+2\kappa_1(T)\|w'\|_1+2\kappa_1(T),
\]
so $D(w)=O(e^{U/2})$ with a constant depending on $\|w\|_{W^{2,1}}$ and $T$. Boundedness, however,
controls only $D$: it bounds $E_{\mathrm{at}}(w)$ by nothing better than $\sum_{2T<\log n\le U}v_n$, which
grows like $4e^{U/2}$. Only the comb conditions of Corollary~\ref{cor:comb} control the atoms.
\end{corollary}
\begin{proof}
On $[0,U]$, $(w\cosh\frac u2)''=w''\cosh\frac u2+w'\sinh\frac u2+\frac14w\cosh\frac u2$, and
$|\sinh|\le\cosh\le\cosh\frac U2$ there; integrating over $[0,\infty)$ gives the bound for $D_c$. On
$[2T,\infty)$, $((1-w)K)''=-w''K-2w'K'+(1-w)K''$. By Lemma~\ref{lem:K}, $K$ is completely monotone, so
$K$ and $|K'|$ are decreasing and $K''\ge0$; hence $|((1-w)K)''|\le\kappa_0|w''|+2\kappa_1|w'|+K''$
(using $0\le1-w\le1$), and $\int_{2T}^\infty K''=-K'(2T)=\kappa_1$. For the atoms, $0\le w\le1$ gives
$E_{\mathrm{at}}(w)\le\sum_{2T<\log n\le U}v_n$, with equality when $w=1$ at those points, and
$\sum_{n\le x}2\Lambda(n)n^{-1/2}=4\sqrt x\,(1+o(1))$ by the prime number theorem and partial summation.
\end{proof}

\begin{remark}\label{rem:comblp}
Numerically, the comb-vanishing class of Corollary~\ref{cor:comb} does not produce a certificate in the
following finite family: compactly supported $C^2$ cubic splines on a lattice of $48$ cells with $U=6$,
matched to $1$ across the window edge, with the comb equalities imposed exactly and $\beta$ free. A
linear program over $\tau\in[0,200]$ returns negative minima at $T=\frac12\log3$, $0.58$, $0.62$ and
$\log2$, of sizes about $10^{-2}$ to $10^{-1}$ (double precision, not certified), and for the
resulting weights $\tau_{\mathrm{stop}}<200$, so the failure concerns all of $\R$. The binding minimum
sits at $\tau\approx4$--$8$, whereas without the comb conditions it lies near $\tau\approx43$. One
structural observation: consecutive prime powers near $e^u$ are spaced $\approx e^{-u}$ on the
logarithmic scale, so once the lattice step exceeds that spacing every added cell is consumed by the
conditions inside it, and comb vanishing forces an effectively bounded support.
\end{remark}

\subsection{The Kre\u{\i}n extension problem at one window}\label{subsec:KEP}
The dual variable of Section~\ref{sec:Z} is a positive measure prescribed only on the window. This
is an extension problem for a positive definite datum given on an interval, in the sense of
Kre\u{\i}n \cite{Krein1940}, \cite[Ch.~2]{Sasvari1994}. In this subsection we make the identification
at one window. The classical theory concerns continuous positive
definite functions, whereas the datum below is a distribution with atoms; we therefore import
nothing from it, and every statement is proved from Theorems~\ref{thm:Z} and \ref{thm:Zexh} and
elementary Fourier analysis.

\begin{definition}[the window datum, and $\mathrm{KEP}(T)$]\label{def:KEP}
Let $g$ be the even distribution defined on even $h\in C_c^\infty(\R)$ by
\begin{equation}\label{eq:gdatum}
\langle g,h\rangle=C_0\,h(0)+\int_0^\infty\big(h(0)-h(u)\big)K(u)\,du-\sum_{n\ge2}v_n\,h(\log n)
+2\int_\R h(u)\cosh\tfrac u2\,du,
\end{equation}
that is, the right-hand side of the explicit formula in the normalization of the proof of
Theorem~\ref{thm:Z}(ii): the L\'evy pair $C_0\delta_0-\frac12\mathrm{Pf}\,K(|u|)$, whose Fourier
transform is the archimedean multiplier $m$ of \eqref{eq:m}; the polar term $2\cosh\frac u2$; and
the prime atoms $-\frac{v_n}2\big(\delta_{\log n}+\delta_{-\log n}\big)$, $n\ge2$. For $T>0$ put
\[
\begin{aligned}
\mathfrak S(T)=\Big\{\nu\ \text{a positive symmetric tempered measure on }\R:\ \\
\langle\nuh-g,h\rangle=0\ \text{for every even }h\in C_c^\infty(-2T,2T)\Big\}.
\end{aligned}
\]
The \emph{Kre\u{\i}n extension problem} $\mathrm{KEP}(T)$ is \emph{solvable} if
$\mathfrak S(T)\ne\emptyset$, \emph{determinate} if $\#\mathfrak S(T)=1$, and \emph{indeterminate}
if $\#\mathfrak S(T)\ge2$.
\end{definition}
Symmetry is imposed because even test functions
determine only the even part of a representing measure, so without it every problem would be
indeterminate for a trivial reason. We define ``the datum is positive definite on the window'' to
mean solvability; by the Schwartz--Bochner theorem this is equivalent to saying that the restriction
of $g$ to $(-2T,2T)$ extends to a positive definite tempered distribution on $\R$.

\begin{lemma}[one negative square in the polar term]\label{lem:polar}
At any $n\ge2$ distinct points $t_1,\dots,t_n$ the Gram matrix of the kernel
$(s,t)\mapsto2\cosh\frac{s-t}2$ is $ab^{\mathsf T}+ba^{\mathsf T}$ with $a_i=e^{t_i/2}$ and
$b_i=e^{-t_i/2}$; it has rank $2$ and signature $(1,1)$. Hence $2\cosh\frac u2$ is positive definite
on no interval, and the solvability of $\mathrm{KEP}(T)$ is automatic at no window, however small.
\end{lemma}
\begin{proof}
$2\cosh\frac{t_i-t_j}2=e^{t_i/2}e^{-t_j/2}+e^{-t_i/2}e^{t_j/2}$, which is the stated rank-two matrix.
It vanishes on $\{a,b\}^\perp$ and on $\mathrm{span}\{a,b\}$ has eigenvalues $a\cdot b\pm|a||b|$;
$a\cdot b=n$ while $|a||b|>n$ strictly by Cauchy--Schwarz, the vectors $a,b$ not being parallel when
the $t_i$ are distinct and $n\ge2$. So one eigenvalue is positive and one negative.
\end{proof}
The datum \eqref{eq:gdatum} is therefore the sum of a positive definite part ($m\ge0$), a part with
exactly one negative square (the pole), and a negative part (the prime atoms). This is the reason a
polar multiplier $\beta_0$ must be carried at all (Proposition~\ref{prop:P}).

\begin{theorem}[a certificate is an absolutely continuous solution]\label{thm:KEP1}
Let $T\ge\frac12\log2$ and let $w\in\Wcl^{\mathrm{reg}}$ satisfy $f_w\ge c$ on $\R$.
\emph{(i)} If $c\ge0$ then $\nu_w:=f_w(\tau)\,d\tau/2\pi$ belongs to $\mathfrak S(T)$; in particular
$\mathrm{KEP}(T)$ is solvable, by an absolutely continuous solution.
\emph{(ii)} If $c>0$ then $\mathfrak S(T)$ contains a continuum of pairwise distinct absolutely
continuous measures; in particular $\mathrm{KEP}(T)$ is indeterminate.
\end{theorem}
\begin{proof}
(i) For $w\in\Wcl^{\mathrm{reg}}$ the product $w\cdot g$ is a tempered distribution and
Theorem~\ref{thm:Z} says $\widehat{w\cdot g}=f_w$, a continuous function with
$f_w(\tau)=\log|\tau|+O(1)$ (Lemma~\ref{lem:mbounds}, Theorem~\ref{thm:Dtail}), hence tempered;
Fourier inversion gives $\nuh_w=w\cdot g$, which agrees with $g$ on $(-2T,2T)$ because $w\equiv1$
there. Positivity of $\nu_w$ is $f_w\ge0$, and symmetry is the evenness of $w$ and $g$.
(ii) Fix an even $\chi\in C_c^\infty(\R)$, $\chi\ge0$, $\chi\not\equiv0$, with
$\supp\chi\subset\{|u|\ge\max(2T,\log2)\}$. For $\varepsilon\in\R$ put $w_\varepsilon=w+\varepsilon\chi
\in\Wcl^{\mathrm{reg}}$; then $f_{w_\varepsilon}=f_w+\varepsilon r_\chi$ with
\[
r_\chi(\tau)=2\int\chi(u)\cosh\tfrac u2\cos\tau u\,du-\sum_{n\ge2}v_n\chi(\log n)\cos(\tau\log n)
-\tfrac12\int_\R\chi(u)K(u)\cos\tau u\,du ,
\]
a bounded continuous function, say $|r_\chi|\le R$ ($\chi$ is compactly supported away from the
singularity of $K$, and $\sum_nv_n|\chi(\log n)|<\infty$). For $|\varepsilon|<c/R$ one has
$f_{w_\varepsilon}>0$, so $\nu_{w_\varepsilon}\in\mathfrak S(T)$ by (i). They are pairwise distinct:
$\nu_{w_\varepsilon}-\nu_{w_{\varepsilon'}}=\frac{\varepsilon-\varepsilon'}{2\pi}r_\chi\,d\tau$, and
$r_\chi\equiv0$ would force $\chi\cdot g=0$, whereas on $\supp\chi\subset\{|u|\ge\log2\}$ the
absolutely continuous part of $g$ is $\chi(u)d_*(u)$ with $d_*(u)=2\cosh\frac u2-\frac12K(u)>0$
(positivity and monotonicity of $d_*$ on $[0.3,\infty)$ are established in the proof of
Theorem~\ref{thm:Zexh}), which is nonnegative and not identically zero and which the singular prime
atoms cannot cancel.
\end{proof}

\begin{theorem}[solvability implies Weil positivity at that window]\label{thm:KEP2}
If $\mathfrak S(T)\ne\emptyset$ then $\lamW\ge0$.
\end{theorem}
\begin{proof}
Let $\nu\in\mathfrak S(T)$ and $F\in C_c^\infty(\IT)$ real, and put $h=F\star\widetilde F$ with
$\widetilde F(u)=\overline{F(-u)}$. Then $h$ is even, $\supp h\subset(-2T,2T)$ and
$\widehat h=|\Fh|^2\ge0$ is rapidly decreasing, so the pairing against the positive tempered measure
$\nu$ is legitimate. Since $h$ vanishes off the window, only the contacts $n\in\NT$ appear in
\eqref{eq:gdatum}, and comparison with \eqref{eq:Aform} and \eqref{eq:W} gives
$\langle g,h\rangle=\mathfrak a_T(F)+2\Rel\big(A(F)\overline{B(F)}\big)=W_T(F)$. By
Definition~\ref{def:KEP}, $W_T(F)=\int\widehat h\,d\nu=\int|\Fh|^2d\nu\ge0$. For a complex test, apply the real inequality to its real and imaginary parts and add, using the real even symbol and \eqref{eq:W}. Finally
$C_c^\infty(\IT)$ is a form core: for $F$ in the form domain, the dilates $F_\lambda(t)=F(t/\lambda)$,
$\lambda\uparrow1$, are supported in $\lambda\IT$ and converge to $F$ in $L^2$ with
$\Fh_\lambda(\tau)=\lambda\Fh(\lambda\tau)$, so that
$\int\log(2+|\tau|)|\Fh_\lambda-\Fh|^2\to0$ (the weight is doubling), and mollifying $F_\lambda$ at a
scale below $(1-\lambda)T$ keeps the support inside $\IT$ and converges in the same norm. Hence
$W_T\ge0$ on the whole form domain and $\lamW\ge0$.
\end{proof}

\begin{theorem}[the dichotomy at one window]\label{thm:KEPdich}
For every $T\ge\frac12\log2$,
\[
\lamW>0\ \Longrightarrow\ \mathrm{KEP}(T)\ \text{is indeterminate}\ \Longrightarrow\ \lamW\ge0 ,
\]
so certificates and indeterminacy coincide except on the set $\{T:\lamW=0\}$.
\end{theorem}
\begin{proof}
If $\lamW>0$, Theorem~\ref{thm:Zexh}(i) gives $w\in\Wcl^{\mathrm{reg}}$ with $\inf f_w>0$, and
Theorem~\ref{thm:KEP1}(ii) applies; the second implication is Theorem~\ref{thm:KEP2}.
\end{proof}

Off the set $\{T:\lamW=0\}$ the solutions of the extension problem and the certificates therefore
determine each other. A Nevanlinna-type parametrization of $\mathfrak S(T)$ in the indeterminate case
would describe the set of certificates, but it presupposes one known solution.


\section{The archimedean form on primitive tests}\label{sec:archstr}

With the primes deleted, the Weil form is an explicit convolution operator. This section studies
its thresholds, its negative index, and the number of even modes that the prime form can compete
with.

\subsection{Thresholds and nodal structure}\label{subsec:arch}

Put, for $F\in L^2(\IT)$,
\[
\Ainf(F)=2A(F)B(F)+\frac1{2\pi}\int_\R m(\tau)|\Fh(\tau)|^2\,d\tau
=2A(F)B(F)+C_0\|F\|^2+\Lev(\cF),
\]
the complete Weil form with the prime comb deleted (the second equality is Theorem~\ref{thm:A} with
$\NT=\emptyset$), and write $\PiT(F)=-\sum_{n\in\NT}v_n\Rel\cF(\log n)$ for the prime form (the real part is redundant for real $F$), so that
$W_T=\Ainf+\PiT$. Let $\laminf(T)$ be the bottom of $\Ainf$ on the unit sphere of $L^2(\IT)$ and
$\laminfP(T)$ its bottom on the unit sphere of $\PT$, where the pole term vanishes identically and
$\Ainf=C_0\|F\|^2+\Lev(\cF)$.

\begin{proposition}[two archimedean thresholds]\label{prop:arch}
\emph{(i)} $\Ainf$ does not depend on $T$ and the feasible sets increase with $T$, so both
$\laminf$ and $\laminfP$ are nonincreasing (for $\laminfP$ strictly, Lemma~\ref{lem:strict}).
\emph{(ii)} $\laminf(T)\le\laminfP(T)\le C_0+\varphi\big((3\pi/2T)^2\big)$; the first inequality is
inclusion of feasible sets and the second is the proof of Corollary~\ref{cor:mono} with the comb deleted. Hence, with
$T_0(0)=0.74921$ the root of $\varphi((3\pi/2T)^2)=-C_0$, both are nonpositive at $T=T_0(0)$ and negative for $T>T_0(0)$, and $\laminfP(T)\downarrow C_0$ as
$T\to\infty$, since $\Lev\ge0$ and $\varphi((3\pi/2T)^2)\downarrow0$.
\emph{(iii)} For $2T<\log2$ the contact set is empty, so $W_T=\Ainf$ and $\laminf(T)\ge0$ by the
known positivity of the complete form at autocorrelation support $\le\log2$
\textup{\cite{Yoshida1992,Bombieri2000}} (and, certified, \cite[Thm.~1.2]{Zhu2026}).
\emph{(iv)} Consequently the thresholds
$T^*_\infty=\sup\{T:\laminf(T)>0\}$ and $T^*_{\infty,\mathrm P}=\sup\{T:\laminfP(T)>0\}$ exist, and
\[
T^*_\infty\ \le\ T^*_{\infty,\mathrm P}\ \le\ 0.74921 .
\]
\end{proposition}
\begin{proof}
(i) $K$, $C_0$ and the two Tate moments are $T$-free, and an infimum over a growing set is
nonincreasing. (ii) $\PT\subset L^2(\IT)$, so the restricted infimum is the larger one; the proof of
Corollary~\ref{cor:mono} bounds the bottom of $\Lev$ on the codimension-two subspace $\PT$ by
$\lambda_3\le\varphi((3\pi/2T)^2)$ and never mentions the primes, and the comb contributes $0$.
(iii) and (iv) are immediate; a nonincreasing function changes sign once, and the thresholds are
suprema. No continuity is used. For $\laminf$ it holds by the proof of Lemma~\ref{lem:contT}, and for
$\laminfP$ it is not asserted.
\end{proof}
\begin{remark}
In (iii) the strict inequality $\laminf>0$ is not re-proved here, and that is the only reason
(iv) is stated for $T^*_\infty$ as an upper bound: strict positivity at $T=\frac12\log2$,
available in \cite{Yoshida1992}, would add the lower bound $\frac12\log2\le T^*_\infty$.
\end{remark}
\begin{remark}
The archimedean form by itself was studied by Connes and Consani \cite{ConnesConsani2021}, who prove
its positivity on a window by a Toeplitz/prolate-spheroidal compression and relate the archimedean
place to the trace formula. Proposition~\ref{prop:arch} looks at the same form from the other side:
it locates the widths beyond which the archimedean form alone stops being positive, first on
$L^2(I_T)$ and then on $\PT$, so that past $T^*_{\infty,\mathrm P}$ the prime form $\PiT$ has to
supply positivity rather than merely be dominated.
\end{remark}

The bottom eigenfunction of the archimedean form on $\PT$ is never of one sign, and this is forced
by the two Tate conditions alone.

\begin{proposition}[nodal structure of $\PT$]\label{prop:nodes}
Recall that $\mT(F)=C_0\|F\|^2+\Lev(\cF)$ is the archimedean form without the pole term on
$L^2(\IT)$, so that $\Ainf=\mT$ on $\PT$. Then:
\emph{(i)} the bottom of $\mT$ on all of $L^2(\IT)$ is a simple eigenvalue whose
eigenfunction can be taken a.e.\ strictly positive;
\emph{(ii)} say that $F\in L^2(\IT)$ has \emph{at most one sign change} if for some
$t_0\in[-T,T]$ one has $F\ge0$ a.e.\ on $(-T,t_0)$ and $F\le0$ a.e.\ on $(t_0,T)$, or the same
with the two inequalities exchanged (the endpoints $t_0=\pm T$ being the case of one sign).
Then every nonzero $F\in\PT$ has at least two sign changes in this sense; in particular no
nonzero element of $\PT$ is of one sign, and an even element has at least one nodal pair in
$(0,T)$;
\emph{(iii)} consequently the bottom eigenfunction of $\mT$ on $\PT$ is not of one sign,
and the Sturm index of $\Ainf|\PT$ is shifted by one nodal pair relative to $(i)$.
\end{proposition}
\begin{proof}
(i) For $F\in L^2(\IT)$ extended by zero, Theorem~\ref{thm:A} and Fubini give the pure-jump
Dirichlet form with killing
\[
\begin{aligned}
\Lev(\cF)&=\frac14\iint_{\IT\times\IT}K(|x-y|)\,|F(x)-F(y)|^2\,dx\,dy
+\frac12\int_{\IT}\kappa(x)|F(x)|^2\,dx,\\
\kappa(x)&=\int_{|y|>T}K(|x-y|)\,dy<\infty,
\end{aligned}
\]
obtained by splitting $2\Lev(\cF)=\int_\R K(|u|)\big(\|F\|^2-\Rel\cF(u)\big)du
=\frac12\iint_{\R^2}K(|x-y|)|\widetilde F(x)-\widetilde F(y)|^2$ into the square $\IT\times\IT$
and the two strips on which the zero extension $\widetilde F$ vanishes; the divergence of $K$ at
the origin is absorbed by $|F(x)-F(y)|^2$ on the form domain, and no $\delta$-kernel exists,
$\int_0^{2T}K=+\infty$. The form is therefore a regular Dirichlet form of pure-jump type with
nonnegative jump density $K$ and killing $\kappa$, up to the additive constant $C_0\|F\|^2$,
which only shifts the spectrum. By the Beurling--Deny criteria the semigroup $e^{-t\mT}$ is
positivity preserving, and since $K>0$ everywhere it is irreducible, hence positivity improving;
Perron--Frobenius then gives a simple, a.e.\ strictly positive ground state, provided the bottom
is an eigenvalue. It is: $m$ is increasing (Lemma~\ref{lem:mbounds}(i)) and $m(\tau)\to\infty$, so
$\int_{|\tau|>R}|\Fh|^2\le2\pi\,(m(R)-C_0)^{-1}\Lev(\cF)$ uniformly on form-bounded sets, while all
competitors are supported in the fixed compact $\IT$; by the Riesz--Kolmogorov criterion the form
domain embeds compactly in $L^2(\IT)$ and the resolvent is compact.
(ii) The pair $\{e^{t/2},e^{-t/2}\}$ is a Chebyshev system of order $2$ on $(-T,T)$: a nonzero
combination $c_1e^{t/2}+c_2e^{-t/2}$ vanishes only where $e^{t}=-c_2/c_1$, i.e.\ at most once. Let
$0\ne F\in\PT$ and suppose $F$ had at most one sign change, at $t_0$. Choose
$G=c_1e^{t/2}+c_2e^{-t/2}$ vanishing exactly at $t_0$ (any nonzero $G$ of the family if
$t_0=\pm T$), with the sign of $G$ matched to that of $F$ on each side. Then $FG\ge0$ a.e.\ and
$FG$ is not a.e.\ zero, so $\int FG\ne0$; but $F\in\PT$ means $\langle F,e^{t/2}\rangle=\langle
F,e^{-t/2}\rangle=0$, whence $\int FG=0$, a contradiction. For even $F$ the two sign changes are
symmetric, so at least one lies in $(0,T)$.
(iii) Immediate from (ii).
\end{proof}

\emph{Parity sectors.} For even $F$ the two Tate moments coincide,
$A(F)=B(F)=\int F(t)\cosh\frac t2\,dt$, and for odd $F$ one has $A(F)=-B(F)=\int F(t)\sinh\frac t2\,dt$.
Write $\PT^{\mathrm{ev}}$ and $\PT^{\mathrm{odd}}$ for the even and the odd elements of $\PT$
$\lambda_{\mathrm{ev}}(T)$ and
$\lambda_{\mathrm{odd}}(T)$ for the bottoms of $\Ainf=\mT$ on their unit spheres,
$n_{\mathrm{ev}}(T)$ and $n_{\mathrm{odd}}(T)$ for the numbers of negative eigenvalues of $\Ainf$ on them,
and $T^*_{\mathrm{ev}}=\sup\{T:\lambda_{\mathrm{ev}}(T)>0\}$, $T^*_{\mathrm{odd}}=\sup\{T:\lambda_{\mathrm{odd}}(T)>0\}$.
By Lemma~\ref{lem:parity} below, $\laminfP=\min(\lambda_{\mathrm{ev}},\lambda_{\mathrm{odd}})$
and $n_-(\Ainf|\PT)=n_{\mathrm{ev}}+n_{\mathrm{odd}}$. For a parity label $\mathrm{par}\in\{\mathrm{ev},\mathrm{odd}\}$
we write $\PT^{\mathrm{par}}$, $\lambda_{\mathrm{par}}$, $n_{\mathrm{par}}$, $T^*_{\mathrm{par}}$.

\begin{lemma}[parity splitting]\label{lem:parity}
$\PT=\PT^{\mathrm{ev}}\oplus\PT^{\mathrm{odd}}$, the sum being orthogonal both in
$L^2(\IT)$ and for $\mT$, which equals $\Ainf$ on $\PT$. Hence
\[
n_-(\Ainf|\PT)=n_{\mathrm{ev}}(T)+n_{\mathrm{odd}}(T),\qquad
\laminfP(T)=\min\big(\lambda_{\mathrm{ev}}(T),\lambda_{\mathrm{odd}}(T)\big).
\]
Here $n_-$ counts negative eigenvalues with multiplicity; for the closed,
lower-bounded form $\mT$, which has compact resolvent on each of these closed subspaces (proof of
Proposition~\ref{prop:nodes}(i)), this is the largest dimension of a subspace of the form domain on
which the form is negative definite.
\end{lemma}
\begin{proof}
$A(JF)=B(F)$ and $B(JF)=A(F)$, so for $F\in\PT$ both $F_{\mathrm e}=\frac12(F+JF)$ and
$F_{\mathrm o}=\frac12(F-JF)$ have vanishing Tate moments: $\PT$ is $J$-invariant and splits. The
parts are $L^2$-orthogonal, and
$\mT(F_{\mathrm e},F_{\mathrm o})=\frac1{2\pi}\int m\,\widehat{F_{\mathrm e}}\,\overline{\widehat{F_{\mathrm o}}}=0$,
because $m$ is even, $\widehat{F_{\mathrm e}}$ is even and $\widehat{F_{\mathrm o}}$ is odd. Negative
indices add over sums that are orthogonal for the norm and for the form, and bottoms take the minimum.
\end{proof}
The definition of $\Ainf$ writes the pole term as $2A(F)B(F)$, i.e.\ on real $F$, while $L^2(\IT)$ is
complex. Nothing depends on the choice: for $F=u+iv$ with $u,v$ real, $\mT(F)=\mT(u)+\mT(v)$
and $F\in\PT$ if and only if $u,v\in\PT$ (as in the proof of Lemma~\ref{lem:strict}), so the counts and
the bottoms are the same over $\R$ and over $\C$.

\begin{lemma}[zero extension]\label{lem:zeroext}
For $0<T<T'$, extension by zero maps $\PT^{\mathrm{par}}$ isometrically into
$\mathcal P^{\mathrm{par}}_{T'}$ with the same values of $\mT$. Hence $n_{\mathrm{par}}$ is
nondecreasing and $\lambda_{\mathrm{par}}$ nonincreasing in $T$, with no range condition at either window.
\end{lemma}
\begin{proof}
The norm, the Tate moments, the parity and $\cF$ are unchanged, and $\mT(F)=C_0\|F\|^2+\Lev(\cF)$
depends on $F$ only through $\cF$ on $\R$ (Theorem~\ref{thm:A}); $\Ainf$ does not depend on $T$
(Proposition~\ref{prop:arch}(i)).
\end{proof}

\begin{lemma}[strict monotonicity]\label{lem:strict}
For every $T>0$ the bottoms $\lambda_{\mathrm{ev}}(T)$, $\lambda_{\mathrm{odd}}(T)$ and
$\laminfP(T)$ are attained, and all three are strictly decreasing on $(0,\infty)$. Consequently, for each
sector, $\lambda_{\mathrm{par}}(T)>0$ for $T<T^*_{\mathrm{par}}$ and $\lambda_{\mathrm{par}}(T)<0$ for
$T>T^*_{\mathrm{par}}$; $\lambda_{\mathrm{par}}(T)<0$ if and only if $n_{\mathrm{par}}(T)\ge1$;
$T^*_{\infty,\mathrm P}=\min(T^*_{\mathrm{ev}},T^*_{\mathrm{odd}})$; and $\laminfP$ is positive on
$(0,T^*_{\infty,\mathrm P})$ and negative on $(T^*_{\infty,\mathrm P},\infty)$.
\end{lemma}
\begin{proof}
\emph{Attainment.} $\mT-C_0\|\cdot\|^2=\Lev(c_\cdot)\ge0$ is a closed form whose domain embeds
compactly in $L^2(\IT)$ (proof of Proposition~\ref{prop:nodes}(i)); $\PT^{\mathrm{par}}$ and $\PT$ are
closed, $A$ and $B$ being continuous on $L^2(\IT)$. A minimizing sequence on the unit sphere has an
$L^2$-convergent subsequence whose limit has norm $1$ and lies in the subspace, and lower semicontinuity
of the closed form gives attainment. Since the symbol is real and even, $\mT(u+iv)=\mT(u)+\mT(v)$
for real $u,v$, and $u+iv\in\PT^{\mathrm{par}}$ if and only if $u,v\in\PT^{\mathrm{par}}$ ($A,B$ are real on
real functions), so the minimizer may be taken real.

\emph{Strictness.} Monotonicity is Proposition~\ref{prop:arch}(i), sector by sector: extension by zero
maps $\mathcal P^{\mathrm{par}}_{T'}$ isometrically into $\PT^{\mathrm{par}}$ for $T'<T$ with the same
values of $\mT$ (Lemma~\ref{lem:zeroext}). Suppose $T'<T$ and
$\lambda_{\mathrm{par}}(T)=\lambda_{\mathrm{par}}(T')=:\lambda$, and let $F^*$ be a real minimizer on
$\mathcal P^{\mathrm{par}}_{T'}$. Extended by zero it is a minimizer on $\PT^{\mathrm{par}}$, so
$\mT(F^*,G)=\lambda\langle F^*,G\rangle$ for every $G\in\PT^{\mathrm{par}}$ in the form domain. Take
real $G_0\in C_c^\infty(T',T)$ with $\int G_0e^{y/2}dy=\int G_0e^{-y/2}dy=0$ and $G=G_0+\epsilon JG_0$,
$\epsilon=\pm1$ the parity sign. Then $G\in\PT^{\mathrm{par}}$, $\langle F^*,G\rangle=0$ by disjoint
supports, and $\mT(F^*,G)=2\mT(F^*,G_0)$, since $\mT$ is $J$-invariant and
$JF^*=\epsilon F^*$. Polarizing the Dirichlet form of Proposition~\ref{prop:nodes}(i),
$\Lev(\cF)=\frac14\iint_{\R^2}K(|x-y|)|\widetilde F(x)-\widetilde F(y)|^2$, for real $F,G$ with
supports at positive distance gives
$\mT(F,G)=-\frac12\iint K(|x-y|)F(x)G(y)\,dx\,dy$, absolutely convergent. Hence
$\Phi(y):=\int_{I_{T'}}K(|x-y|)F^*(x)\,dx$ satisfies $\int\Phi G_0=0$ for every such $G_0$, so
$\Phi=\alpha e^{y/2}+\beta e^{-y/2}$ on $(T',T)$. For $y>T'\ge x$, Lemma~\ref{lem:K} and Tonelli give
$\Phi(y)=2\sum_{k\ge0}c_ke^{-(2k+\frac12)y}$ with $c_k=\int F^*(x)e^{(2k+\frac12)x}dx$,
$|c_k|\le\|F^*\|_1e^{(2k+\frac12)T'}$; so $e^{y/2}\Phi(y)=2\sum_kc_kw^k$, $w=e^{-2y}$, is a convergent
power series for $0<w<e^{-2T'}$, and the identity extends to all $y>T'$. As $y\to\infty$, $\Phi\to0$,
so $\alpha=0$; then $2\sum_kc_kw^k\equiv\beta$, so $c_k=0$ for $k\ge1$ (the value $c_0=\beta/2$ is not
needed). With $u=e^{2x}$, $c_k=\int f(u)u^k\,du$ for
$f(u)=\frac12F^*(\frac12\log u)u^{-3/4}\in L^1[e^{-2T'},e^{2T'}]$, so all moments of $uf(u)$ vanish and
$uf=0$ by Weierstrass, whence $F^*=0$, contradicting $\|F^*\|=1$. For $\laminfP$ use Lemma~\ref{lem:parity}.

\emph{Consequences.} If $T>T^*_{\mathrm{par}}$, pick $T'\in(T^*_{\mathrm{par}},T)$: then
$\lambda_{\mathrm{par}}(T')\le0$ and $\lambda_{\mathrm{par}}(T)<\lambda_{\mathrm{par}}(T')$. If
$T<T^*_{\mathrm{par}}$, some $T''>T$ has $\lambda_{\mathrm{par}}(T'')>0$, and monotonicity gives
$\lambda_{\mathrm{par}}(T)>0$. A negative vector forces $\lambda_{\mathrm{par}}<0$, and conversely an
attained negative bottom spans a negative line. The last two statements follow from
$\laminfP=\min(\lambda_{\mathrm{ev}},\lambda_{\mathrm{odd}})$, the sets $\{\lambda_{\mathrm{par}}>0\}$ being
intervals.
\end{proof}
The same argument, with the pole term of $\Ainf$ in place of the two multipliers, shows that $\laminf$
and the bottom of $\mT$ on $L^2(\IT)$ are strictly decreasing as well (attainment as above, the pole
term being a bounded form of rank two); this is not used below.

\section{The two Tate directions}\label{sec:complete}

\begin{lemma}[non-attainment for the complete form]\label{lem:Onat}
Let $T\ge\frac12\log2$ and $w\in\Wcl$. Then $\inf_\R f_w<\lamW$.
\end{lemma}
\begin{proof}
Weak duality gives $c:=\inf f_w\le\lamW$ (the paragraph after Definition~\ref{def:W}). Suppose $c=\lamW$. The
bottom of $W_T$ on the unit sphere of $L^2(\IT)$ is attained at some $F_\star$ (\cite[Thm.~3]{Bombieri2000},
or the compactness argument in the proof of Theorem~\ref{thm:O}, which uses only that the symbol tends to
$+\infty$). Then $\varrho:=\sT+\nuh_w+2\gh0-\lamW\ge0$ is continuous, and
$\frac1{2\pi}\int\varrho|\Fh_\star|^2=W_T(F_\star)-\lamW=0$ by Proposition~\ref{prop:P}, since
$\frac1{2\pi}\int\nuh_w|\Fh_\star|^2=\int c_{F_\star}d\nu_w=0$. As $\Fh_\star$ is entire and nonzero,
$\varrho\equiv0$ on $\R$. Taking inverse Fourier transforms on $0<|u|<2T$, where $\nu_w=0$ and $g_0=\cosh\frac u2$,
and where the constant contributes nothing: $2\cosh\frac u2=\frac12K(|u|)+\sum_{n\in\NT}\frac{v_n}2(\delta_{\log n}+\delta_{-\log n})$
there, impossible since $K(u)\sim1/u$ as $u\downarrow0$ while $\cosh$ is bounded.
\end{proof}

\begin{theorem}[flat-top certificates and the complete bottom]\label{thm:Al}
\emph{(a)} If $f_w\ge0$ on $\R$ for some $w\in\Wcl$, regular or not, then $W_T\ge0$ on the form domain;
hence $\lamW\ge0$ and $\lamA\ge0$.
\emph{(b)} Let $T\ge\frac12\log2$. Then there is $w\in\Wcl$ with $f_w\ge0$ on $\R$ if and only if
$\lamW>0$, and $w$ can then be taken in $\Wcl^{\mathrm{reg}}$ with $\inf f_w>0$. Likewise, if $\lamP>0$
there are $w\in\Wcl^{\mathrm{reg}}$ and $\beta\in\R^2$ with $\inf_\tau[f_w+(\beta_0-2)\gh0+\beta_1\gh1]>0$.
\end{theorem}
\begin{proof}
(a) $(\nu_w,2,0)$ is dual feasible and $W_T\ge\inf f_w$ on $L^2(\IT)$ (the paragraph after
Definition~\ref{def:W}), and $W_T=\mathfrak a_T$ on $\PT$. (b) ``$\Leftarrow$'' and the regular choice:
Theorem~\ref{thm:Zexh}(i). ``$\Rightarrow$'': $\lamW\ge\inf f_w\ge0$ by weak duality, and $\lamW>\inf f_w$ by
Lemma~\ref{lem:Onat}. The last sentence is Theorem~\ref{thm:Zexh}(ii).
\end{proof}

\begin{definition}[the alignment inequality]\label{def:Al}
For $w\in\Wcl$ let $\mathrm{Al}_T(w)$ be the statement: for every $\tau\in\R$,
\[
-2\Rel\sum_{n\ge2}\frac{\Lambda(n)w(\log n)}{n^{\frac12+i\tau}}\ \ge\
-\Big[C_0+\int_0^\infty\big(1-w(u)\cos\tau u\big)K(u)\,du\Big]-2\int_\R w(u)\cosh\tfrac u2\cos(\tau u)\,du .
\]
The left side is the smoothed prime sum on the critical line, $-\sum_nv_nw(\log n)\cos(\tau\log n)$; by
\eqref{eq:fw}, $\mathrm{Al}_T(w)$ is exactly the statement $f_w\ge0$ on $\R$ of Theorem~\ref{thm:Al}.
\end{definition}

\begin{lemma}[the alignment inequality lives on a band]\label{lem:Alband}
For $w\in\Wcl$ let $\tau_{\mathrm{env}}(w)$ be the unique positive root of
$m_{\mathrm{lo}}(\tau)-\Aconst-\|\nu_w\|_{\TV}-2k_0(\tau)=0$, with $m_{\mathrm{lo}}$ as in Lemma~\ref{lem:mbounds}(ii)
and $k_0$ as in Lemma~\ref{lem:gbounds}. Then $f_w>0$ on $(\tau_{\mathrm{env}}(w),\infty)$, so $\mathrm{Al}_T(w)$ is
equivalent to its restriction to $|\tau|\le\tau_{\mathrm{env}}(w)$.
\end{lemma}
\begin{proof}
$m_{\mathrm{lo}}$ increases from $-\infty$ to $+\infty$ and $k_0$ decreases (Lemmas~\ref{lem:mbounds}, \ref{lem:gbounds}),
so the left side is strictly increasing and the root is unique. Proposition~\ref{prop:env}(i) for the feasible
triple $(\nu_w,2,0)$ gives $f_w\ge m_{\mathrm{lo}}-\Aconst-\|\nu_w\|_{\TV}-2k_0$, and $f_w$ is even.
\end{proof}

\begin{remark}[the per-window quantity]\label{rem:excess}
By Theorem~\ref{thm:Zexh}(i), $\sup_{w\in\Wcl^{\mathrm{reg}}}\inf_\tau f_w=\lamW$ for $T\ge\frac12\log2$: the largest
uniform margin with which the alignment inequality can hold is $\lamW$ itself, a number determined by
the finite contact weights and $m$. The trivial envelope $-\sum_nv_n|w(\log n)|$ of the prime
sum certifies $\mathrm{Al}_T(w)$ wherever the right side of Definition~\ref{def:Al} is at most it, and by
Proposition~\ref{prop:env}(ii) the constant $\Aconst$ of that envelope cannot be lowered; on the complement the
actual phases $\tau\log p$ modulo $2\pi$ enter. For $w\in\Wcl^{\mathrm{reg}}$, by
Theorem~\ref{thm:Z}(ii), $\mathrm{Al}_T(w)$ at $\tau$ is the inequality $\sum_\rho\wh(\tau-\gamma_\rho)\ge0$.
\end{remark}

\begin{remark}[the prime side of a primitive test is the Chebyshev remainder]\label{rem:Cheb}
Let $\psi_\Lambda(x)=\sum_{n\le x}\Lambda(n)$ and $r_\Lambda(u):=\psi_\Lambda(e^u)-e^u$. For real $F\in\PT\cap\mathcal D_T$,
\begin{equation}\label{eq:Cheb}
\mathfrak a_T(F)=C_0\|F\|^2+\Lev(\cF)+2\int_0^{2T}\cF(u)e^{-u/2}\,du-2\int_{(0,2T)}e^{-u/2}\cF(u)\,dr_\Lambda(u).
\end{equation}
Indeed $\sum_{n\in\NT}v_n\cF(\log n)=2\int_{(0,2T)}e^{-u/2}\cF(u)\,d\psi_\Lambda(e^u)$, since $u\mapsto\psi_\Lambda(e^u)$
jumps by $\Lambda(n)$ at $u=\log n$ and $\cF(2T)=0$; split $d\psi_\Lambda(e^u)=e^udu+dr_\Lambda(u)$. The identity
$\int_{-2T}^{2T}\cF(u)e^{u/2}du=A(F)B(F)=0$ on $\PT$ (substitute $s=t-u$) and the evenness of $\cF$ give
$\int_0^{2T}\cF e^{u/2}du=-\int_0^{2T}\cF e^{-u/2}du$; insert into \eqref{eq:Aform}. So on primitive tests the whole
arithmetic of the comb is the last term of \eqref{eq:Cheb}, and the unconditional control of $r_\Lambda$ is the
theory of zero-free regions. If $\lamA\ge0$, then for every $\eta>\lamA$ there are tests $F\in\PT$ with
$0\le\mathfrak a_T(F)\le\eta\|F\|^2$, and $\lamA$ is at most $4\cdot10^9Te^{-15.449T}$ for $T\ge1$
(Corollary~\ref{cor:margin1}); for such $F$, \eqref{eq:Cheb} says that $2\int e^{-u/2}\cF\,dr_\Lambda$ equals
$C_0\|F\|^2+\Lev(\cF)+2\int_0^{2T}\cF e^{-u/2}du$ up to $\eta\|F\|^2$.
\end{remark}


\emph{The per-window converse.} At a single window, the converse of Theorem~\ref{thm:Al}(a) is the
implication $\mathrm{WP}_T\Rightarrow\lamW\ge0$. It reduces to the signs of two numbers. Fix $T$. We work
with real functions. This suffices for bottoms and negative indices, as in the remark after
Lemma~\ref{lem:parity}: for $F=u+iv$ with $u,v$ real, $W_T(F)=W_T(u)+W_T(v)$ and
$\mathfrak a_T(F)=\mathfrak a_T(u)+\mathfrak a_T(v)$. Let $\mathcal D_T$ be the form domain of \eqref{eq:form}
and $\mathcal D^{\mathrm{par}}_T$ its elements of parity $\mathrm{par}\in\{\mathrm{ev},\mathrm{odd}\}$. With
$\chi_{\mathrm{par}}$ as in \S\ref{subsec:conv}, put
$\ell_{\mathrm{par}}(F)=\langle F,\chi_{\mathrm{par}}\rangle$, $\varepsilon_{\mathrm{ev}}=+1$ and
$\varepsilon_{\mathrm{odd}}=-1$. Then $A=\ell_{\mathrm{ev}}+\ell_{\mathrm{odd}}$ and
$B=\ell_{\mathrm{ev}}-\ell_{\mathrm{odd}}$ on real $F$ (the parity paragraph after
Proposition~\ref{prop:nodes}). Let $\mathcal Q^{\mathrm{par}}_T:=\PT^{\mathrm{par}}\cap\mathcal D_T
=\ker(\ell_{\mathrm{par}}|\mathcal D^{\mathrm{par}}_T)$, and write $\lambda^{\mathrm{par}}_{\mathrm A}(T)$ for the
bottom of $\mathfrak a_T$ on its unit sphere. This is the bottom with the primes, not the prime-free
$\lambda_{\mathrm{par}}$ of \S\ref{subsec:arch}. Write $\lambda^{\mathrm W}_{T,\mathrm{par}}$ for the bottom of
$W_T$ on the unit sphere of $\mathcal D^{\mathrm{par}}_T$. As in Lemma~\ref{lem:parity}, $n_-$ is the largest
dimension of a subspace of the form domain on which the form is negative definite. The letters $h$,
$\mathfrak g$, $\mathfrak G$, $S^{\mathrm{Sch}}$, $\mathcal N$ are local to this section.

\begin{theorem}[the parity Schur complement of the complete form]\label{thm:parityschur}
Fix $T>0$.
\emph{(o) Parity splitting.} On real $F=F_{\mathrm e}+F_{\mathrm o}\in\mathcal D_T$,
\[
W_T(F)=\big[\mathfrak a_T(F_{\mathrm e})+2\ell_{\mathrm{ev}}(F_{\mathrm e})^2\big]
+\big[\mathfrak a_T(F_{\mathrm o})-2\ell_{\mathrm{odd}}(F_{\mathrm o})^2\big],
\]
the sum being orthogonal for $L^2$, for $\mathfrak a_T$ and for $W_T$. Hence
$\lamA=\min_{\mathrm{par}}\lambda^{\mathrm{par}}_{\mathrm A}(T)$ and $\lamW=\min_{\mathrm{par}}\lambda^{\mathrm W}_{T,\mathrm{par}}$.

\emph{(a) Non-degenerate sector.} Suppose $0$ is not in the spectrum of the self-adjoint operator of the
closed form $\mathfrak a_T|\mathcal Q^{\mathrm{par}}_T$. This holds in particular if
$\lambda^{\mathrm{par}}_{\mathrm A}(T)>0$. Then there is a unique $h_{\mathrm{par}}\in\mathcal D^{\mathrm{par}}_T$
with $\ell_{\mathrm{par}}(h_{\mathrm{par}})=1$ and $\mathfrak a_T(h_{\mathrm{par}},G)=0$ for all
$G\in\mathcal Q^{\mathrm{par}}_T$. Put $\mathfrak g_{\mathrm{par}}(T):=W_T(h_{\mathrm{par}})=\mathfrak a_T(h_{\mathrm{par}})+2\varepsilon_{\mathrm{par}}$.
Then:
\begin{itemize}
\item[(a1)] $\mathcal D^{\mathrm{par}}_T=\mathcal Q^{\mathrm{par}}_T\oplus\R h_{\mathrm{par}}$, orthogonally for
$\mathfrak a_T$ and for $W_T$, and $W_T(G+xh_{\mathrm{par}})=\mathfrak a_T(G)+\mathfrak g_{\mathrm{par}}x^2$ for
$G\in\mathcal Q^{\mathrm{par}}_T$, $x\in\R$;
\item[(a2)] $n_-(W_T|\mathcal D^{\mathrm{par}}_T)=n_-(\mathfrak a_T|\mathcal Q^{\mathrm{par}}_T)+[\mathfrak g_{\mathrm{par}}<0]$.
The radical of $W_T$ on $\mathcal D^{\mathrm{par}}_T$ is $\R h_{\mathrm{par}}$ if $\mathfrak g_{\mathrm{par}}=0$ and $0$
otherwise;
\item[(a3)] if $\lambda^{\mathrm{par}}_{\mathrm A}(T)>0$, then
$\mathfrak a_T(h_{\mathrm{par}})=\min\{\mathfrak a_T(F):F\in\mathcal D^{\mathrm{par}}_T,\ \ell_{\mathrm{par}}(F)=1\}$ and
$\mathfrak g_{\mathrm{par}}=\min\{W_T(F):F\in\mathcal D^{\mathrm{par}}_T,\ \ell_{\mathrm{par}}(F)=1\}$;
\item[(a4)] $\lambda^{\mathrm W}_{T,\mathrm{par}}\le\mathfrak g_{\mathrm{par}}/\|h_{\mathrm{par}}\|^2$ and
$\|h_{\mathrm{par}}\|\ge1/\|\chi_{\mathrm{par}}\|$. If $\lambda^{\mathrm{par}}_{\mathrm A}(T)>0$ \emph{and}
$\mathfrak g_{\mathrm{par}}>0$, then
$\lambda^{\mathrm W}_{T,\mathrm{par}}\ge\frac12\min\big(\lambda^{\mathrm{par}}_{\mathrm A}(T),\mathfrak g_{\mathrm{par}}/\|h_{\mathrm{par}}\|^2\big)$.
If in addition $2\lambda^{\mathrm W}_{T,\mathrm{par}}<\lambda^{\mathrm{par}}_{\mathrm A}(T)$, then
$\lambda^{\mathrm W}_{T,\mathrm{par}}\le\mathfrak g_{\mathrm{par}}/\|h_{\mathrm{par}}\|^2\le2\lambda^{\mathrm W}_{T,\mathrm{par}}$.
\end{itemize}

\emph{(b) Degenerate sector.} Suppose $\lambda^{\mathrm{par}}_{\mathrm A}(T)=0$, and let
$\mathcal N=\ker(\mathfrak a_T|\mathcal Q^{\mathrm{par}}_T)$. It is nonzero and finite-dimensional. For each
$F_0\in\mathcal N$ there is a unique $\alpha(F_0)\in\R$ with $\mathfrak a_T(F_0,G)=\alpha(F_0)\ell_{\mathrm{par}}(G)$
for all $G\in\mathcal D^{\mathrm{par}}_T$. If $\alpha\not\equiv0$ on $\mathcal N$, then $W_T$ is unbounded below on
$\mathcal D^{\mathrm{par}}_T$, so $\lamW<0$. If $\alpha\equiv0$ on $\mathcal N$ (the \emph{range condition}), then
$\mathcal N$ lies in the radical of $W_T$ on $\mathcal D^{\mathrm{par}}_T$ and (a) applies on the $L^2$-orthogonal
complement of $\mathcal N$ in $\mathcal D^{\mathrm{par}}_T$. With the $h_{\mathrm{par}},\mathfrak g_{\mathrm{par}}$ so obtained,
$W_T\ge0$ on $\mathcal D^{\mathrm{par}}_T$ if and only if $\mathfrak g_{\mathrm{par}}\ge0$.

\emph{(c) Assembled.} Put $\mathfrak G_T:=\mathrm{diag}(\mathfrak g_{\mathrm{ev}}(T),\mathfrak g_{\mathrm{odd}}(T))$. It is
defined whenever each sector is as in (a), or as in (b) with the range condition. If $0$ is not in the
spectrum of the operator of $\mathfrak a_T|(\PT\cap\mathcal D_T)$, then
$n_-(W_T)=n_-(\mathfrak a_T|\PT)+n_-(\mathfrak G_T)$, and the radical of $W_T$ has dimension $n_0(\mathfrak G_T)$. At every window,
without further hypotheses,
\[
\lamW>0\iff\big[\lamA>0\ \text{and}\ \mathfrak G_T\succ0\big],
\]
\[
\lamW\ge0\iff\big[\mathrm{WP}_T,\ \text{the range condition in every sector with }
\lambda^{\mathrm{par}}_{\mathrm A}(T)=0,\ \text{and}\ \mathfrak G_T\succeq0\big].
\]
In particular, if $\lamA>0$, then $\lamW\ge0$ if and only if $\mathfrak g_{\mathrm{ev}}(T)\ge0$ and
$\mathfrak g_{\mathrm{odd}}(T)\ge0$.
\end{theorem}
\begin{proof}
(o) The symbol $\sT$ is even. For real even $F_{\mathrm e}$ and real odd $F_{\mathrm o}$, $\widehat{F_{\mathrm e}}$ is even and
$\widehat{F_{\mathrm o}}$ is odd, so $\mathfrak a_T(F_{\mathrm e},F_{\mathrm o})=0$ as in the proof of
Lemma~\ref{lem:parity}. Also $\ell_{\mathrm{ev}}$ vanishes on odd functions and $\ell_{\mathrm{odd}}$ on even ones, so
$A(F)B(F)=(\ell_{\mathrm{ev}}(F_{\mathrm e})+\ell_{\mathrm{odd}}(F_{\mathrm o}))(\ell_{\mathrm{ev}}(F_{\mathrm e})-\ell_{\mathrm{odd}}(F_{\mathrm o}))
=\ell_{\mathrm{ev}}(F_{\mathrm e})^2-\ell_{\mathrm{odd}}(F_{\mathrm o})^2$, with no cross term. Polarizing gives the
orthogonality. $F\in\PT$ if and only if $\ell_{\mathrm{ev}}(F_{\mathrm e})=\ell_{\mathrm{odd}}(F_{\mathrm o})=0$, and bottoms over an
orthogonal sum take the minimum.

(a) Write $\mathcal Q=\mathcal Q^{\mathrm{par}}_T$, $\ell=\ell_{\mathrm{par}}$ and $\varepsilon=\varepsilon_{\mathrm{par}}$. On
$\mathcal D^{\mathrm{par}}_T$, $W_T(F,G)=\mathfrak a_T(F,G)+2\varepsilon\,\ell(F)\ell(G)$ by (o).

\emph{Solvability.} First, $\mathfrak a_T=\mT+\PiT$ with $\|\PiT\|\le\Aconst$, since $|\cF(\log n)|\le\|F\|^2$ for each $n$ (this is Proposition~\ref{prop:coer-prime}(i) below). The form
$\mT\ge C_0\|\cdot\|^2$ is closed, with domain $\mathcal D_T$ (Theorem~\ref{thm:A}), and this domain embeds
compactly in $L^2(\IT)$ (proof of Proposition~\ref{prop:nodes}(i)). So $\mathfrak a_T$ is closed and bounded
below by $-M$, where $M:=\Aconst-C_0$. Second, $\mathcal Q$ is closed in the form norm, because $\ell$ and the
parity projection are $L^2$-continuous. Third, $\mathcal Q$ is dense in
$H:=\ker\ell\cap L^2_{\mathrm{par}}(\IT)$. Indeed, let $x\in H$, let $d_n\in\mathcal D^{\mathrm{par}}_T$ with $d_n\to x$ in
$L^2$ (smooth functions of parity $\mathrm{par}$ are dense), and let $e\in\mathcal D^{\mathrm{par}}_T$ with $\ell(e)=1$.
Then $d_n-\ell(d_n)e\in\mathcal Q$ tends to $x$. Hence $\mathfrak a_T|\mathcal Q$ is a closed, semibounded, densely
defined form in $H$ with compact resolvent. Its spectrum is discrete, and its negative spectral subspace
$\mathcal Q_-$ is finite-dimensional and $\mathfrak a_T$-orthogonal to $\mathcal Q_+:=\mathcal Q\cap\mathcal Q_-^\perp$. If
$0$ is not in the spectrum, then $\mathfrak a_T\ge\delta\|\cdot\|^2$ on $\mathcal Q_+$ for some $\delta>0$. Hence
$\mathfrak a_T(x)\ge\frac{\delta}{\delta+M+1}\big(\mathfrak a_T(x)+(M+1)\|x\|^2\big)$ there, so $\mathfrak a_T$ is coercive
on $\mathcal Q_+$ in the form norm and the Lax--Milgram lemma applies on $\mathcal Q_+$. On $\mathcal Q_-$ the form is a
negative definite finite matrix. So for every form-bounded linear functional $f$ on $\mathcal Q$ there is a unique
$Y\in\mathcal Q$ with $\mathfrak a_T(Y,\cdot)=f$ on $\mathcal Q$. Uniqueness holds because the operator has trivial
kernel. Now fix $H_0\in\mathcal D^{\mathrm{par}}_T$ with $\ell(H_0)=1$, for instance a smooth bump of parity
$\mathrm{par}$. The functional $f=\mathfrak a_T(H_0,\cdot)$ is form-bounded, by Cauchy--Schwarz for the nonnegative
form $\mathfrak a_T+M\|\cdot\|^2$. Put $h:=H_0-Y$. Then $\ell(h)=1$ and $\mathfrak a_T(h,\cdot)=0$ on $\mathcal Q$. Two such
$h$ differ by an element of $\mathcal Q$ that is $\mathfrak a_T$-orthogonal to $\mathcal Q$, which is $0$.

(a1) Any $F\in\mathcal D^{\mathrm{par}}_T$ decomposes as $F=(F-\ell(F)h)+\ell(F)h$ with $F-\ell(F)h\in\mathcal Q$. For
$G\in\mathcal Q$, $W_T(h,G)=\mathfrak a_T(h,G)+2\varepsilon\ell(h)\ell(G)=0$. Expand, using $\ell(G)=0$ and $\ell(h)=1$.

(a2) By (a1), $\mathcal D^{\mathrm{par}}_T=(\mathcal Q_-\oplus[\R h\text{ if }\mathfrak g<0])\oplus(\mathcal Q_+\oplus[\R h\text{ if }\mathfrak g\ge0])$,
orthogonally for $W_T$. The first summand is negative definite and has the stated dimension. The second
summand carries $W_T\ge0$. A negative definite subspace meets the second summand only in $0$, so its
projection onto the first summand along the second is injective. For the radical: if
$W_T(G+xh,\cdot)=0$, then $\mathfrak a_T(G,\cdot)=0$ on $\mathcal Q$, so $G=0$, and then $\mathfrak g x=0$.

(a3) If the bottom is positive, then $\mathfrak a_T(G+h)=\mathfrak a_T(G)+\mathfrak a_T(h)\ge\mathfrak a_T(h)$ and
$W_T(G+h)=\mathfrak a_T(G)+\mathfrak g\ge\mathfrak g$.

(a4) The upper bound is the Rayleigh quotient at $h$, and $1=\ell(h)\le\|h\|\|\chi_{\mathrm{par}}\|$. Put
$\lambda:=\lambda^{\mathrm{par}}_{\mathrm A}(T)$ and assume $\lambda>0$ and $\mathfrak g>0$. With
$\mu:=\min(\lambda,\mathfrak g/\|h\|^2)$,
\[
W_T(G+xh)=\mathfrak a_T(G)+\mathfrak gx^2\ge\lambda\|G\|^2+\frac{\mathfrak g}{\|h\|^2}\|xh\|^2
\ge\mu\big(\|G\|^2+\|xh\|^2\big)\ge\tfrac12\mu\|G+xh\|^2 .
\]
If also $2\lambda^{\mathrm W}_{T,\mathrm{par}}<\lambda$, then $\frac12\mu\le\lambda^{\mathrm W}_{T,\mathrm{par}}<\frac12\lambda$ forces
$\mu=\mathfrak g/\|h\|^2$. The hypothesis $\mathfrak g>0$ is used in the second inequality and cannot be dropped.

(b) The bottom $0$ is attained (compact resolvent), so $\mathcal N\ne0$, and $\mathcal N$ is finite-dimensional. For
$F_0\in\mathcal N$ and $G\in\mathcal Q$, Cauchy--Schwarz for the nonnegative form $\mathfrak a_T|\mathcal Q$ gives
$\mathfrak a_T(F_0,G)=0$. So $\mathfrak a_T(F_0,\cdot)$ is a linear functional on $\mathcal D^{\mathrm{par}}_T$ vanishing on
$\ker\ell$, hence a multiple $\alpha(F_0)\ell$. Since $\ell(F_0)=0$, we get $W_T(F_0,H)=\alpha(F_0)\ell(H)$ and
$W_T(F_0)=0$. So $W_T(xF_0+H)=W_T(H)+2x\alpha(F_0)\ell(H)$, which is unbounded below in $x$ if
$\alpha(F_0)\ne0$ and $\ell(H)=1$. Now suppose $\alpha\equiv0$. Then $\mathcal N$ is $W_T$-orthogonal to
$\mathcal D^{\mathrm{par}}_T$. Put $\mathcal D':=\mathcal D^{\mathrm{par}}_T\cap\mathcal N^\perp$ and $\mathcal Q':=\mathcal Q\cap\mathcal N^\perp$.
Then $\mathcal Q'=\ker(\ell|\mathcal D')$, and $\ell\ne0$ on $\mathcal D'$: subtract from $H_0$ its $L^2$-projection onto
$\mathcal N\subset\mathcal Q$, which $\ell$ annihilates. On $\mathcal Q'$ the form is bounded below by the next eigenvalue,
which is $>0$. So (a) applies with $\mathcal D',\mathcal Q'$ in place of $\mathcal D^{\mathrm{par}}_T,\mathcal Q$. Moreover
$\mathcal D^{\mathrm{par}}_T=\mathcal N\oplus\mathcal D'$, $W_T$-orthogonally. Conclude by (a1).

(c) By (o), everything adds over the two sectors, and the inertia formula is (a2) in each sector. If
$\lamW>0$, then $\lamA\ge\lamW>0$ (recall $\lamW\le\lamA$), so both sectors are as in (a), and
$\mathfrak g_{\mathrm{par}}=W_T(h_{\mathrm{par}})\ge\lamW\|h_{\mathrm{par}}\|^2>0$. The converse is (a4) in each sector. If $\lamW\ge0$,
then $\lamA\ge\lamW\ge0$, which is $\mathrm{WP}_T$. In a sector with bottom $0$, (b) excludes $\alpha\not\equiv0$, and
in every sector $\mathfrak g_{\mathrm{par}}=W_T(h_{\mathrm{par}})\ge0$. Conversely, (a1) and $\mathfrak a_T\ge0$ on
$\mathcal Q^{\mathrm{par}}_T$ give $W_T\ge0$ on a sector with positive bottom when $\mathfrak g_{\mathrm{par}}\ge0$, and (b) does the
same on a sector with bottom $0$.
\end{proof}

The range condition says that the column $\rho(G):=\mathfrak a_T(H_0,G)$ of the bordered form, with
$\ell_{\mathrm{par}}(H_0)=1$, vanishes on $\ker(\mathfrak a_T|\mathcal Q^{\mathrm{par}}_T)$, since $\rho(F_0)=\alpha(F_0)$. This is the
range condition of a bordered Schur complement. Because $\mathfrak a_T|\mathcal Q^{\mathrm{par}}_T$ has a spectral gap above
$0$, it is equivalent to $\rho$ lying in the range of the square root of the operator. The theorem takes
the complement over $\PT$ and not over $L^2(\IT)$. So only $\mathfrak a_T|\PT$ has to be invertible, which is
the natural hypothesis for a converse whose premise is on the primitive side.

\begin{proposition}[$\mathfrak G_T$ is the polar capacity]\label{prop:polarcap}
Assume $\mathrm{WP}_T$, and the range condition of Theorem~\ref{thm:parityschur}(b) in every
sector with $\lambda^{\mathrm{par}}_{\mathrm A}(T)=0$. Let
$S^{\mathrm{Sch}}_T:=\mathrm{diag}\big(\mathfrak a_T(h_{\mathrm{ev}}),\mathfrak a_T(h_{\mathrm{odd}})\big)$. This is the compression of
$\mathfrak a_T$ to $\mathrm{span}(h_{\mathrm{ev}},h_{\mathrm{odd}})$, the $\mathfrak a_T$-orthogonal complement of the primitive part,
i.e.\ the Schur complement of $\mathfrak a_T$ over $\PT$. Let $D_{\mathrm{pol}}:=\mathrm{diag}(2,-2)$. Then
\[
\mathfrak G_T=S^{\mathrm{Sch}}_T+D_{\mathrm{pol}} .
\]
For $(\beta,\alpha)\in\R^2$ put
$E_T(\beta,\alpha):=\inf\{W_T(F):\ F\in\mathcal D_T\ \text{real},\ \Phi_F(0)=B(F)=\beta,\ \Phi_F(1)=A(F)=\alpha\}$. Then
\[
E_T(\beta,\alpha)=(\alpha,\beta)\,\mathfrak G^{AB}_T\,(\alpha,\beta)^{\mathsf T},\qquad
\mathfrak G^{AB}_T=\frac14\begin{pmatrix}\mathfrak g_{\mathrm{ev}}+\mathfrak g_{\mathrm{odd}}&\mathfrak g_{\mathrm{ev}}-\mathfrak g_{\mathrm{odd}}\\
\mathfrak g_{\mathrm{ev}}-\mathfrak g_{\mathrm{odd}}&\mathfrak g_{\mathrm{ev}}+\mathfrak g_{\mathrm{odd}}\end{pmatrix},
\]
and the infimum is attained. In these coordinates $D_{\mathrm{pol}}$ becomes $\big(\begin{smallmatrix}0&1\\1&0\end{smallmatrix}\big)$,
the form $2\alpha\beta=2\Phi_F(0)\Phi_F(1)$, with the one negative square of Lemma~\ref{lem:polar}. Hence
$\lamW\ge0$ if and only if $E_T\ge0$ on $\R^2$. Under $\mathrm{WP}_T$, the complete form is nonnegative
exactly when the least complete Weil energy needed to prescribe the values of $\Phi_F$ at the pole
$s\in\{0,1\}$ is nonnegative.
\end{proposition}
\begin{proof}
$\mathfrak G_T=S^{\mathrm{Sch}}_T+D_{\mathrm{pol}}$ is the definition $\mathfrak g_{\mathrm{par}}=\mathfrak a_T(h_{\mathrm{par}})+2\varepsilon_{\mathrm{par}}$. In the
degenerate case with the range condition, $\mathfrak a_T(h_{\mathrm{par}})$ does not depend on the choice of $h_{\mathrm{par}}$
modulo $\mathcal N$. For real $F$, $\Phi_F(1)=A(F)$ and $\Phi_F(0)=B(F)$ (proof of Lemma~\ref{lem:roughEF}). The
constraints are $\ell_{\mathrm{ev}}(F_{\mathrm e})=c:=\frac{\alpha+\beta}2$ and $\ell_{\mathrm{odd}}(F_{\mathrm o})=s:=\frac{\alpha-\beta}2$,
and they split over the sectors, as does $W_T$ by Theorem~\ref{thm:parityschur}(o). In each sector,
Theorem~\ref{thm:parityschur}(a1) (or (b)) gives $W_T(G+xh_{\mathrm{par}})=\mathfrak a_T(G)+\mathfrak g_{\mathrm{par}}x^2$ with $x$ the
prescribed value of $\ell_{\mathrm{par}}$ and $\mathfrak a_T(G)\ge0$ by $\mathrm{WP}_T$. So the sector minimum is
$\mathfrak g_{\mathrm{par}}x^2$, attained at $G=0$. Hence $E_T=\mathfrak g_{\mathrm{ev}}c^2+\mathfrak g_{\mathrm{odd}}s^2$, which is the displayed
matrix. The form $2c^2-2s^2=2\alpha\beta$ gives $D_{\mathrm{pol}}$ in the new coordinates. The last statement is
Theorem~\ref{thm:parityschur}(c), since $E_T\ge0$ if and only if $\mathfrak G_T\succeq0$.
\end{proof}

\subsection{The optimal multipliers}\label{subsec:mult}
Remark~\ref{rem:Kprim} leaves open whether $\lamP=\lamW$, that is, whether the multipliers $\beta$ of
near-optimal certificates of the primitive problem tend to the polar value $(2,0)$ of
Proposition~\ref{prop:P}. The question has a variational answer. For $F\in L^2(\IT)$ put
$A'(F)=\int tF(t)e^{t/2}dt$, $B'(F)=\int tF(t)e^{-t/2}dt$ and
\[
\mathfrak m(F):=\Big(\frac1{2\pi}\int\gh0|\Fh|^2,\ \frac1{2\pi}\int\gh1|\Fh|^2\Big)
=\big(\Rel(A\bar B),\ \Rel(A'\bar B-A\bar B')\big),
\]
the two moments of Lemma~\ref{lem:W} (the second expression is the computation in the proof of
Theorem~\ref{thm:A}), and for $\beta\in\R^2$
\[
V_T(\beta):=\inf\Big\{\frac1{2\pi}\int(\sT+\beta\cdot\widehat g)|\Fh|^2\,d\tau:\ F\in\mathcal D_T,\ \|F\|=1\Big\}
=\inf_{F\in\mathcal D_T,\ \|F\|=1}\big[\mathfrak a_T(F)+\beta\cdot\mathfrak m(F)\big].
\]
By Proposition~\ref{prop:P}, $V_T(2,0)=\lamW$. For real $F=F_{\mathrm e}+F_{\mathrm o}$ the first moment is
$\Rel(A\bar B)=\ell_{\mathrm{ev}}(F_{\mathrm e})^2-\ell_{\mathrm{odd}}(F_{\mathrm o})^2$ (proof of
Theorem~\ref{thm:parityschur}(o)), and for $F=u+iv$ with $u,v$ real one has $\mathfrak m(F)=\mathfrak m(u)+\mathfrak m(v)$,
$W_T(F)=W_T(u)+W_T(v)$ and $\|F\|^2=\|u\|^2+\|v\|^2$, since $\gh0,\gh1,\sT$ are even and
$\Iml(\widehat u\,\overline{\widehat v})$ is odd.

\begin{theorem}[the optimal multipliers]\label{thm:mult}
Fix $T>0$.
\emph{(i)} $V_T$ is finite, concave and continuous on $\R^2$, $V_T(\beta)\to-\infty$ as $|\beta|\to\infty$, and
\[
\lamP=\max_{\beta\in\R^2}V_T(\beta).
\]
Let $\mathcal M_T\subset\R^2$ be the set of maximizers; it is nonempty, compact and convex.
\emph{(ii)} If $(\nu_k,\beta^{(k)})$ are certificates of values $c_k\to\lamP$, then
$\operatorname{dist}(\beta^{(k)},\mathcal M_T)\to0$. Hence every such sequence satisfies $\beta^{(k)}\to(2,0)$ if and
only if $\mathcal M_T=\{(2,0)\}$, and $\lamP=\lamW$ if and only if $(2,0)\in\mathcal M_T$.
\emph{(iii)} Let $\mathcal Z_T$ be the set of real unit vectors $F\in\mathcal D_T$ with $W_T(F)=\lamW$; it is nonempty
and compact in $L^2(\IT)$. Then $(2,0)\in\mathcal M_T$ if and only if
$0\in\operatorname{conv}\{\mathfrak m(F):F\in\mathcal Z_T\}$.
\emph{(iv)} Suppose $\lamW<\lamA$. Then every $F\in\mathcal Z_T$ of definite parity has
$\Rel(A\bar B)=\ell_{\mathrm{ev}}(F)^2>0$ if $F$ is even and $\Rel(A\bar B)=-\ell_{\mathrm{odd}}(F)^2<0$ if $F$ is odd.
Consequently, if moreover $\lambda^{\mathrm W}_{T,\mathrm{ev}}\ne\lambda^{\mathrm W}_{T,\mathrm{odd}}$, then $\lamW<\lamP$, and
the multipliers of every maximizing sequence of certificates stay at a positive distance from $(2,0)$.
\emph{(v)} If $\lamW=\lamA$, then $\lamP=\lamW$ and $(2,0)\in\mathcal M_T$.
\end{theorem}
\begin{proof}
(i) For a unit vector $F$, $|\Rel(A\bar B)|\le\|e^{t/2}\|\,\|e^{-t/2}\|=2\sinh T$ and
$|\Rel(A'\bar B-A\bar B')|\le2\|te^{t/2}\|\,\|e^{-t/2}\|\le4T\sinh T$, so $|\mathfrak m(F)|\le C_T:=2\sinh T\,(1+2T)$.
Hence $V_T(\beta)\ge C_0-\Aconst-C_T|\beta|>-\infty$, and $V_T(\beta)<\infty$ since $\mathcal D_T\ne0$. As an infimum
of affine functions of $\beta$, $V_T$ is concave; a finite concave function on $\R^2$ is continuous. For the
behaviour at infinity, Lemma~\ref{lem:W} gives $r>0$ and $V_0<\infty$ such that every point of the ball of radius
$r$ is the moment vector $\mathfrak m(F)$ of some unit $F\in\mathcal D_T$ with $\mathfrak a_T(F)\le V_0$ (the ball is covered by
convex combinations, at the level of the measures $|\Fh|^2d\tau/2\pi$, of four measures of finite value, and such a
measure is $|\Fh|^2d\tau/2\pi$ for a unit $F$ by the proof of Proposition~\ref{prop:weak}). Taking
$\mathfrak m(F)=-r\beta/|\beta|$ gives $V_T(\beta)\le V_0-r|\beta|$.
Next, $\lamP\ge V_T(\beta)$ for every $\beta$: if $\mu$ is feasible for (P) then $\int(\sT+\beta\cdot\widehat g)d\mu=\int\sT d\mu$,
and the feasible set of (P) is contained in $\{|\Fh|^2d\tau/2\pi:\|F\|=1\}$ (Proposition~\ref{prop:weak}). Conversely,
by Theorem~\ref{thm:D1} there are, for every $\varepsilon>0$, a certificate $(\nu,\beta)$ of value $c>\lamP-\varepsilon$,
and $c\le V_T(\beta)$, because for a unit $F$ one has $\frac1{2\pi}\int\nuh|\Fh|^2=\int\cF\,d\nu=0$, so that
$\frac1{2\pi}\int(\sT+\beta\cdot\widehat g)|\Fh|^2=\frac1{2\pi}\int f|\Fh|^2\ge c$. Hence $\sup V_T=\lamP$, the
supremum is attained because $V_T$ is continuous and tends to $-\infty$ at infinity, and $\mathcal M_T$, a superlevel
set of a concave function, is closed, bounded and convex.
(ii) As just shown, $c_k\le V_T(\beta^{(k)})\le\lamP$, so $V_T(\beta^{(k)})\to\lamP$. If $\operatorname{dist}(\beta^{(k)},\mathcal M_T)\ge\delta>0$
along a subsequence, that subsequence is bounded (coercivity) and has a limit point $\beta^*$ with
$V_T(\beta^*)=\lamP$ by continuity, i.e.\ $\beta^*\in\mathcal M_T$, a contradiction. The two equivalences follow, the
second one because $V_T(2,0)=\lamW\le\lamP$.
(iii) By Lemma~\ref{lem:compact} the bottom of $W_T$ on the unit sphere is attained, so $\mathcal Z_T\ne\emptyset$
(a complex minimizer has a real or an imaginary part that is a minimizer, by the additivity recalled above).
Compactness: unit vectors with $W_T(F)\le\lamW+\eta$ have $\mathfrak a_T(F)\le\lamW+\eta+4\sinh T$, so by
Lemma~\ref{lem:compact} they form a relatively compact set; $W_T$ is lower semicontinuous on $L^2(\IT)$ (a closed
form bounded below, plus a continuous term), so a limit $F$ of minimizers satisfies $\|F\|=1$ and
$W_T(F)\le\lamW$, hence is a minimizer. The map $\mathfrak m$ is continuous on $L^2(\IT)$, so
$\mathfrak m(\mathcal Z_T)$ and its convex hull are compact.
\emph{Sufficiency.} If $0=\sum_i\theta_i\mathfrak m(F_i)$ with $F_i\in\mathcal Z_T$, $\theta_i\ge0$, $\sum\theta_i=1$, then for
every $\beta'$, $V_T(\beta')\le\mathfrak a_T(F_i)+\beta'\cdot\mathfrak m(F_i)=\lamW+(\beta'-(2,0))\cdot\mathfrak m(F_i)$ for each $i$;
averaging, $V_T(\beta')\le\lamW=V_T(2,0)$, so $(2,0)\in\mathcal M_T$.
\emph{Necessity.} Suppose $0\notin\operatorname{conv}\mathfrak m(\mathcal Z_T)$. By strict separation there are $d\in\R^2$ and
$\delta>0$ with $d\cdot\mathfrak m(F)\ge\delta$ for all $F\in\mathcal Z_T$, and then also for all complex unit minimizers,
since for $F=u+iv$ one has $\mathfrak m(F)=\|u\|^2\mathfrak m(u/\|u\|)+\|v\|^2\mathfrak m(v/\|v\|)$ with $u/\|u\|,v/\|v\|\in\mathcal Z_T$
whenever nonzero. We claim that there is $\eta>0$ with $d\cdot\mathfrak m(F)\ge\delta/2$ for every unit $F\in\mathcal D_T$
with $W_T(F)\le\lamW+\eta$: otherwise there are unit $F_k$ with $W_T(F_k)\to\lamW$ and $d\cdot\mathfrak m(F_k)<\delta/2$;
by the compactness above a subsequence converges in $L^2$ to a unit minimizer $F^*$, and by continuity
$d\cdot\mathfrak m(F^*)\le\delta/2<\delta$, a contradiction. Now for $t>0$ and $\beta=(2,0)$,
$V_T(\beta+td)=\inf_F[W_T(F)+t\,d\cdot\mathfrak m(F)]$ over unit $F\in\mathcal D_T$: if $W_T(F)\le\lamW+\eta$ the bracket is
$\ge\lamW+t\delta/2$, and otherwise it is $\ge\lamW+\eta-tC_T|d|$. For $0<t<\eta/(2C_T|d|)$ this gives
$V_T(\beta+td)\ge\lamW+\min(t\delta/2,\eta/2)>\lamW=V_T(\beta)$, so $\beta\notin\mathcal M_T$.
(iv) Let $F\in\mathcal Z_T$ be even, $x=\ell_{\mathrm{ev}}(F)$. Then $A=B=x$ (the parity paragraph before
Theorem~\ref{thm:parityschur}) and $W_T(F)=\mathfrak a_T(F)+2x^2$. If $x=0$ then $F\in\PT\cap\mathcal D_T$ and
$W_T(F)=\mathfrak a_T(F)\ge\lamA>\lamW$, which is impossible; so $x\ne0$ and $\Rel(A\bar B)=x^2>0$. The odd case is the
same with $A=-B=\ell_{\mathrm{odd}}(F)$. Now let $F=F_{\mathrm e}+F_{\mathrm o}\in\mathcal Z_T$ be arbitrary. By
Theorem~\ref{thm:parityschur}(o), $W_T(F)=W_T(F_{\mathrm e})+W_T(F_{\mathrm o})$ with
$W_T(F_{\mathrm{par}})\ge\lambda^{\mathrm W}_{T,\mathrm{par}}\|F_{\mathrm{par}}\|^2\ge\lamW\|F_{\mathrm{par}}\|^2$, so each nonzero component
is a minimizer in its sector, and if the sector bottoms differ the component in the sector with the larger bottom
vanishes. Hence, when $\lambda^{\mathrm W}_{T,\mathrm{ev}}\ne\lambda^{\mathrm W}_{T,\mathrm{odd}}$, all of $\mathcal Z_T$ lies in one
sector, the first coordinates of $\mathfrak m(\mathcal Z_T)$ are all positive or all negative, $0\notin\operatorname{conv}\mathfrak m(\mathcal Z_T)$,
and (iii), (ii) give $\lamW<\lamP$ and the statement about the multipliers.
(v) $\lamW\le\lamP\le\lamA$ (Theorem~\ref{thm:A}, Remark~\ref{rem:Kprim}) and (ii).
\end{proof}

\section{Lower bounds on finite windows}\label{sec:log5}

Put $T_5:=\frac12\log5=0.804718\ldots$. Since $\NT=\{n\ge2:\log n<2T\}$, the contact set is
$\{2,3,4\}$ for $\frac12\log4\le T\le T_5$, and $5$ enters only for $T>T_5$; at $T_5$,
$\Aconst=v_2+v_3+v_4=2.94186\ldots$. The window $T_5$ is therefore the largest one with this contact set,
and it lies slightly beyond the range $T\le0.8$ of \cite[Thm.~1.2]{Zhu2026}; the ranges $T\le1$ and $T\le17/16$ are claimed, with much
smaller constants $2^{-151}$ and $2^{-49162}$, in \cite[Thms.~A, B]{Liu2026}. The notation
$\chi_{\mathrm{par}},\ell_{\mathrm{par}},\varepsilon_{\mathrm{par}},\mathcal Q^{\mathrm{par}}_T,
\mathfrak g_{\mathrm{par}},\mathfrak G_T$ is that of \S\ref{sec:complete}, and
$\lambda^{\mathrm W}_{T,\mathrm{par}}$ is the bottom of $W_T$ on the unit sphere of
$\mathcal D^{\mathrm{par}}_T$ (Theorem~\ref{thm:parityschur}(o)).

\begin{theorem}[the primitive form at $\frac12\log5$]\label{thm:log5A}
$\lambda^{\mathrm{ev}}_{\mathrm A}(T_5)\ge2.5\cdot10^{-12}$ and
$\lambda^{\mathrm{odd}}_{\mathrm A}(T_5)\ge3.0\cdot10^{-10}$. Hence
$\lamA\ge2.5\cdot10^{-12}$ for every $0<T\le\frac12\log5$.
\end{theorem}

\noindent The proof includes a computer verification in ball arithmetic (\S\ref{subsec:encllog5}).

\begin{proof}
By Theorem~\ref{thm:parityschur}(o), $\lambda_{\mathrm A}(T_5)$ is the minimum of the two
parity bottoms, which are bounded below by Proposition~\ref{prop:lowercert} with $c=0$ and the
certificate data of \S\ref{subsec:encllog5}. The bound for $T<T_5$ follows from
Lemma~\ref{lem:monoT}.
\end{proof}

\begin{theorem}[the complete form at $\frac12\log5$]\label{thm:log5W}
$W_{T_5}\ge10^{-18}\|\cdot\|^2$ on $\mathcal D^{\mathrm{ev}}_{T_5}$ and $W_{T_5}\ge10^{-16}\|\cdot\|^2$ on
$\mathcal D^{\mathrm{odd}}_{T_5}$. Hence $\lamW\ge10^{-18}$ for every $0<T\le\frac12\log5$, and
$\mathfrak G_T\succ0$ there.
\end{theorem}

\noindent The proof includes a computer verification in ball arithmetic (\S\ref{subsec:encllog5}).

\begin{proof}
By Theorem~\ref{thm:parityschur}(o), $\lambda^{\mathrm W}_{T_5}=\min_{\mathrm{par}}\lambda^{\mathrm W}_{T_5,\mathrm{par}}$,
and each parity is bounded below by Proposition~\ref{prop:lowercert} with $c=-\frac12\varepsilon_{\mathrm{par}}$.
Lemma~\ref{lem:monoT} gives the bound for $T<T_5$. Since $\lamW>0$, Theorem~\ref{thm:parityschur}(c)
gives $\mathfrak G_T\succ0$.
\end{proof}

\begin{corollary}[indeterminacy for $T\le\frac12\log5$]\label{cor:KEPlog5}
$\mathrm{KEP}(T)$ is indeterminate for every $0<T\le\frac12\log5$.
\end{corollary}
\begin{proof}
By Theorem~\ref{thm:log5W}, $\lamW>0$ for $T\le T_5$, so Theorem~\ref{thm:KEPdich} applies for
$\frac12\log2\le T\le T_5$. For $0<T<\frac12\log2$, note that $\mathcal W_{T'}\subseteq\Wcl$ for $T'>T$ (a weight
flat on $[-2T',2T']$ is flat on $[-2T,2T]$, and the decay condition does not involve $T$), so a weight
$w$ with $\inf f_w>0$ at $T'=\frac12\log2$ is admissible at $T$, and $f_w$ does not depend on the window
by \eqref{eq:fw}; the perturbation of Theorem~\ref{thm:KEP1}(ii) applies, its $\chi$ being supported in
$\{|u|\ge\log2\}$, outside the window.
\end{proof}

\subsection{The time--band limiting operator}
\begin{lemma}[a Ky Fan ceiling for the time--band limiting operator]\label{lem:KyFan}
Let $S_c=Q_TB_sQ_T$ on $L^2(\IT)$, where $Q_T$ is restriction to $\IT$ and $B_s$ the band projection
to $|\tau|<s$, and put $c=Ts$. Let $j_n$ be the spherical Bessel functions and
$G_k(x)=\sum_{j<k}(2j+1)j_j(x)^2\in[0,1]$. Then, in the orthonormal Legendre basis
$e_j(t)=\sqrt{(2j+1)/2T}\,P_j(t/T)$,
\[
\langle S_ce_j,e_j\rangle=\frac{2(2j+1)}\pi\int_0^cj_j(x)^2dx,\qquad
\sum_{j\ge0}\langle S_ce_j,e_j\rangle=\frac{2c}\pi=\operatorname{tr}S_c ,
\]
and for every $k\ge0$
\[
\lambda_k(S_c)\ \le\ \Lambda_k(c):=\frac2\pi\int_0^c\big(1-G_k(x)\big)\,dx .
\]
In particular $\Lambda_0(c)=2c/\pi$ is the trace and $\Lambda_k(c)=O(c^{2k+1})$ as $c\to0$.
\end{lemma}
\begin{proof}
From $\int_{-1}^1P_j(u)e^{-ixu}du=2(-i)^jj_j(x)$ one gets
$\widehat{e_j}(\tau)=\sqrt{(2j+1)/2T}\,\cdot2T(-i)^jj_j(\tau T)$, so
$|\widehat{e_j}(\tau)|^2=2T(2j+1)j_j(\tau T)^2$ and Plancherel gives the displayed diagonal entry.
The addition theorem $\sum_{j\ge0}(2j+1)j_j(x)^2=1$ gives the trace identity, so the basis loses no
trace. By the Ky Fan maximum principle $\sum_{j<k}\lambda_j(S_c)\ge\sum_{j<k}\langle
S_ce_j,e_j\rangle$; since $S_c=(B_sQ_T)^*(B_sQ_T)\ge0$ its eigenvalues are nonnegative and
non-increasing, whence
$\lambda_k\le\sum_{n\ge k}\lambda_n=\operatorname{tr}S_c-\sum_{j<k}\lambda_j\le\sum_{j\ge k}\langle
S_ce_j,e_j\rangle=\Lambda_k(c)$.
\end{proof}

\begin{corollary}[quadrature-free form]\label{cor:P1}
$|j_n(x)|\le x^n/(2n+1)!!$ for $x\ge0$, hence
$\lambda_k(S_c)\le\Lambda_k(c)\le\Lambda^*_k(c):=\frac2\pi\sum_{j\ge k}c^{2j+1}/\big((2j+1)!!\big)^2$, a convergent
series of rationals in $c$: the bound is evaluated by arithmetic alone, with no quadrature.
\end{corollary}


\begin{lemma}[bordered inertia]\label{lem:border}
Let $q$ be a real quadratic form on a real vector space $V$, $f$ a linear functional on
$V$ and $v_0\in V$ with $f(v_0)=1$ and $q(\cdot,v_0)$ finite. On $V\oplus\R$ put
$b(v,\xi)=q[v]+2\xi f(v)$. Then $n_-(b)=n_-(q|_{\ker f})+1$, $n_-$ being the largest dimension of a
subspace on which the form is negative definite.
\end{lemma}
\begin{proof}
Write $v=w+tv_0$ with $w\in\ker f$. Then $b=q[w]+2tq(w,v_0)+t^2q[v_0]+2\xi t$, and the invertible
substitution $\eta=\xi+q(w,v_0)$ gives $b=q[w]+(t^2q[v_0]+2t\eta)$: a $b$-orthogonal sum of
$q|_{\ker f}$ and a binary form of determinant $-1$, whose inertia is $(1,1)$.
\end{proof}

\subsection{The certificates at \texorpdfstring{$T=\frac12\log5$}{T = (1/2) log 5}}\label{subsec:encllog5}
Use the Legendre basis and the band projection of Lemma~\ref{lem:KyFan} in each real parity space. Here the lower form is a frequency truncation of the full symbol $\sT$, and the operator subtracted from a
multiple of the identity is indefinite.

\emph{The lower form.} Fix $\tau_N>0$ and a number $s_\infty\le m(\tau_N)-\Aconst$. Since $m$ is increasing
with $m(0)=C_0$ and the comb is at least $-\Aconst$ (Lemma~\ref{lem:mbounds}(i),(iii)), the even function
$\underline s$ equal to $\sT$ on $|\tau|\le\tau_N$ and to $s_\infty$ on $|\tau|>\tau_N$ satisfies
$\underline s\le\sT$. Write $g=s_\infty-\underline s=(s_\infty-\sT)\mathbf 1_{|\tau|<\tau_N}=g_+-g_-$ with
$g_\pm\ge0$, and $\mathcal G_{(\pm)}=Q_Tg_{(\pm)}(D)Q_T$. Then
\[
\mathfrak a_T(F)\ \ge\ s_\infty\|F\|^2-\langle\mathcal GF,F\rangle,\qquad
0\le\mathcal G_\pm\le\omega_\pm S_{c_*},\quad c_*=T\tau_N,
\]
with $\omega_+=s_\infty-C_0+\Aconst$ and $\omega_-=m(\tau_N)+\Aconst-s_\infty$. For $\mu\in\R$ and
$c\in\{0,\frac12,-\frac12\}$ put, on $L^2_{\mathrm{par}}(\IT)\oplus\R$,
\[
\mathfrak b_{\mu,c}(F,\xi)=(s_\infty-\mu)\|F\|^2-\langle\mathcal GF,F\rangle+2\xi\,\ell_{\mathrm{par}}(F)+c\,\xi^2 .
\]

\begin{proposition}[the lower certificate]\label{prop:lowercert}
Let $k_1<k_2$, let $E$ be the span of the $e_j$, $j<k_1$, of parity $\mathrm{par}$, $Y$ its orthogonal
complement in $L^2_{\mathrm{par}}(\IT)$, and $P_Y$ the orthogonal projection onto $Y$.
\emph{(a)} On $Y$, $\mathfrak b_{\mu,c}\ge q_Y:=s_\infty-\mu-\omega_+\Lambda_{k_1}(c_*)$.
\emph{(b)} Suppose $q_Y>0$, and let $\mathsf R$ be the symmetric matrix on $E\oplus\R$ with
\[
\begin{gathered}
\mathsf R_{ij}=(s_\infty-\mu)\delta_{ij}-\langle\mathcal Ge_i,e_j\rangle
-\frac{\langle P_Y\mathcal Ge_i,P_Y\mathcal Ge_j\rangle}{q_Y},\\
\mathsf R_{i\xi}=\langle e_i,\chi_{\mathrm{par}}\rangle+\frac{\langle P_Y\mathcal Ge_i,P_Y\chi_{\mathrm{par}}\rangle}{q_Y},\qquad
\mathsf R_{\xi\xi}=c-\frac{\|P_Y\chi_{\mathrm{par}}\|^2}{q_Y}.
\end{gathered}
\]
Then $n_-(\mathfrak b_{\mu,c})\le n_-(\mathsf R)$, and the same holds for any symmetric $\mathsf R'\le\mathsf R$.
\emph{(c)} If $n_-(\mathsf R')\le1$ with $c=0$, then $\mathfrak a_T\ge\mu$ on the unit sphere of
$\PT^{\mathrm{par}}$. If $n_-(\mathsf R')=0$ with $c=\frac12$ and $\mathrm{par}=\mathrm{odd}$, or $n_-(\mathsf R')\le1$ with
$c=-\frac12$ and $\mathrm{par}=\mathrm{ev}$, then $W_T\ge\mu\|\cdot\|^2$ on $\mathcal D^{\mathrm{par}}_T$.
\end{proposition}
\begin{proof}\sloppy
(a) $-P_Y\mathcal GP_Y\ge-P_Y\mathcal G_+P_Y$ and
$\|P_Y\mathcal G_+P_Y\|\le\omega_+\operatorname{tr}(P^\perp S_{c_*}P^\perp)=\omega_+\Lambda_{k_1}(c_*)$.
(b) For $i<k_1$ and $y\in Y$, $\mathfrak b_{\mu,c}(e_i,y)=-\langle P_Y\mathcal Ge_i,y\rangle$ and the border pairs
with $y$ through $+\langle P_Y\chi_{\mathrm{par}},y\rangle$. Completing the square in the $Y$ variable gives $\mathfrak b_{\mu,c}\cong\mathfrak s\oplus\mathfrak b_{\mu,c}|_Y$ with
$\mathfrak b_{\mu,c}|_Y\ge q_Y>0$, so $n_-(\mathfrak b_{\mu,c})=n_-(\mathfrak s)$, and the subtracted term of the Schur
complement $\mathfrak s$ is at most $q_Y^{-1}$ times the Gram matrix of the $Y$-columns
$-P_Y\mathcal Ge_i$ and $P_Y\chi_{\mathrm{par}}$; hence $\mathfrak s\ge\mathsf R\ge\mathsf R'$.
(c) Put $\mathfrak q_\mu(F)=\mathfrak b_{\mu,0}(F,0)$, so that $\mathfrak a_T-\mu\|\cdot\|^2\ge\mathfrak q_\mu$. For
$c=0$, Lemma~\ref{lem:border} gives $n_-(\mathfrak q_\mu|_{\ker\ell_{\mathrm{par}}})=n_-(\mathfrak b_{\mu,0})-1=0$. For
$c=\frac12$, $\mathfrak b_{\mu,1/2}=\mathfrak q_\mu-2\ell_{\mathrm{par}}^2+\frac12(\xi+2\ell_{\mathrm{par}})^2$, so
$\mathfrak q_\mu-2\ell_{\mathrm{odd}}^2\ge0$, and $W_T=\mathfrak a_T-2\ell_{\mathrm{odd}}^2$ on odd $F$
(Theorem~\ref{thm:parityschur}(o)). For $c=-\frac12$,
$\mathfrak b_{\mu,-1/2}=\mathfrak q_\mu+2\ell_{\mathrm{par}}^2-\frac12(\xi-2\ell_{\mathrm{par}})^2$; the substitution
$\eta=\xi-2\ell_{\mathrm{par}}(F)$ is a congruence, so
$n_-(\mathfrak q_\mu+2\ell_{\mathrm{ev}}^2)=n_-(\mathfrak b_{\mu,-1/2})-1=0$, and $W_T=\mathfrak a_T+2\ell_{\mathrm{ev}}^2$ on even
$F$. The statements pass from $L^2_{\mathrm{par}}(\IT)$ to the form domain because $\mathfrak q_\mu$ is bounded.
\end{proof}

\emph{The computation.} For $i-l$ even,
\[
\langle\mathcal Ge_i,e_l\rangle=\frac{2T}\pi\sqrt{(2i+1)(2l+1)}\,(-1)^{(l-i)/2}\int_0^{\tau_N}g(\tau)j_i(\tau T)j_l(\tau T)\,d\tau
\]
(proof of Lemma~\ref{lem:KyFan}), and the entries with $i-l$ odd vanish. The integral is computed by composite
Gauss--Legendre quadrature with $46$ nodes on panels of length $\frac12$, with the remainder bound for
functions analytic in a Bernstein ellipse \cite[Thm.~19.3]{Trefethen2013}, $\rho=3.2$. The ellipse lies
within distance $0.439<\frac12$ of the panel, while the poles of $\psi(\frac14\pm\frac{i\tau}2)$ lie at
$|\Iml\tau|\ge\frac12$. The modulus on the ellipse is bounded with $|j_n(z)|\le e^{|\Iml z|}$,
$|\cos(z\log n)|\le\cosh(\Iml z\log n)$ and, for $\psi$, the recurrence $\psi(w)=\psi(w+8)-\sum_{k<8}(w+k)^{-1}$
together with Binet's bound $|\psi(u)|\le\log|u|+\frac\pi2+\frac1{2\Rel u}+\frac1{12(\Rel u)^2}$ for
$\Rel u\ge8$. The resulting remainder, $1.5\cdot10^{-43}$ per entry (below $2\cdot10^{-42}$ under either
point-count convention for the rule), is added as a ball to every entry. The inner products over $Y$ are
exact sums over $k_1\le l<k_2$ of the given parity plus tail boxes: by Cauchy--Schwarz for the positive
operators $\mathcal G_\pm$, $\|P^\perp_{k_2}\mathcal G_\pm e\|^2\le\|P^\perp_{k_2}\mathcal G_\pm P^\perp_{k_2}\|
\langle\mathcal G_\pm e,e\rangle\le\omega_\pm^2\Lambda_{k_2}(c_*)$, so
$\|P^\perp_{k_2}\mathcal Ge\|\le(\omega_++\omega_-)\Lambda_{k_2}(c_*)^{1/2}$, and
$\|P^\perp_{k_2}\chi_{\mathrm{par}}\|^2=\sinh T\pm T-\sum_{l<k_2}\langle e_l,\chi_{\mathrm{par}}\rangle^2$ with
$\langle e_l,\chi_{\mathrm{par}}\rangle=\sqrt{(2l+1)/2T}\,2T\,i_l(T/2)$. The tails $\Lambda_k$ are bounded by
$\Lambda^*_k$ of Corollary~\ref{cor:P1}. Everything is midpoint--radius ball arithmetic at $640$ bits
(Arb \cite{Johansson2017}, through \texttt{python-flint}), including the digamma and Bessel values and the
Gauss--Legendre nodes and weights; every Bessel ball is asserted to have radius below $10^{-60}$, and
comparisons are ball predicates.

\emph{Inertia.} In $\mathsf R$ the term $\|P_Y\chi_{\mathrm{par}}\|^2/q_Y$ is of order $10^{-190}$; it is replaced
by an exact upper bound $\delta'$, which gives $\mathsf R'\le\mathsf R$. When $n_-(\mathsf R')=0$ is to be shown,
interval Cholesky of $\mathsf R'$ is applied. When $n_-(\mathsf R')\le1$ is to be shown, a double-precision
Householder matrix $H$ that sends the midpoint of the column $(\mathsf R'_{i\xi})_i$ to a multiple of $e_0$ is
used as an exact matrix; $\|H^{\mathsf T}H-I\|_\infty<1$ is verified in balls, so $\mathrm{diag}(H,1)$ is a
congruence. In the new coordinates the $2\times2$ block $B$ on $(e_0,\xi)$ has $\det B<0$, certified, so its
inertia is $(1,1)$, and by Haynsworth's inertia additivity \cite{Haynsworth1968}
$n_-(\mathsf R')=1+n_-(S)$ for the Schur complement $S$ of $B$, whose positivity is verified by interval
Cholesky. The only non-rigorous objects are $H$ and the midpoints, which enter only through this exact
congruence.

\emph{Data.} $T=\frac12\log5$, $\tau_N=240$, $c_*=193.13$, $k_1=280$, $k_2=330$ ($\dim E=140$ in each parity).
From $m(\tau_N)=3.64276\ldots$: $s_\infty=0.700787608\ldots$, $\omega_+=9.0150$, $\omega_-=5.8839$,
$\Lambda^*_{k_1}(c_*)\le4.4\cdot10^{-20}$, $\Lambda^*_{k_2}(c_*)\le3.9\cdot10^{-70}$, $q_Y\ge0.70079$, tail
box $(\omega_++\omega_-)\Lambda^*_{k_2}(c_*)^{1/2}\le2.9\cdot10^{-34}$; the entry radii of $\mathsf R$ are at
most $4.4\cdot10^{-41}$. The verified cases are:
\begin{center}\small
\begin{tabular}{@{}llccc@{}}
\toprule
form & sector & certified $\mu$ & minimal pivot at that $\mu$ & fails at\\
\midrule
$\mathfrak a_T|\PT$ ($c=0$) & even & $2.5\cdot10^{-12}$ & $8.47\cdot10^{-3}$ & $3\cdot10^{-12}$\\
$\mathfrak a_T|\PT$ ($c=0$) & odd & $3.0\cdot10^{-10}$ & $2.03\cdot10^{-2}$ & $4\cdot10^{-10}$\\
$W_T$ ($c=-\frac12$) & even & $10^{-18}$ & $4.47\cdot10^{-4}$ & $10^{-16}$\\
$W_T$ ($c=\frac12$) & odd & $10^{-16}$ & $8.14\cdot10^{-14}$ & $10^{-14}$\\
\bottomrule
\end{tabular}
\end{center}
\subsection{The margin decays}\label{subsec:margin}
The trivial half of Zhu's conjectural law $-\log\lamW\asymp e^{2T}$ \cite[Conj.~12.1]{Zhu2026}, namely that
the margin tends to $0$ at an explicit rate, can be proved from three facts about the zeros of $\zeta$.
These are also the zero inputs of the subsequent window computation cited in
\S\ref{subsec:power-window}.
\begin{itemize}
\item[(Z1)] $\rho_1=\frac12+i\gamma_1$, $\gamma_1=14.134725141734\ldots$, is a simple zero on the critical line,
and $\zeta$ has no non-trivial zero $\rho\ne\rho_1,\bar\rho_1$ with $|\Iml\rho|\le21$
\cite{RosserYoheSchoenfeld1968,PlattTrudgian2021}. Only $\gamma_1\in[14.13,14.14]$ and the gap $(\gamma_1,21]$
are used.
\item[(Z2)] Every non-trivial zero $\rho$ with $|\Iml\rho|\le H_{\mathrm{PT}}:=3\cdot10^{12}$ lies on the
critical line \cite{PlattTrudgian2021}.
\item[(Z3)] For $t\ge2$, $\big|N(t)-\frac t{2\pi}\log\frac t{2\pi e}-\frac78\big|\le0.137\log t+0.443\log\log t+1.588$,
where $N(t)$ counts the zeros with $0<\Iml\rho\le t$ with multiplicity \cite[Thm.~19]{Rosser1941}.
\end{itemize}
Besides these, we use that the non-trivial zeros lie in $0<\Rel s<1$, and Lemma~\ref{lem:roughEF}.
\begin{lemma}[a crude zero count]\label{lem:Ncount}
$N(t)\le\frac t{2\pi}\log t$ for $t\ge21$.
\end{lemma}
\begin{proof}
$\frac t{2\pi}\log\frac t{2\pi e}=\frac t{2\pi}\log t-\frac{1+\log2\pi}{2\pi}t$ and
$\frac{1+\log2\pi}{2\pi}=0.45166\ldots$, so by (Z3)
$\frac t{2\pi}\log t-N(t)\ge0.45166\,t-\frac78-1.588-0.137\log t-0.443\log\log t$, which is $6.11\ldots$ at
$t=21$ and has derivative $0.45166-0.137/t-0.443/(t\log t)>0$ for $t\ge21$.
\end{proof}

\begin{lemma}[zero sums by counting]\label{lem:zerosum}
For $p\ge2$ and $H\ge21$, with the sum over the ordinates $\gamma=\Iml\rho>0$ of the non-trivial zeros
counted with multiplicity,
\[
\Sigma_p(H):=\sum_{\gamma>H}\gamma^{-p}\ \le\ \frac p{2\pi}\,H^{1-p}\Big[\frac{\log H}{p-1}+\frac1{(p-1)^2}\Big].
\]
\end{lemma}
\begin{proof}
$\sum_{\gamma>H}\gamma^{-p}=\int_{H^+}^\infty t^{-p}dN(t)=\big[t^{-p}N(t)\big]_H^\infty+p\int_H^\infty t^{-p-1}N(t)\,dt$;
the term at $\infty$ vanishes and the one at $H$ is $-H^{-p}N(H)\le0$ by Lemma~\ref{lem:Ncount}, which
also bounds the integral by $\frac p{2\pi}\int_H^\infty t^{-p}\log t\,dt=\frac p{2\pi}H^{1-p}\big[\frac{\log H}{p-1}+\frac1{(p-1)^2}\big]$.
\end{proof}

\begin{lemma}[the test family]\label{lem:testfam}
Fix $T>0$ and an integer $k\ge5$; let $b_k=\frac k2\mathbf 1_{[-1/k,1/k]}$ and $\omega_k=b_k^{*k}$ (the
$k$-fold convolution).
\emph{(i)} $\omega_k\ge0$, $\int\omega_k=1$, $\supp\omega_k\subseteq[-1,1]$, $\omega_k\in C^{k-2}(\R)$ with
$\omega_k^{(j)}(\pm1)=0$ for $j\le k-2$, and $\|\omega_k\|_2^2\ge\frac12$.
\emph{(ii)} $\widehat\omega_k(x)=\big(\sin(x/k)/(x/k)\big)^k$; $|\widehat\omega_k(x)|\le\min\{1,(k/|x|)^k\}$ for real
$x$, and $|\widehat\omega_k(z)|\le(k/|\Rel z|)^ke^{|\Iml z|}$ for complex $z$.
\emph{(iii)} Put $G(t)=\omega_k(t/T)$, $\varpi(c)=(c^2-\frac14)(c^2+\gamma_1^2)$, $\partial=d/dt$, and
\[
F:=\varpi(\partial)G=G''''+(\gamma_1^2-\tfrac14)G''-\tfrac{\gamma_1^2}4G .
\]
Then $F$ is real, supported in $[-T,T]$, bounded, and $F\in\PT$: four integrations by parts, the
boundary terms carrying $G^{(j)}(\pm T)=0$ for $j\le3$, give $\int Fe^{ct}dt=\varpi(c)\int Ge^{ct}dt$
for every $c\in\C$, and $\varpi(\pm\frac12)=0$.
\emph{(iv)} $\Fh(\tau)=(\tau^2+\frac14)(\tau^2-\gamma_1^2)\,T\,\widehat\omega_k(T\tau)$,
$\Phi_F(s)=\varpi(s-\frac12)\Phi_G(s)$ and $\Phi_G(\frac12+\beta+i\gamma)=T\,\widehat\omega_k\big(T(\gamma-i\beta)\big)$;
in particular $\Phi_F(\rho_1)=\Phi_F(\bar\rho_1)=0$.
\emph{(v)} With $Y:=T\gamma_1/\sqrt2$ and $\varepsilon_k(Y):=\frac1\pi k^{2k}Y^{1-2k}/(2k-1)$,
$\|F\|^2\ge\frac{\gamma_1^4}{64}\,T\,\big(\frac12-\varepsilon_k(Y)\big)$ whenever $\varepsilon_k(Y)<\frac12$.
\end{lemma}
\begin{proof}
(i) $\omega_k$ is a spline of degree $k-1$ supported in $[-1,1]$ and vanishing there to order $k-1$;
$1=(\int_{-1}^1\omega_k)^2\le2\int\omega_k^2$. (ii) $\widehat b_k(x)=\sin(x/k)/(x/k)$; $|\sin y/y|\le\min(1,1/|y|)$
on $\R$, and $|\sin(a+ib)|\le\cosh b\le e^{|b|}$. (iii) $G\in C^3$ with $G'''$ absolutely continuous
and $G''''\in L^\infty$ ($k=5$), and $G^{(j)}(\pm T)=0$ for $j\le3$; $\int G''e^{ct}=c^2\int Ge^{ct}$.
(iv) is (iii) with $c=i\tau$, resp.\ $c=s-\frac12$, and $\varpi(\pm i\gamma_1)=0$. (v) On
$|\tau|\le\gamma_1/\sqrt2$ one has $(\tau^2+\frac14)^2\ge\frac1{16}$ and $(\tau^2-\gamma_1^2)^2\ge\gamma_1^4/4$,
so $\|F\|^2=\frac1{2\pi}\int|\Fh|^2\ge\frac{\gamma_1^4}{64}\,T\,\frac1{2\pi}\int_{|x|\le Y}|\widehat\omega_k|^2$, and
$\frac1{2\pi}\int_{|x|\le Y}|\widehat\omega_k|^2\ge\|\omega_k\|_2^2-\frac1\pi\int_Y^\infty(k/x)^{2k}dx\ge\frac12-\varepsilon_k(Y)$.
\end{proof}

\begin{theorem}[the margin decays]\label{thm:margin}
Let $T>0$, $k\ge5$ an integer, $p:=2k-8$, $Y:=T\gamma_1/\sqrt2$, and suppose $\varepsilon_k(Y)<\frac12$. Then
\begin{equation}\label{eq:Bmargin}
\lamW\ \le\ \lamP\ \le\ \lamA\ \le\ \mathcal B(T,k):=
\frac{2\cdot1.0012\cdot T^2(k/T)^{2k}\big[\Sigma_p(21)+e^T\,\Sigma_p(H_{\mathrm{PT}})\big]}
{(14.13^4\,T/64)\,\big(\frac12-\varepsilon_k(Y)\big)},
\end{equation}
with $\Sigma_p(H)$ bounded by Lemma~\ref{lem:zerosum}. Consequently $\limsup_{T\to\infty}\lamA\le0$, and the
same for $\lamP$ and $\lamW$.
\end{theorem}
\begin{proof}
Take $F$ of Lemma~\ref{lem:testfam}. (1) $F\in\PT$, so $\lamA\le\mathfrak a_T(F)/\|F\|^2$, and
$\lamW\le\lamP\le\lamA$ by Theorem~\ref{thm:A} and Remark~\ref{rem:Kprim}.
(2) For $|\tau|\ge21$, $(\tau^2+\frac14)^2\le1.0012\,\tau^4$ and $(\tau^2-\gamma_1^2)^2\le\tau^4$, so
$|\Fh(\tau)|^2\le1.0012\,\tau^8T^2(k/(T|\tau|))^{2k}=O(|\tau|^{8-2k})$ with $8-2k\le-2$: $F$ is in the
form domain, which is hypothesis (a) of Lemma~\ref{lem:roughEF}; hypothesis (b) is verified in step (6).
Hence $\mathfrak a_T(F)=\sum_\rho\Phi_F(\rho)\Phi_F(1-\rho)$.
(3) $\rho_1,\bar\rho_1$ contribute $0$ (Lemma~\ref{lem:testfam}(iv)); by (Z1) no other zero has $|\Iml\rho|\le21$.
(4) A zero with $21<|\gamma|\le H_{\mathrm{PT}}$ lies on the line by (Z2) and contributes
$|\Fh(\gamma)|^2\le1.0012\,T^2(k/T)^{2k}|\gamma|^{8-2k}$.
(5) A zero $\rho=\frac12+\beta+i\gamma$ with $|\gamma|>H_{\mathrm{PT}}$ has $|\beta|<\frac12$;
$|\varpi(\beta+i\gamma)|\le(\gamma^2+\beta^2+\frac14)(\gamma^2+\beta^2+\gamma_1^2)\le(1+10^{-22})\gamma^4$, $\varpi$ is even,
and $|\Phi_G(\rho)|,|\Phi_G(1-\rho)|\le T(k/(T|\gamma|))^ke^{T|\beta|}$ by Lemma~\ref{lem:testfam}(ii),(iv);
so $|\Phi_F(\rho)\Phi_F(1-\rho)|\le1.0012\,T^2(k/T)^{2k}|\gamma|^{8-2k}e^T$.
(6) Summing (4)--(5) over $\pm\gamma$ with Lemma~\ref{lem:zerosum} ($p=2k-8\ge2$) gives
$\sum_\rho|\Phi_F(\rho)\Phi_F(1-\rho)|\le2\cdot1.0012\,T^2(k/T)^{2k}\big[\Sigma_p(21)+e^T\Sigma_p(H_{\mathrm{PT}})\big]<\infty$,
which is hypothesis (b) of Lemma~\ref{lem:roughEF} and bounds $\mathfrak a_T(F)$.
(7) Divide by Lemma~\ref{lem:testfam}(v) with $\gamma_1^4\ge14.13^4$. The limit statement is
Corollary~\ref{cor:margin1}.
\end{proof}

\begin{corollary}[closed form]\label{cor:margin1}
For every $T\ge1$,
\begin{equation}\label{eq:closedmargin}
\lamW\ \le\ \lamP\ \le\ \lamA\ \le\ 4\cdot10^9\,T\,e^{-15.449\,T},\qquad 15.449<\tfrac{42}e=15.4494\ldots\footnotemark
\end{equation}
\footnotetext{The exponent $\frac{42}{e}$ belongs to the test family of Lemma~\ref{lem:testfam}: one factor removes only the first zero, so the cutoff is $21$, just below the second ordinate $\gamma_2=21.022\ldots$, and $\frac{42}e$ comes from optimizing $k$ in $(k/21T)^{2k}$. It is a property of that family, not a sharp rate for $\lamA$.}
\end{corollary}
\begin{proof}
Take $k:=\lceil21T/e\rceil\ge8$ and write $k=21T/e+\theta$, $\theta\in[0,1)$. Then
$(k/21T)^{2k}=\exp\big(2k[\log(1+\theta e/21T)-1]\big)\le\exp(-42T/e+2\theta^2e/21T)\le e^{-15.4494\,T+0.259}$.
Next $(k/T)^{2k}\Sigma_p(21)\le21^9(k/21T)^{2k}\cdot\frac1{2\pi}\frac{2k-8}{2k-9}\big[\log21+\frac1{2k-9}\big]\le21^9\cdot0.5797\,(k/21T)^{2k}$
for $k\ge8$. The bound for the $H_{\mathrm{PT}}$-term is at most $10\,e^T(21/H_{\mathrm{PT}})^{2k-9}$ times that for the first, with
$21/H_{\mathrm{PT}}=7\cdot10^{-12}$ and $2k-9\ge15.449\,T-9$, hence less than $10^{-70}$ of it for $T\ge1$.
Finally $\varepsilon_k(Y)\le0.01566$ for all $T\ge1$: for fixed $k$, $\varepsilon_k$ decreases in $T$, so it
suffices to evaluate it at the left end of each $k$-range, namely $T=1$ for $k=8$ ($\varepsilon_8=0.0060$)
and $T=(k-1)e/21$ for $k\ge9$, where the largest value is $\varepsilon_9=0.01566$ at $T=8e/21$, the values
decreasing thereafter. So $\frac12-\varepsilon_k\ge0.484$, and
$\mathcal B(T,k)\le\frac{2\cdot1.0012\cdot0.5797\cdot21^9e^{0.259}}{(14.13^4/64)\cdot0.484}\,T\,e^{-15.4494\,T}(1+10^{-70})
\le3.97\cdot10^9\,T\,e^{-15.4494\,T}$.
\end{proof}

\begin{remark}[a cited doubly exponential ceiling]\label{rem:moriceiling}
Corollary~\ref{cor:margin1} is singly exponential in $T$ and is uninformative for $T\lesssim1.5$ (its right side
exceeds $1$ there). Mori \cite{MoriUpper} states, for the complete form, $\lamW\le3.50\cdot10^{24}\mu^{9/2}\log\mu\,
e^{-4\pi\mu}$ for $50\le\mu\le10^4$ and $\lamW\le7.58\cdot10^{24}\mu^5(\log\mu)^3e^{-4\pi\mu}$ for $\mu\ge50$, with
$\mu=e^{2T}$, by an explicit prolate test vector; that is, a ceiling of order $\exp(-4\pi e^{2T})$, doubly exponential in
$T$, consistent with the envelope threshold of Proposition~\ref{prop:env}. The proof uses zero data (verification of
the zeros up to a height and a zero-free region). We have not checked it line by line; we record it as plausible,
unverified, and cite it. It bounds $\lamW\le\lamA$ and so does not replace Corollary~\ref{cor:margin1} as a bound on
$\lamA$.
\end{remark}

\begin{corollary}[the structure of the margin]\label{cor:margin2}
\emph{(i)} There is no $c>0$ with $\lamA\ge c$ for all $T$, nor for $\lamP$ or $\lamW$; for $T\ge1$ every
lower bound on any of the three lies below $4\cdot10^9Te^{-15.449T}$.
\emph{(ii)} The limits $L^{\mathrm A}:=\lim_{T\to\infty}\lamA$, $L^{\mathrm{prim}}:=\lim_{T\to\infty}\lamP$ and
$L^{\mathrm W}:=\lim_{T\to\infty}\lamW$ exist in $[-\infty,0]$, and $L^{\mathrm W}\le L^{\mathrm{prim}}\le L^{\mathrm A}$.
In particular \eqref{eq:target} is equivalent to $L^{\mathrm A}=0$.
\end{corollary}
\begin{proof}
(i) is Corollary~\ref{cor:margin1}. (ii) Monotone limits (Lemma~\ref{lem:monoT}) bounded above by
Theorem~\ref{thm:margin}; the last sentence follows from (ii) and monotonicity.
\end{proof}

\subsection{A reported larger window and its inputs}\label{subsec:power-window}
The extended window certificate gives, using (Z1)--(Z3) of \S\ref{subsec:margin},
\begin{equation}\label{eq:power-window}
 W_T(F)\ge e^{-120110026807}\|F\|_2^2,
 \qquad 0<T\le\frac{14011}{5000}=2.8022,
\end{equation}
for every $F$ in the complete window form domain. Here $T$ is the support radius in the additive
variable, not its diameter: $\supp F\subset[-T,T]$, and $N=e^{2T}$.
By Plancherel,
$\int|\widehat F|^2=2\pi\|F\|_2^2$.
On $\PT$ the polar term vanishes, so the same lower bound applies to $\lamA$.

We quote this computer-assisted bound from the accompanying certificate. Its analytic reduction uses the complete
explicit formula, a compactly supported spectral filter, a weighted continuous Schur bound for the
prime-shift operator, and bounds for the corrections due to zeros above
$G=3\cdot10^{12}$. The zero inputs are precisely (Z1), the isolated first pair;
(Z2), the finite-height verification up to $G$; and (Z3), the explicit counting bound with multiplicities.
No hypothesis on zeros above $G$ is imposed: their correction is bounded using (Z3).
The lower bound is uniform over the form domain, rather than a value on a finite test space.

For \eqref{eq:power-window} the recorded parameters are
$U=57/20$, $H=2673507086718$, $R=2014985964$, $m=9115062736$, $b=13$ and
$N_{\mathrm{cut}}=120122483651$. Here $N_{\mathrm{cut}}$ is the filter order, distinct from $N=e^{2T}$.
The verifier checks $17$ inequalities in $192$-bit Arb arithmetic, including the support,
all $2048$ continuous Schur rows, the mask bounds and both correction tails.
All $17$ checks were independently reproduced with python-flint~0.9.0 at $192$ bits.
Smaller radii follow by extension by zero. The endpoint $2.8022$ belongs
to the quantitative certificate \eqref{eq:power-window}. The all-window
sign using \eqref{eq:dirichlet-input} is proved in the conclusion.

\section{Reduction to piecewise exponentials}\label{sec:coer}

This section localizes the positivity question for $\mathfrak a_T$ on $\PT$ to an explicit subspace of
dimension $2k-2$, the \emph{piecewise Tate exponentials}, by proving that the archimedean form is coercive,
with explicit constants, on the complementary subspace of functions that satisfy the two Tate conditions on
each of $k$ consecutive blocks. The coercivity is a high-pass phenomenon ($m(\tau)\sim\log|\tau|$); the prime
form is handled by a norm bound. All forms are on $L^2(\IT)$ with its metric, $\mathfrak a_T=\mT+\PiT$ on
$\mathcal D_T$, and $\mT(F)=C_0\|F\|^2+\Lev(\cF)$ does not depend on $T$ except through its domain. In this
section $r=T/k$ denotes the half-length of a block.

\subsection{Blocks and the Tate-orthogonal subspace}\label{subsec:coer-blocks}

\begin{definition}\label{def:coer-V}
Let $k\ge1$ and let $B_1,\dots,B_k$ be the $k$ consecutive open intervals of length $2r$,
$r:=T/k$, that partition $\IT$ up to endpoints, $B_j=(-T+2(j-1)r,\,-T+2jr)$. Put
\[
\mathcal V_k:=\Big\{F\in L^2(\IT):\ \int_{B_j}F(t)e^{t/2}dt=\int_{B_j}F(t)e^{-t/2}dt=0\ \ (1\le j\le k)\Big\},
\qquad E_k:=\PT\ominus\mathcal V_k ,
\]
and $e_j^\pm:=e^{\pm t/2}\mathbf 1_{B_j}$. (Lemma~\ref{lem:coer-decomp} holds for any partition into
intervals; equal blocks are used from Lemma~\ref{lem:coer-lowpass} on.)
\end{definition}

\begin{lemma}[the decomposition $\PT=\mathcal V_k\oplus E_k$]\label{lem:coer-decomp}
\emph{(i)} $\mathcal V_k$ is a closed subspace of $\PT$. \emph{(ii)} The orthogonal complement of
$\mathcal V_k$ in $L^2(\IT)$ is $\operatorname{span}\{e_j^\pm\}$, of dimension $2k$.
\emph{(iii)} $E_k=\operatorname{span}\{e_j^\pm\}\cap\PT=\{\sum_jc_j^+e_j^++c_j^-e_j^-:\ A=B=0\}$ has
dimension $2k-2$, and $\PT=\mathcal V_k\oplus E_k$ (orthogonal in $L^2(\IT)$).
\emph{(iv)} $E_k\subset\mathcal D_T$, and $\PT\cap\mathcal D_T=(\mathcal V_k\cap\mathcal D_T)\oplus E_k$.
\end{lemma}
\begin{proof}
(i) The $2k$ conditions are continuous linear functionals; summing them over $j$ gives $A(F)=B(F)=0$.
(ii) $\mathcal V_k$ is the joint kernel of the functionals $F\mapsto\langle F,e_j^\pm\rangle$, and the $2k$
functions $e_j^\pm$ are linearly independent (disjoint supports; $e^{t/2},e^{-t/2}$ independent on an
interval). (iii) $E_k=\mathcal V_k^\perp\cap\PT$ since $\mathcal V_k\subset\PT$. The two functionals
$A,B$ restricted to $\operatorname{span}\{e_j^\pm\}$ are independent: on the two-dimensional space
$\operatorname{span}\{e_1^+,e_1^-\}$ their matrix is
$\begin{psmallmatrix}\int_{B_1}e^t&|B_1|\\|B_1|&\int_{B_1}e^{-t}\end{psmallmatrix}$, of determinant
$\int_{B_1}e^t\int_{B_1}e^{-t}-|B_1|^2>0$ (Cauchy--Schwarz, strict since $e^t$ is not constant). Hence
$\dim E_k=2k-2$. If $F\in\PT$ and $F=v+w$ with $v\in\mathcal V_k$, $w\in\mathcal V_k^\perp$, then
$w=F-v\in\PT$, so $w\in E_k$. (iv) $\widehat{e_j^\pm}(s)=\big[e^{(\pm\frac12+is)t}/(\pm\frac12+is)\big]_{\partial B_j}$,
so $|\widehat{e_j^\pm}(s)|\le2e^{T/2}(\frac14+s^2)^{-1/2}$ and $\int\log(2+|s|)|\widehat{e_j^\pm}|^2<\infty$.
Since $\mathcal D_T$ is a linear space containing $E_k$, the $\mathcal V_k$-component of an $F\in\PT\cap\mathcal D_T$
lies in $\mathcal D_T$.
\end{proof}

For a window $r>0$ we write $\lambda_{\mathrm A}(r)$ for the bottom of $\mathfrak a_r$ on the unit sphere of
$\mathcal P_r\cap\mathcal D_r$; it is attained (Lemma~\ref{lem:compact}) and continuous in $r$ (Lemma~\ref{lem:contA}).
For $2r\le\log2$ the contact set $\mathcal N_r$ is empty, so $\mathfrak a_r=\mT$ on $L^2(I_r)$ and $\lambda_{\mathrm A}(r)$ is a
purely archimedean quantity. We use $|\cF(u)|\le\cF(0)=\|F\|^2$ and $\cF(u)=0$ for $|u|\ge2T$.

\begin{lemma}[the reduced kernel]\label{lem:coer-K1}
Put $K_1(u):=K(u)-2e^{-u/2}=\dfrac{2e^{-5u/2}}{1-e^{-2u}}=2\sum_{j\ge1}e^{-(2j+\frac12)u}$ for $u>0$, with $K$ as in
Lemma~\ref{lem:K}. Then $0<K_1(u)\le u^{-1}$ for all $u>0$, and $K_1(u)=u^{-1}-\frac32+O(u)$ as $u\downarrow0$. Let
$F,G\in L^2(\IT)$ be real, $\supp F\subset\overline B$, $\supp G\subset\overline{B'}$, with $B,B'$ disjoint intervals and $B'$
to the right of $B$. Then
\begin{equation}\label{eq:coer-cross}
\mT(F+G)=\mT(F)+\mT(G)-\iint F(x)G(y)K(y-x)\,dx\,dy ,
\end{equation}
and if $\int F(x)e^{x/2}dx=0$ or $\int G(y)e^{-y/2}dy=0$, the kernel $K$ in \eqref{eq:coer-cross} may be replaced by $K_1$.
\end{lemma}
\begin{proof}
The three expressions for $K_1$ agree by Lemma~\ref{lem:K} ($K=2e^{-u/2}\sum_{j\ge0}e^{-2ju}$). $K_1\le u^{-1}$ is
$e^{5u/2}-e^{u/2}\ge2u$: the function $g(u)=e^{5u/2}-e^{u/2}-2u$ has $g(0)=g'(0)=0$ and
$g''=\frac{25}4e^{5u/2}-\frac14e^{u/2}>0$. The expansion follows from $\frac1{2\sinh(u/2)}=u^{-1}-\frac u{24}+O(u^3)$,
$\frac1{2\cosh(u/2)}=\frac12+O(u^2)$ and $2e^{-u/2}=2-u+O(u^2)$. For \eqref{eq:coer-cross}, apply \eqref{eq:Aform} with
$\NT=\emptyset$ to $F+G$: $c_{F+G}(u)=\cF(u)+c_G(u)+\int F(t)G(t-u)dt+\int G(t)F(t-u)dt$; for $u>0$ the first cross
integral vanishes ($t\in B$, $t-u\in B'$ is impossible), while the second equals $\iint F(x)G(y)\mathbf 1_{y-x=u}$; integrating
against $K(u)du$ over $u>0$ gives \eqref{eq:coer-cross}, the term $\|F+G\|^2=\|F\|^2+\|G\|^2$ having no cross part.
Finally $2e^{-(y-x)/2}=2e^{-y/2}e^{x/2}$, so $\iint F(x)G(y)\,2e^{-(y-x)/2}dxdy=2\int Fe^{x/2}\int Ge^{-y/2}$, which vanishes
under either Tate condition.
\end{proof}
The cross term in \eqref{eq:coer-cross} carries the factor $1$ for the unordered pair $\{B,B'\}$; the polarized form is
$\mT(F,G)=-\frac12\iint F(x)G(y)K(y-x)\,dxdy$.

\begin{lemma}[a block sector is a short window]\label{lem:coer-sector}
Let $B\subset\IT$ be an interval of length $2r$ with $2r\le\log2$. Then
\[
\inf\Big\{\mathfrak a_T(F):\ F\in\mathcal D_T,\ \supp F\subset\overline B,\ \int_BFe^{\pm t/2}=0,\ \|F\|=1\Big\}=\lambda_{\mathrm A}(r),
\]
independently of $T$ and of the position of $B$; the infimum is attained. In particular, for $F\in\mathcal V_k$ with
$r=T/k$, $2r\le\log2$, and $F_j:=F\mathbf 1_{B_j}$,
\begin{equation}\label{eq:coer-diag}
\sum_{j=1}^k\mT(F_j)\ \ge\ \lambda_{\mathrm A}(r)\,\|F\|^2 .
\end{equation}
\end{lemma}
\begin{proof}
Let $B=I_r+a$. For $F$ supported in $\overline B$ and $F_a(t):=F(t+a)$, supported in $\overline{I_r}$, one has $c_{F_a}=\cF$,
hence $\mT(F_a)=\mT(F)$ by \eqref{eq:Aform}, and $\int F_ae^{\pm t/2}=e^{\mp a/2}\int Fe^{\pm t/2}$, so the two Tate
conditions on $B$ are those of $\mathcal P_r$ for $F_a$. Since $\cF(\log n)=0$ for $\log n\ge2r$ and $\log2\ge2r$, also
$\PiT(F)=0$, so $\mathfrak a_T(F)=\mT(F)=\mathfrak a_r(F_a)$ and the infimum is $\lambda_{\mathrm A}(r)$, attained since
$\lambda_{\mathrm A}(r)$ is. \eqref{eq:coer-diag} is the sum over the blocks, each $F_j$ satisfying the Tate conditions of
$B_j$, together with $\sum_j\|F_j\|^2=\|F\|^2$.
\end{proof}


\subsection{A lower bound for the short window}\label{subsec:coer-short}

\begin{lemma}\label{lem:coer-K}
$K(u)\ge u^{-1}$ for $0<u\le2$.
\end{lemma}
\begin{proof}
By Lemma~\ref{lem:K}, $K(u)-u^{-1}=\frac{u-2\sinh(u/2)}{2u\sinh(u/2)}+\frac1{2\cosh(u/2)}$. From $\tanh x\le x$,
$2\sinh(u/2)\le u\cosh(u/2)$, so $u-2\sinh(u/2)\ge-u(\cosh(u/2)-1)=-2u\sinh^2(u/4)$, and the first term is
$\ge-\frac{\sinh^2(u/4)}{\sinh(u/2)}=-\frac12\tanh\frac u4\ge-\frac u8$. Hence $K(u)-u^{-1}\ge\frac1{2\cosh(u/2)}-\frac u8\ge0$
whenever $u\cosh(u/2)\le4$, which holds for $u\le2$ ($2\cosh1=3.09$).
\end{proof}

\begin{lemma}[the tail]\label{lem:coer-tail}
For $r>0$, $\int_{2r}^\infty K(u)\,du\ge\log\frac2r+\frac\pi2-r$.
\end{lemma}
\begin{proof}
By Theorem~\ref{thm:S}(iii) the integral is $\log\coth\frac r2+\pi-2\arctan e^r$; $\coth x\ge x^{-1}$ bounds the first
term, and $\frac{d}{dr}(\pi-2\arctan e^r)=-\frac1{\cosh r}\ge-1$ with value $\frac\pi2$ at $r=0$ bounds the second.
\end{proof}

\begin{lemma}[the truncated logarithm]\label{lem:coer-w}
Put $\mathsf w(s):=\int_0^2\frac{1-\cos sv}v\,dv=\log(2|s|)+\gamma-\operatorname{Ci}(2|s|)$ for $s\ne0$, $\mathsf w(0)=0$. Then $\mathsf w$ is
even, increasing in $|s|$, and for every $0<R\le\pi$ and all $s\in\R$, $\mathsf w(s)\ge\mathsf w(R)\min\{1,s^2/R^2\}$.
\end{lemma}
\begin{proof}
$\mathsf w'(s)=\int_0^2\sin(sv)\,dv=\frac{1-\cos2s}s\ge0$ for $s>0$. Next $\mathsf w(s)/s^2=\int_0^2v\,\frac{1-\cos sv}{(sv)^2}dv$, and
$x\mapsto\frac{1-\cos x}{x^2}=\frac12\big(\frac{\sin(x/2)}{x/2}\big)^2$ is decreasing on $(0,2\pi)$; for $0<s\le R\le\pi$ and $v\le2$
one has $sv\le2\pi$, so $\mathsf w(s)/s^2\ge\mathsf w(R)/R^2$. For $s\ge R$ use monotonicity.
\end{proof}

\begin{lemma}[a double Poincar\'e inequality]\label{lem:coer-poinc}
Let $h\in L^2(\R)$ with $\supp h\subset[-1,1]$ and $h''\in L^2(\R)$. Then $\|h''\|^2\ge\pi^2\|h'\|^2\ge\frac{\pi^4}4\|h\|^2$.
\end{lemma}
\begin{proof}
$h,h''\in L^2(\R)$ give $h\in H^2(\R)\subset C^1(\R)$, and the support condition gives $h(\pm1)=h'(\pm1)=0$. The second
inequality is the Dirichlet Poincar\'e inequality for $h\in H^1_0(-1,1)$ (first eigenvalue $(\pi/2)^2$). For the first,
$g:=h'\in H^1_0(-1,1)$ with $\int_{-1}^1g=h(1)-h(-1)=0$. Expand $g=\sum_{n\ge1}c_n\phi_n$ in the orthonormal Dirichlet
eigenbasis $\phi_n(x)=\sin\frac{n\pi(x+1)}2$ of $(-1,1)$, so that $\|g\|^2=\sum c_n^2$ and $\|g'\|^2=\sum(\frac{n\pi}2)^2c_n^2$. Since
$\int\phi_n=\frac4{n\pi}$ for odd $n$ and $0$ for even $n$, the constraint reads $c_1=-\sum_{n\ge3\ \mathrm{odd}}c_n/n$. The
even-indexed terms satisfy $(\frac{n\pi}2)^2\ge\pi^2$, and for the odd-indexed ones
\[
\sum_{n\ \mathrm{odd}}\Big(\frac{n^2\pi^2}4-\pi^2\Big)c_n^2=\frac{\pi^2}4\Big[\sum_{n\ge3\ \mathrm{odd}}(n^2-4)c_n^2-3c_1^2\Big]\ \ge\ 0,
\]
because by Cauchy--Schwarz $c_1^2\le\mathcal S\sum_{n\ge3\ \mathrm{odd}}(n^2-4)c_n^2$ with
$\mathcal S=\sum_{n\ge3\ \mathrm{odd}}\frac1{n^2(n^2-4)}=\frac14\big(\frac43-\frac{\pi^2}8\big)=0.0249<\frac13$ (partial fractions and
$\sum_{n\ \mathrm{odd}}n^{-2}=\pi^2/8$). Hence $\|g'\|^2\ge\pi^2\|g\|^2$.
\end{proof}

\begin{theorem}[the short window]\label{thm:coer-A}
Let $0<r\le1$ with $2r\le\log2$, and put $a:=r^2/4$. For every $R\in(0,\pi]$,
\begin{equation}\label{eq:coer-A}
\lambda_{\mathrm A}(r)\ \ge\ C_0+\log\frac2r+\frac\pi2-r+\mathsf w(R)\Big[1-\frac{16(R^2+a)^3}{27\pi^4R^2}\Big].
\end{equation}
In particular, with $R=2$ ($\mathsf w(2)=\log4+\gamma-\operatorname{Ci}(4)=2.1045$),
$\lambda_{\mathrm A}(r)\ge\log\frac1r-1.209-r$ for $0<r\le0.1$, so $\lambda_{\mathrm A}(r)\ge\log\frac1r-1.31$ on $(0,0.1]$.
\end{theorem}
\begin{proof}
Let $F\in\mathcal P_r\cap\mathcal D_r$, $\|F\|=1$. Since $\mathcal N_r=\emptyset$ and $\cF(u)=0$ for $u\ge2r$, \eqref{eq:Aform} gives
\[
\mathfrak a_r(F)=C_0+\int_{2r}^\infty K+\int_0^{2r}\big(1-\Rel\cF(u)\big)K(u)\,du .
\]
The tail is bounded by Lemma~\ref{lem:coer-tail}. Put $G(y):=r^{1/2}F(ry)$, so that $\supp G\subset[-1,1]$, $\|G\|=1$,
$\cF(u)=c_G(u/r)$, and $\int_0^{2r}(1-\Rel\cF)K=\int_0^2(1-\Rel c_G(v))\,rK(rv)\,dv\ge\int_0^2\frac{1-\Rel c_G(v)}v\,dv$, using
$1-\Rel c_G\ge0$ and Lemma~\ref{lem:coer-K} ($rv\le2r\le2$). By Plancherel, $1-\Rel c_G(v)=\frac1{2\pi}\int(1-\cos sv)|\widehat G(s)|^2ds$,
so by Tonelli the last integral equals $\frac1{2\pi}\int\mathsf w|\widehat G|^2$, and by Lemma~\ref{lem:coer-w}, for $0<R\le\pi$,
\[
\frac1{2\pi}\int\mathsf w|\widehat G|^2\ \ge\ \mathsf w(R)\Big[1-\frac1{2\pi}\int\Big(1-\frac{s^2}{R^2}\Big)^+|\widehat G(s)|^2ds\Big].
\]
The Tate conditions $\int Fe^{\pm t/2}=0$ read $\int Ge^{\pm ry/2}=0$, i.e.\ $\widehat G(\mp ir/2)=0$. $\widehat G$ is entire of
exponential type $\le1$ and in $L^2(\R)$, so $H(s):=\widehat G(s)/(s^2+a)$ is entire of exponential type $\le1$ and in $L^2(\R)$
($s^2+a\ge a>0$ on $\R$); by the Paley--Wiener theorem $H=\widehat h$ with $h\in L^2(\R)$, $\supp h\subset[-1,1]$, and
$G=-h''+ah$, so $h''\in L^2$. Then
$1=\|G\|^2=\frac1{2\pi}\int(s^2+a)^2|H|^2=\|h''\|^2+2a\|h'\|^2+a^2\|h\|^2\ge\|h''\|^2\ge\frac{\pi^4}4\|h\|^2$
(Lemma~\ref{lem:coer-poinc}). On the other hand $(1-\frac{s^2}{R^2})^+(s^2+a)^2\le M_R:=\frac{4(R^2+a)^3}{27R^2}$ (the
maximum of $(1-x/R^2)(x+a)^2$ on $0\le x\le R^2$ is at $x=(2R^2-a)/3$), so
$\frac1{2\pi}\int(1-s^2/R^2)^+|\widehat G|^2\le M_R\|h\|^2\le\frac{4M_R}{\pi^4}$. Collecting terms gives \eqref{eq:coer-A}. For
$R=2$ and $a\le\frac1{400}$ the bracket is $\ge0.9025$ and $\mathsf w(2)\cdot0.9025\ge1.899$; with
$C_0+\log2+\frac\pi2=-3.1082$ this gives $\log\frac1r-1.209-r$.
\end{proof}

\begin{remark}\label{rem:coer-A}
(1) The only numerical input is the value $\mathsf w(2)$, an evaluation of an elementary function
($\operatorname{Ci}(4)=-0.1410$). Optimizing $R$ in \eqref{eq:coer-A} numerically gives $R\approx2.17$ and
$\log\frac1r-1.18-r$. (2) Numerically, $\lambda_{\mathrm A}(r)\approx\log\frac1r-0.5$ for $r\in[0.02,0.1]$ (Galerkin upper
bounds, not certified); the loss of the proof is in $\int_{|s|<R}|H|^2\le\|h\|^2$ and in the constant $\pi^4/4$.
(3) Theorem~\ref{thm:coer-A} complements Lemma~\ref{lem:contA} at the left end: $\lambda_{\mathrm A}(r)\to+\infty$ as $r\to0$ at
the rate $\log\frac1r$, consistently with Corollary~\ref{cor:mono}, which gives
$\lambda_{\mathrm A}(r)\le C_0+\varphi((3\pi/2r)^2)=\log\frac1r+O(1)$ (Theorem~\ref{thm:S}(i), Lemma~\ref{lem:mbounds}(ii)).
(4) The theorem is about $\lambda_{\mathrm A}(r)$, the bottom over $\mathcal P_r$ (two Tate conditions on $F$), not about
$\lambda^{\mathrm{prim}}_r$ of Definition~\ref{def:PD}; no statement about the latter is made here.
\end{remark}

\begin{remark}[precedents]\label{rem:coer-prec}
Coercivity of the Weil form on short windows in finite codimension is classical in qualitative form: Yoshida
\cite[Lemma~3]{Yoshida1992} proves that for given $a_0,\mu>0$ there is $N$ such that the Weil form is $\ge\mu\|\phi\|^2$ for
$\phi$ supported in $[-a,a]$, $0<a\le a_0$, orthogonal to $N$ Fourier modes (with no rate in $a$); Connes and Consani
\cite[Thm.~3.6, Cor.~3.8, Rem.~3.9]{ConnesConsani2021} write the archimedean part as $-2\cdot\mathrm{Id}$ plus a
Hilbert--Schmidt operator and obtain positivity on short windows after finitely many linear conditions. The leading
term $\log(1/r)$ of the short-window bottom has been observed numerically for Suzuki's operator without primes in
\cite{Kim2026}. Finite truncations of forms with a difference kernel are also studied in \cite{CvS2025}. What is new in
Theorems~\ref{thm:coer-A} and \ref{thm:coer-B} is the rate $\log(1/r)$ with an explicit constant, under the two Tate
conditions, and the block structure.
\end{remark}


\subsection{Coercivity of the archimedean form on $\mathcal V_k$}\label{subsec:coer-B}
For $R>0$ let $P_R$ be the orthogonal projection of $L^2(\R)$ onto the band $\{|s|\le R\}$,
$\widehat{P_RF}=\Fh\mathbf 1_{|s|\le R}$.

\begin{lemma}[low-pass leakage of $\mathcal V_k$]\label{lem:coer-lowpass}
Let $F\in\mathcal V_k$ with equal blocks of half-length $r$, and $R>0$. Then
\[
\|P_RF\|^2\le\varepsilon_{\mathrm{lp}}(Rr)\,\|F\|^2,\qquad \varepsilon_{\mathrm{lp}}(\theta):=\frac{\cosh^2r}{6}\Big(\theta^2+\frac{r^2}4\Big)^2 .
\]
\end{lemma}
\begin{proof}
$G:=P_RF$ is entire. Fix a block $B$ with centre $c$ and put, on $B$,
$u(x):=G(c)\cosh\frac{x-c}2+2G'(c)\sinh\frac{x-c}2\in\operatorname{span}\{e^{x/2},e^{-x/2}\}$, so that
$u(c)=G(c)$, $u'(c)=G'(c)$, $u''=u/4$. Then $\varrho:=G-u$ satisfies $\varrho(c)=\varrho'(c)=0$ and
$\varrho''-\frac\varrho4=G''-\frac G4=:H$, whence $\varrho(x)=\int_c^x2\sinh\frac{x-t}2\,H(t)\,dt$. By Cauchy--Schwarz and
$2\sinh\frac z2\le z\cosh\frac z2\le z\cosh r$ ($0\le z\le r$),
$|\varrho(x)|^2\le\cosh^2r\,\frac{|x-c|^3}3\,\|H\|_{L^2(B)}^2$, so $\int_B\varrho^2\le\cosh^2r\,\frac{r^4}6\|H\|_{L^2(B)}^2$.
Since $F_B:=F\mathbf 1_B$ is orthogonal to $e^{\pm x/2}\mathbf 1_B$, $\langle F_B,G\rangle=\langle F_B,\varrho\rangle\le\|F_B\|\,\|\varrho\|_{L^2(B)}$.
Summing over the blocks and using Cauchy--Schwarz once more,
\begin{align*}
\|G\|^2=\langle F,G\rangle&\le\|F\|\Big(\sum_B\|\varrho\|_{L^2(B)}^2\Big)^{1/2}\le\|F\|\cosh r\,\frac{r^2}{\sqrt6}\,\|G''-\tfrac14G\|_{L^2(\R)}\\
&\le\|F\|\cosh r\,\frac{r^2}{\sqrt6}\Big(R^2+\frac14\Big)\|G\|,
\end{align*}
the first equality because $P_R$ is an orthogonal projection, the last inequality because
$\widehat G$ is supported in $|s|\le R$. Squaring gives the claim, as $r^4(R^2+\frac14)^2=((Rr)^2+\frac{r^2}4)^2$.
\end{proof}

\begin{theorem}[coercivity on $\mathcal V_k$]\label{thm:coer-B}
Let $k\ge1$, $r=T/k\le1$, $\kappa:=\cosh^2r$, $a:=r^2/4$, $\varepsilon_{\mathrm{lp}}$ as in
Lemma~\ref{lem:coer-lowpass}, and $\theta_*:=(\sqrt{6/\kappa}-a)^{1/2}$ (the root of $\varepsilon_{\mathrm{lp}}(\theta_*)=1$); $m_{\mathrm{lo}}$ is
the lower bound of Lemma~\ref{lem:mbounds}(ii).
Then for every $F\in\mathcal V_k\cap\mathcal D_T$,
\begin{equation}\label{eq:coer-B}
\mT(F)\ \ge\ \Big(\log\frac1r-C_B(r)\Big)\|F\|^2,\qquad
C_B(r):=\log\frac1r-C_0\,\varepsilon_{\mathrm{lp}}(r)-\frac{2\kappa}3\big[\Phi_r(\theta_*)-\Phi_r(r)\big],
\end{equation}
where $\Phi_r$ is the elementary primitive
\[
\Phi_r(\theta):=\Big(\frac{\theta^4}4+\frac{a\theta^2}2\Big)\big(\log\theta-\log(2\pi r)\big)-\frac{\theta^4}{16}-\frac{a\theta^2}4
-r\Big(\frac{\theta^3}3+a\theta\Big)-\frac{r^2\theta^2}4-\frac{ar^2}2\log\theta
\]
of $(\theta^3+a\theta)\,m_{\mathrm{lo}}(\theta/r)$. Moreover $C_B(r)\to\frac14-\log\frac{6^{1/4}}{2\pi}=1.6399$ as
$r\to0$, and $C_B(0.02)=1.6573$, $C_B(0.0619)=1.6957$, $C_B(0.1)=1.7331$ (evaluations of the closed form).
\end{theorem}
\begin{proof}
Let $\|F\|=1$ and $\mu:=\frac1{2\pi}|\Fh|^2ds$, a probability measure, so that $\mT(F)=\int m\,d\mu$ by
\eqref{eq:form} with $\NT=\emptyset$. By Lemma~\ref{lem:mbounds}(i) and \eqref{eq:m}, $m$ is even, real-analytic,
increasing on $[0,\infty)$ with $m(0)=C_0$; hence $m(s)=C_0+\int_0^{|s|}m'(\rho)d\rho$ with $m'\ge0$ and, by Tonelli,
\[
\int m\,d\mu=C_0+\int_0^\infty m'(\rho)\,\mu(|s|\ge\rho)\,d\rho\ \ge\ C_0+\int_0^{\rho_*}m'(\rho)\big(1-\varepsilon_{\mathrm{lp}}(\rho r)\big)d\rho,
\qquad\rho_*:=\theta_*/r,
\]
since $\mu(|s|\ge\rho)=1-\|P_{\rho}F\|^2+\mu(|s|=\rho)\ge1-\varepsilon_{\mathrm{lp}}(\rho r)$ by Lemma~\ref{lem:coer-lowpass} and the
integrand is nonnegative on $[0,\rho_*]$, where $\varepsilon_{\mathrm{lp}}(\rho r)\le1$. Integrating by parts on $[0,\rho_*]$, with
vanishing boundary terms ($1-\varepsilon_{\mathrm{lp}}(\theta_*)=0$ and $m(0)-C_0=0$),
\[
\int_0^{\rho_*}m'(\rho)(1-\varepsilon_{\mathrm{lp}}(\rho r))d\rho=\int_0^{\theta_*}\big(m(\theta/r)-C_0\big)\varepsilon_{\mathrm{lp}}'(\theta)\,d\theta,
\]
so that
\[
\int m\,d\mu\ \ge\ C_0\,\varepsilon_{\mathrm{lp}}(0)+\int_0^{\theta_*}m(\theta/r)\,\varepsilon_{\mathrm{lp}}'(\theta)\,d\theta ,
\]
using $\int_0^{\theta_*}\varepsilon_{\mathrm{lp}}'=1-\varepsilon_{\mathrm{lp}}(0)$. Here $\varepsilon_{\mathrm{lp}}'(\theta)=\frac{2\kappa}3\theta(\theta^2+a)\ge0$. By
Lemma~\ref{lem:mbounds}(i),(ii), $m\ge\max\{C_0,m_{\mathrm{lo}}\}$ on $(0,\infty)$; we use $m\ge C_0$ on $[0,r]$ (i.e.\ $\rho\le1$)
and $m\ge m_{\mathrm{lo}}$ on $[r,\theta_*]$ (note $m_{\mathrm{lo}}(1)=-3.34>C_0$, and $m_{\mathrm{lo}}$ is increasing, so no further loss occurs
there). Hence
\[
\int m\,d\mu\ \ge\ C_0\,\varepsilon_{\mathrm{lp}}(0)+C_0\big(\varepsilon_{\mathrm{lp}}(r)-\varepsilon_{\mathrm{lp}}(0)\big)+\int_r^{\theta_*}m_{\mathrm{lo}}(\theta/r)\,\varepsilon_{\mathrm{lp}}'(\theta)\,d\theta
=C_0\,\varepsilon_{\mathrm{lp}}(r)+\frac{2\kappa}3\big[\Phi_r(\theta_*)-\Phi_r(r)\big],
\]
which is \eqref{eq:coer-B}; the primitive is checked by differentiation. As $r\to0$: $\kappa\to1$, $a\to0$,
$\theta_*\to6^{1/4}$, $\varepsilon_{\mathrm{lp}}(r)\to0$, $\frac{2\kappa}3\Phi_r(r)\to0$, and
$\frac{2\kappa}3\Phi_r(\theta_*)=\log\theta_*-\log(2\pi r)-\frac14+o(1)$ (the $a$- and $r$-terms are $o(1)$),
whence $C_B(r)\to\log2\pi-\log6^{1/4}+\frac14$. The three values are evaluations of the closed form (double precision).
\end{proof}

\begin{remark}\label{rem:coer-B}
(1) Using $m\ge\max\{C_0,m_{\mathrm{lo}}\}$ with the exact crossover instead of the split at $\rho=1$ changes $C_B$ by less
than $5\cdot10^{-5}$ on $(0,0.1]$ (numerical). (2) Numerically $C_B(r)\le1.74$ on $(0,0.1]$; monotonicity of $C_B$ is not
proved and is not used, since Corollary~\ref{cor:coer-L1} needs $C_B$ only at the specific value $r=T/k$. (3) The bound
loses about $1$ against the Galerkin value of the bottom of $\mT$ on $\mathcal V_k$ (about $1.1$ against $2.1$ at
$T=\frac12\log5$, $k=13$; numerical); the loss is in Lemma~\ref{lem:coer-lowpass}. (4) The theorem uses the whole
multiplier $m$ and the Tate conditions only through the geometry of $\mathcal V_k$; it does not separate the block-diagonal
part \eqref{eq:coer-diag} from the coupling \eqref{eq:coer-cross}. Remark~\ref{rem:coer-nogo} explains why that separation
is not viable.
\end{remark}




\begin{proposition}[the norm of the prime form]\label{prop:coer-prime}
\emph{(i)} $|\PiT(F)|\le\Aconst\|F\|^2$ for all $F\in L^2(\IT)$ and all $T$.
\emph{(ii)} For $\ell>0$ let $L(T,\ell):=\lceil2T/\ell\rceil$. Then $|\Rel\cF(\ell)|\le\cos\frac{\pi}{L(T,\ell)+1}\,\|F\|^2$,
hence
\[
\|\PiT\|\ \le\ \sum_{n\in\NT}v_n\cos\frac\pi{L(T,\log n)+1}
=\begin{cases}
\Aconst/2, & T\le\log2,\\[2pt]
\dfrac{v_2}{\sqrt2}+\dfrac{v_3+v_4}2=1.6740, & \log2<T\le\tfrac12\log5 .
\end{cases}
\]
Here $\frac{v_2}{\sqrt2}+\frac{v_3+v_4}2=\log2+\frac{\log3}{\sqrt3}+\frac{\log2}2$; in the first case the bound is
$\frac{v_2}2=0.4901$ for $T\le\frac12\log3$ and $\frac{v_2+v_3}2=1.1244$ for $\frac12\log3<T\le\log2$. The bounds hold on all of $L^2(\IT)$, in particular on $\mathcal V_k$.
\end{proposition}
\begin{proof}
(i) is $|\cF(\log n)|\le\|F\|^2$. (ii) Let $V_\ell$ be the shift $F\mapsto F(\cdot-\ell)\mathbf 1_{\IT}$ on
$L^2(\IT)$, so that $\Rel\cF(\ell)=\frac12\langle(V_\ell+V_\ell^*)F,F\rangle$. The map
$F\mapsto\big(F(-T+y+j\ell)\big)_{0\le j<n(y)}$, $y\in[0,\ell)$, $n(y):=\#\{j\ge0:-T+y+j\ell<T\}\le L(T,\ell)$, is unitary
from $L^2(\IT)$ onto $\int^\oplus_{[0,\ell)}\C^{n(y)}dy$, and it carries $V_\ell$ to the fibrewise shift $j\mapsto j+1$.
Hence $V_\ell+V_\ell^*$ acts on the fibre $\C^{n}$ as the adjacency matrix of the path with $n$ vertices, whose
eigenvalues are $2\cos\frac{j\pi}{n+1}$, $1\le j\le n$; its norm is $2\cos\frac\pi{n+1}\le2\cos\frac\pi{L+1}$. For
$T\le\log2$ one has $2T/\log n\le2$ for all $n\ge2$, so $L=2$ and $\cos\frac\pi3=\frac12$. For $\log2<T\le\frac12\log5$:
$\NT=\{2,3,4\}$, $L(T,\log2)=3$ since $2<2T/\log2\le\log5/\log2=2.32$, and $L(T,\log3)=L(T,\log4)=2$ since
$2T\le\log5<2\log3$; $\cos\frac\pi4=\frac1{\sqrt2}$. The values use $v_2=\sqrt2\log2$, $v_3=2\log3/\sqrt3$, $v_4=\log2$.
\end{proof}
For $\frac12\log2<T<\frac12\log3$ the bound $\Aconst/2$ is the exact norm: on a set of fibres of positive
measure the fibre has length $2$ and the prime form acts there as $-v_2\begin{psmallmatrix}0&1/2\\1/2&0\end{psmallmatrix}$.
Once $T>\log2$ the orbit of $\log2$ has length $3$ and the path norm exceeds $1$, so the bound
$\|\PiT\|\le\Aconst/2$ is specific to $T\le\log2$.

\begin{corollary}[coercivity of the truncated form on $\mathcal V_k$]\label{cor:coer-L1}
Let $k\ge1$ with $r=T/k\le1$, and put $c(T,k):=\inf\{\mathfrak a_T(F):F\in\mathcal V_k\cap\mathcal D_T,\ \|F\|=1\}$. Then
\begin{equation}\label{eq:coer-L1}
c(T,k)\ \ge\ \log\frac kT-C_B\Big(\frac Tk\Big)-\|\PiT\| ,
\end{equation}
with $\|\PiT\|$ bounded by Proposition~\ref{prop:coer-prime}. In particular (evaluations of the closed form):
\begin{center}
\begin{tabular}{@{}lllll@{}}
\toprule
$T$ & $k$ & $\dim E_k$ & bound for $\|\PiT\|$ & lower bound for $c(T,k)$\\
\midrule
$\frac12\log3$ & $8$ & $14$ & $0.4901$ & $0.486$\\
$\log2$ & $12$ & $22$ & $1.1244$ & $0.035$\\
$0.69$ & $14$ & $26$ & $1.1244$ & $0.202$\\
$\frac12\log5$ & $25$ & $48$ & $1.6740$ & $0.094$\\
$\frac12\log5$ & $26$ & $50$ & $1.6740$ & $0.134$\\
$\frac12\log5$ & $30$ & $58$ & $1.6740$ & $0.281$\\
$\frac12\log5$ & $85$ & $168$ & $\Aconst=2.9420$ & $0.070$\\
\bottomrule
\end{tabular}
\end{center}
Moreover, for every $T>0$ there is $k_0(T)$ such that $c(T,k)>0$ for all $k\ge k_0(T)$; indeed the right-hand side
of \eqref{eq:coer-L1} tends to $+\infty$ with $k$, because $C_B(T/k)\to1.6399$ (Theorem~\ref{thm:coer-B}).
\end{corollary}
\begin{proof}
$\mathfrak a_T=\mT+\PiT$ on $\mathcal D_T$, Theorem~\ref{thm:coer-B} and Proposition~\ref{prop:coer-prime}.
\end{proof}

\begin{remark}\label{rem:coer-size}
The size of $k_0(T)$: by \eqref{eq:coer-L1} and Proposition~\ref{prop:coer-prime}(i), $c(T,k)>0$ as soon as
$\log\frac kT>C_B(T/k)+\Aconst$, so any $k$ with $k>Te^{\Aconst+C_B(T/k)}$ is admissible, and $C_B(T/k)\to1.6399$ as
$k\to\infty$ (Theorem~\ref{thm:coer-B}). Since $\Aconst\asymp e^T$ by the prime number theorem, the dimension $2k-2$ of $E_k$
obtained in this way grows doubly exponentially in $T$; on the range $T\le\frac12\log5$ it stays below $60$
(Corollary~\ref{cor:coer-L1}). The Galerkin values of $c(T,k)$ exceed the bounds by about $1$: for instance
$c(\frac12\log5,13)\approx1.2$ (numerical), where \eqref{eq:coer-L1} gives a negative number.
\end{remark}




\subsection{The finite reduction}\label{subsec:coer-schur}

\begin{theorem}[reduction to the piecewise Tate exponentials]\label{thm:coer-schur}
Fix $T>0$ and $k\ge1$ with $c:=c(T,k)>0$ (Corollary~\ref{cor:coer-L1}: every $k\ge k_0(T)$). Put
$\mathcal V:=\mathcal V_k\cap\mathcal D_T$ and $M:=\Aconst-C_0$.
\emph{(i)} $\mathfrak a_T$ is an inner product on $\mathcal V$ for which $\mathcal V$ is complete, and for every
$e\in E_k$ the functional $v\mapsto\mathfrak a_T(e,v)$ is bounded on $(\mathcal V,\mathfrak a_T)$; hence there is a unique
$\mathsf Ke\in\mathcal V$ with $\mathfrak a_T(\mathsf Ke,v)=\mathfrak a_T(e,v)$ for all $v\in\mathcal V$. The map $\mathsf K:E_k\to\mathcal V$ is linear,
and $\|\mathsf K\|:=\sup\{\|\mathsf Ke\|:\|e\|=1\}<\infty$.
\emph{(ii)} For $v\in\mathcal V$ and $e\in E_k$,
\begin{equation}\label{eq:coer-schur}
\begin{split}
\mathfrak a_T(v+e)&=\mathfrak a_T(v+\mathsf Ke)+\mathfrak s_{T,k}(e),\\
\mathfrak s_{T,k}(e)&:=\mathfrak a_T(e)-\mathfrak a_T(\mathsf Ke)=\mathfrak a_T(e)-\sup_{u\in\mathcal V}\big[2\Rel\mathfrak a_T(e,u)-\mathfrak a_T(u)\big].
\end{split}
\end{equation}
Let $S_{T,k}$ be the matrix of the real quadratic form $\mathfrak s_{T,k}$ in an orthonormal basis of $E_k$ (a
real symmetric $(2k-2)\times(2k-2)$ matrix) and $\lambda_S$ its least eigenvalue. Then
\begin{equation}\label{eq:coer-sandwich}
\frac{\min\{c,\lambda_S\}}{(1+\|\mathsf K\|)^2}\ \le\ \lamA\ \le\ \min\{c,\lambda_S\}.
\end{equation}
\emph{(iii)} Consequently $\lamA>0$ if and only if $S_{T,k}\succ0$, and $\lamA\ge0$ if and only if $S_{T,k}\succeq0$.
\end{theorem}
\begin{proof}
(i) By Lemma~\ref{lem:mbounds}, $\sigma_T\ge C_0-\Aconst$, so $\mathfrak a_T+M\|\cdot\|^2\ge0$ on $\mathcal D_T$, and
by Lemma~\ref{lem:mbounds}(ii) and $m\ge C_0$ there are constants $c_1,c_2>0$ (depending on $T$) with
$c_1\int\log(2+|s|)|\Fh|^2\le\mathfrak a_T(F)+(M+1)\|F\|^2\le c_2\int\log(2+|s|)|\Fh|^2$ on $\mathcal D_T$. On $\mathcal V$,
$\mathfrak a_T\ge c\|\cdot\|^2$ gives $\mathfrak a_T(v)\ge\frac c{c+M+1}\big(\mathfrak a_T(v)+(M+1)\|v\|^2\big)$, so
$\mathfrak a_T$ is equivalent on $\mathcal V$ to the weighted norm $\int\log(2+|s|)|\widehat v|^2$, for which
$\mathcal D_T$ is complete and $\mathcal V_k\cap\mathcal D_T$ is closed ($\mathcal V_k$ is $L^2$-closed and the weighted norm
dominates the $L^2$ norm). For $e\in E_k\subset\mathcal D_T$ (Lemma~\ref{lem:coer-decomp}(iv)), the Cauchy--Schwarz
inequality for the nonnegative form $\mathfrak a_T+M\langle\cdot,\cdot\rangle$ gives
$|\mathfrak a_T(e,v)|\le(\mathfrak a_T(e)+M\|e\|^2)^{1/2}(\mathfrak a_T(v)+M\|v\|^2)^{1/2}+M\|e\|\|v\|\le C_e\,\mathfrak a_T(v)^{1/2}$,
so $\mathsf Ke$ exists by the Riesz representation theorem; $K$ is linear, and $\|\mathsf K\|<\infty$ since $\dim E_k<\infty$.
(ii) $\mathfrak a_T(v+\mathsf Ke)=\mathfrak a_T(v)+2\Rel\mathfrak a_T(\mathsf Ke,v)+\mathfrak a_T(\mathsf Ke)=\mathfrak a_T(v)+2\Rel\mathfrak a_T(e,v)+\mathfrak a_T(\mathsf Ke)$,
which rearranges to \eqref{eq:coer-schur}; the variational expression is $\mathfrak a_T(\mathsf Ke)=\sup_u[2\Rel\mathfrak a_T(\mathsf Ke,u)-\mathfrak a_T(u)]$
(maximum at $u=\mathsf Ke$). For \eqref{eq:coer-sandwich}: every $F\in\PT\cap\mathcal D_T$ is $v+e$ (Lemma~\ref{lem:coer-decomp}(iv)),
$\|F\|^2=\|v\|^2+\|e\|^2$, and $v\mapsto v+\mathsf Ke$ is a bijection of $\mathcal V$. Taking $e=0$ gives $\lamA\le c$; taking
$v=-\mathsf Ke$ gives $\mathfrak a_T(F)=\mathfrak s_{T,k}(e)$ with $\|F\|^2\ge\|e\|^2$, so $\lamA\le\lambda_S$. Conversely, if
$\lambda_S>0$, then $\mathfrak a_T(v+e)\ge c\|v+\mathsf Ke\|^2+\lambda_S\|e\|^2$, and
$\|v\|^2+\|e\|^2\le(\|v+\mathsf Ke\|+\|\mathsf K\|\|e\|)^2+\|e\|^2\le(1+\|\mathsf K\|)^2(\|v+\mathsf Ke\|^2+\|e\|^2)$, which gives the lower bound; if
$\lambda_S\le0$ the lower bound is implied by $\lamA\le\lambda_S$. (iii) follows from \eqref{eq:coer-sandwich}
and, for the nonstrict case, from $\mathfrak a_T(v+e)\ge\mathfrak s_{T,k}(e)$ together with $v=-\mathsf Ke$.
\end{proof}

\begin{remark}[the inertia tower]\label{rem:inertiatower}
Let $T,k$ be as in Theorem~\ref{thm:coer-schur} and suppose $S_{T,k}$ is nonsingular. Then
$n_-(W_T)=n_-(S_{T,k})+n_-(\mathfrak G_T)$, and the radical of $W_T$ has dimension $n_0(\mathfrak G_T)$ (proved).
Indeed, by \eqref{eq:coer-schur} the map $(w,e)\mapsto w-\mathsf Ke+e$ is a bijection $\mathcal V\times E_k\to\PT\cap\mathcal D_T$
carrying $\mathfrak a_T(w)+\mathfrak s_{T,k}(e)$ to $\mathfrak a_T$, with $\mathfrak a_T\ge c\|\cdot\|^2$ on $\mathcal V$; hence
$n_-(\mathfrak a_T|\PT)=n_-(S_{T,k})$ and the radical of $\mathfrak a_T|(\PT\cap\mathcal D_T)$ is $\{-\mathsf Ke+e:\ e\in\ker S_{T,k}\}=0$.
By compactness of the resolvent, $0$ is then not in the spectrum of the operator of $\mathfrak a_T|(\PT\cap\mathcal D_T)$, nor,
by Theorem~\ref{thm:parityschur}(o), of either sector $\mathcal Q^{\mathrm{par}}_T$; so $\mathfrak G_T$ is defined and
Theorem~\ref{thm:parityschur}(c) gives the claim. The three levels are the coercive block $\mathcal V$ (no negative
squares), the finite matrix $S_{T,k}$ and the $2\times2$ polar block $\mathfrak G_T$. The size of $S_{T,k}$ is $2k-2$ with
$k$ the number of blocks, not the number of contacts: at a contact $T_n$ the matrix varies continuously and is not
bordered, so Cauchy interlacing at contacts is not available in this basis.
\end{remark}



\begin{remark}[what the reduction is and is not]\label{rem:coer-honest}
(1) Theorem~\ref{thm:coer-schur} is the Schur complement of $\mathfrak a_T$ with respect to
$\PT=\mathcal V_k\oplus E_k$, in the form-domain setting required because $\mathfrak a_T$ is unbounded on $L^2$.
Its content is a localization: the sign of $\lamA$ is the sign of an explicit real symmetric matrix of size
$2k-2$, with $k=25$ at $T=\frac12\log5$, indexed by the piecewise Tate exponentials $e^{\pm t/2}\mathbf 1_{B_j}$.
(2) There is no gain in precision. By \eqref{eq:coer-sandwich}, $\lamA\le\lambda_S\le(1+\|\mathsf K\|)^2\lamA$
whenever $\lambda_S\le c$, so $S_{T,k}$ inherits the decay of Theorem~\ref{thm:margin}; at $T=\frac12\log5$ its
least eigenvalue is of the size of $\lambda_{\mathrm A}(\frac12\log5)\approx3.10\cdot10^{-12}$ (numerical). The
entries of $S_{T,k}$ are not in closed form: the correction $\mathfrak a_T(\mathsf Ke)$ is a supremum over $\mathcal V$,
so a Galerkin computation in $\mathcal V$ gives upper bounds for $S_{T,k}$, and a certified lower bound needs a
near-maximizer together with a dual estimate, which is legitimate because $c>0$ is proved.
(3) Finite reductions of the window problem in sufficient (one-directional) form exist in Legendre bases
\cite[Thms.~1.1, 1.4]{Zhu2026}, \cite{Groskin2026}; Theorem~\ref{thm:coer-schur} is an exact reduction in an
explicit basis of explicit dimension, and it does not improve any certified value.
\end{remark}

\begin{remark}[the route by operator norms]\label{rem:coer-nogo}
By Lemmas~\ref{lem:coer-sector} and \ref{lem:coer-K1}, for $F\in\mathcal V_k$ with $2r\le\log2$,
$\mT(F)=\sum_j\mT(F_j)+\mathfrak q(F)$ with $\sum_j\mT(F_j)\ge\lambda_{\mathrm A}(r)\|F\|^2$ and the coupling form
$\mathfrak q(F)=-\sum_{i<j}\iint F_i(x)F_j(y)K_1(|x-y|)\,dxdy$ (the Tate conditions on each block allow $K_1$ in place of $K$),
so that $c(T,k)\ge\lambda_{\mathrm A}(r)-\sup_{\mathcal V_k}|\mathfrak q|/\|\cdot\|^2-\|\PiT\|$. This splitting cannot replace
Theorem~\ref{thm:coer-B}:
\begin{itemize}
\item[(a)] For two adjacent blocks, the bilinear form $(F_1,F_2)\mapsto\iint F_1(x)F_2(y)K_1(|x-y|)dxdy$ on Tate-orthogonal
functions has norm exactly $\pi$. It is $\le\pi$ by $0<K_1(u)\le u^{-1}$ and Hilbert's inequality
$\iint_{(0,\infty)^2}\frac{f(x)g(y)}{x+y}dxdy\le\pi\|f\|\|g\|$ \cite[Thm.~316]{HLP}, in the distances to the common endpoint. It is
$\ge\pi$ by the sequence $f_M(x')=x'^{-1/2}\mathbf 1_{[2re^{-M},2r]}(x')$ in the distance $x'$ to the common endpoint, placed on
both blocks and projected onto the orthogonal complement of $\{e^{\pm x/2}\mathbf 1_B\}$: in logarithmic variables
$\iint f_M(x')f_M(y')/(x'+y')=\int_{[-M,0]^2}\frac{du\,dv}{2\cosh\frac{u-v}2}=\pi M-O(1)$ while $\|f_M\|^2=M$; the bounded part
$K_1-u^{-1}=-\frac32+O(u)$ contributes $O(\|f_M\|_1^2)=O(r)$; and the projection changes $f_M$ by a vector of norm $O(1)$
uniformly in $M$, which changes the pairing by $O(\sqrt M)$. Hence, with the normalization of \eqref{eq:coer-cross},
$\sup_{\mathcal V_2}|\mathfrak q|/\|\cdot\|^2=\frac\pi2$, and $\sup_{\mathcal V_k}|\mathfrak q|/\|\cdot\|^2\ge\frac\pi2$ for every $k\ge2$ and every
$r$.
\item[(b)] With \eqref{eq:coer-A} and $\|\PiT\|\le1.674$ (Proposition~\ref{prop:coer-prime}), the split bound is positive at
$T=\frac12\log5$ only if $\log\frac1r>1.21+r+\frac\pi2+1.674$, which needs $k\ge70$ ($\dim E_k\ge138$), against $k=25$ from
Theorem~\ref{thm:coer-B}. Numerically the true coupling cost at the Galerkin minimizers of $\mT$ on $\mathcal V_k$ is only
about $0.2$: the extremals of the coupling norm are not the minimizers of the form.
\end{itemize}
\end{remark}


\section{The first vanishing of the minimum}
Throughout, $\mathfrak a_T(F,G)$ is the sesquilinear form of \eqref{eq:form}, $H_T$ the operator of
Lemma~\ref{lem:compact}, and $\mathrm{pr}_T$ the orthogonal projection of $L^2(\IT)$ onto $\PT$; $J$ preserves
$\PT$ and commutes with $\mathrm{pr}_T$. The orthogonal complement of $\PT$ in $L^2(\IT)$ is the
two-dimensional \emph{Tate space} $\mathrm{span}\{e^{t/2},e^{-t/2}\}$ (restricted to $\IT$).

\begin{lemma}[first-contact reduction]\label{lem:firstcontact}
If $\lamA<0$ for some $T$, then there is a window $T_*>\frac12\log5$ and a nonzero $F\in\mathcal P_{T_*}$ such
that $\lambda_{\mathrm A}(T)\ge0$ for $T\le T_*$, $\lambda_{\mathrm A}(T)<0$ for $T>T_*$, and
\[
 H_{T_*}\ge0,\qquad H_{T_*}F=0,\qquad JF=\varepsilon F,\quad \varepsilon\in\{\pm1\}.
\]
\end{lemma}
\begin{proof}
Let $T_*=\inf\{T:\lamA<0\}$. By Lemma~\ref{lem:monoT}, $\lamA<0$ for every $T>T_*$ and $\lamA\ge0$ for every
$T<T_*$; by Lemma~\ref{lem:contA}, $\lambda_{\mathrm A}(T_*)=0$. Since $\lambda_{\mathrm A}(\frac12\log5)>0$
(Theorem~\ref{thm:log5A}), monotonicity gives $T_*>\frac12\log5$. By Lemma~\ref{lem:compact} the value $0$ is
an eigenvalue of $H_{T_*}$, and since $\ker H_{T_*}$ is $J$-invariant it contains a nonzero vector of
definite parity.
\end{proof}
\section{A common family of tests for all windows}
There is one sequence of matrices which suffices to test the sign for every
window. Its entries can be computed before choosing any zeros. For $a>1$
write
\[
 L(a)=\frac1a+\frac1{a-1}-\frac12\log\pi+
 \frac12\psi(a/2)+\frac{\zeta'(a)}{\zeta(a)},\qquad
 K(a,b)=\frac{L(a)+L(b)}{a+b-1}.
\]
For real tests $h$ on the half-line, extended by zero, write $\mathcal W[h]$
for \eqref{eq:W} with all prime powers included. We use it only where
the integral, moments and prime series converge absolutely. Put
$h_a(y)=e^{-(a-1/2)y}\mathbf1_{y>0}$. The polarized form satisfies
$\mathcal W[h_a,h_b]=K(a,b)$. Indeed the positive shifts
in the symmetrized autocorrelation contribute the two Dirichlet series for
$\zeta'/\zeta$, divided by $a+b-1$. The moments contribute
$1/((a-1)b)+1/(a(b-1))$. For the archimedean integral the partial fractions
of the two Fourier transforms give Poisson kernels of widths $a-1/2$ and
$b-1/2$, divided by $a+b-1$. Poisson integration of $m$ supplies the
digamma and constant terms displayed above.

\begin{proposition}\label{ladder}
If $[K(32+i,32+j)]_{0\le i,j\le m}$ is positive semidefinite for every
$m\ge0$, then $W_T\ge0$ on $\mathcal D_T$ for every $T>0$.
\end{proposition}
\begin{proof}
Translate a smooth compact real test $h$ into $(0,\infty)$. Translation
preserves the autocorrelation and the product of the Tate moments. For
$x=e^{-y}$ put $f(x)=x^{-63/2}h(-\log x)$ on $(0,1)$. It extends smoothly
by zero at both ends. Choose polynomials $P_j\to f$ in $C^1[0,1]$ by
approximating the derivative and integrating. The approximants
$e^{-63y/2}P_j(e^{-y})$ on $y>0$ belong to the spans in the statement.

For clarity, put $\epsilon_j=\|P_j-f\|_{C^1}$. The error on $y>0$ and its
first derivative there are bounded by $C\epsilon_j e^{-63y/2}$.
After zero extension its transform is bounded by
$C'\epsilon_j/(1+|\tau|)$, by integration by parts including the boundary
value at zero. This gives convergence for the logarithmically weighted
Fourier integral and the exponential bound gives convergence of both Tate
moments. All approximants and the limit are bounded by a common multiple
of $e^{-63y/2}$ on the half-line. Their autocorrelations at $\log n$
therefore differ by $O(\epsilon_j n^{-63/2})$, summable against $v_n$.
Thus their energies tend to $W_T$ of the original test. The matrix
hypothesis gives a nonnegative limit. Splitting into real and imaginary
parts and using the form core finishes the proof.
\end{proof}

The accompanying finite-order estimates give the following comparisons. For up to four distinct real exponents $a>1$ the matrix $K$ is positive
definite. For up to thirteen real poles of rates
$a-1/2\ge63/2$ the complete form exceeds $\|h\|_2^2/4000$ on nonzero
rational Laplace data with proper numerator, including repeated poles
. Proposition~\ref{ladder} requires the sign also
when the degree grows without bound.

\clearpage
\section{Conclusion}
The comparison with zeros determines precisely what an all-window
intersection inequality would prove. The statement concerns the radial
classes on the colimit of the semilocal correspondence algebras.

\begin{theorem}
The following are equivalent:
\begin{enumerate}
\item The Riemann hypothesis holds for the Riemann zeta function.
\item $W_T(F)\ge0$ for every $T>0$ and every $F\in\mathcal D_T$.
\item $\lambda_{\rm A}(T)\ge0$ for every $T>0$.
\end{enumerate}
\end{theorem}
\begin{proof}
Assume first that every nontrivial zero has real part $1/2$. For a real
smooth test $F$ we have $1-\rho=\bar\rho$ and
$\Phi_F(\bar\rho)=\overline{\Phi_F(\rho)}$. Equations \eqref{comparison}
and \eqref{eq:W} now give
\[
 W_T(F)=\sum_\rho|\Phi_F(\rho)|^2\ge0.
\]
Absolute convergence was established above. Splitting over real and
imaginary parts gives the assertion for complex smooth tests, and
approximation in the form norm gives (2). On a primitive test the polar
term vanishes, giving (3).

Conversely, if there is a zero off the line, the detection theorem \cite[Theorem~F]{Paper45} supplies a real smooth compactly
supported primitive test $F$ for which the sum in \eqref{comparison} is
strictly negative. Choose $T$ with $\supp F\subset I_T$. Both Tate moments
vanish, hence $\mathfrak a_T(F)=W_T(F)<0$. Normalizing $F$ contradicts (3).
\end{proof}

Consider the following condition on $\zeta$:
\begin{equation}\label{eq:dirichlet-input}
 \zeta(s)\ne0\qquad (\Re s>\tfrac12,\ s\ne1).
\end{equation}

There is a one-sided arithmetic estimate sufficient for this condition.
With $\gamma_E$ Euler's constant, put
\[
 G(x)=\sum_{n\le x}\Lambda(n)(x/n-1)-x\log x+(1+\gamma_E)x,
 \qquad x\ge1.
\]
Suppose that for every $1/2<\theta<1$ there exist $C_\theta,X_\theta$
such that $G(x)\le C_\theta x^\theta$ for $x\ge X_\theta$.
Then \eqref{eq:dirichlet-input} follows. We give the argument because
the direction of this bound matters.

Fix $\theta$ and increase $C=C_\theta$ so that $H(x)=Cx^\theta-G(x)$
is nonnegative on $[1,\infty)$. This is possible by continuity on
compact intervals. For $\Re s>1$ termwise integration gives
\[
 M(s):=\int_1^\infty H(x)x^{-s-1}\,dx
 =\frac C{s-\theta}
 +\frac{\zeta'(s)}{s(s-1)\zeta(s)}
 +\frac1{(s-1)^2}-\frac{1+\gamma_E}{s-1}.
\]
Indeed the integral of $(x/n-1)x^{-s-1}$ from $n$ to infinity is
$n^{-s}/[s(s-1)]$, and the Euler product gives absolute convergence
there. At $s=1$ both principal parts cancel by the Laurent expansion
of $\zeta$. The right side is regular at every real $s>\theta$:
for $0<s<1$ the alternating zeta series is positive by grouping
successive pairs, and equals $(1-2^{1-s})\zeta(s)$.

For completeness, if the integral had real abscissa of convergence
$b>\theta$, its continuation would be regular on a disk of radius
$r>0$ about $b$. Choose $0<d<r/4$. The Taylor series at $b+d$
converges at $b-d$. Its terms there are the integrals of
$H(x)x^{-b-d-1}(2d\log x)^k/k!$, by differentiation inside the
convergence half-plane. Monotone convergence would then give a finite
integral at $b-d$, a contradiction. Thus $M$ is holomorphic on
$\Re s>\theta$. By meromorphic continuation the displayed identity
holds there. A zero of $\zeta$ away from $1$ would give a pole of
$\zeta'/\zeta$ of residue its multiplicity, impossible in this
half-plane. Taking $\theta$ arbitrarily close to $1/2$ proves the claim.
The outstanding arithmetic step in this argument is the upper estimate
for $G(x)$ at all sufficiently large $x$ for each $\theta$.


Writing $\psi(u)=\sum_{n\le u}\Lambda(n)$ makes the cancellation
required of the primes precise. For $1\le x<y$,
\[
 \frac{G(y)}y-\frac{G(x)}x=\int_x^y\frac{\psi(u)-u}{u^2}\,du .
\]
Indeed differentiation gives the integrand between consecutive
integers, and each new summand starts at zero. The classical prime
number theorem with remainder $\psi(u)-u=O(u/\log^2 u)$ implies
that $G(y)/y$ has a finite limit $L$. Its value follows from the
regularity at $1$ of the Mellin transform of $G$ computed above:
$h\int_1^\infty (G(y)/y)y^{-1-h}\,dy\to0$ as $h\downarrow0$.
This limit equals $L$, since the weight has mass one and its mass on
each compact interval tends to zero. Thus the desired upper bound
is equivalent to
\[
 \int_x^\infty\frac{\psi(u)-u}{u^2}\,du\ge -C_\theta x^{\theta-1}.
\]
It suffices to establish this lower bound with upper limit $2x$ for
all sufficiently large $x$: summing over $[2^jx,2^{j+1}x]$ costs
the factor $(1-2^{\theta-1})^{-1}$. The prime number theorem used
to identify the integration constant gives only a logarithmic saving.

In fact the analytic implication also holds with a weaker arithmetic
hypothesis. Fix $1/2<\theta<1$. Suppose there are arbitrarily large $Y$
such that
\[
 G(x)\le C_\theta x^\theta\qquad(Y\le x\le Y^{3/2})
\]
with the same $C_\theta$ on all these intervals. Then $\zeta$ has no
zero in $\Re s>\theta$.

To see this, write $g(u)=G(e^u)$ and let $P(s)$ be its Laplace
transform. The preceding Mellin calculation gives its meromorphic
continuation to $\Re s>0$. Its poles there are precisely the nontrivial
zeros of $\zeta$, all simple as poles of the logarithmic derivative,
including at a multiple zero. It has no real pole there. The elementary
bound $\Lambda(n)\le\log n$ gives $|g(u)|\ll e^{2u}$ for $u\ge0$.

Put $U=\log Y$, let $\sigma>2$ be fixed, and take
$k=\lfloor5\sigma U/4\rfloor$. Set $A_k(z)=(-1)^kP^{(k)}(z)/k!$.
On $[U,3U/2]$ write $g=g_+-g_-$ with $g_\pm\ge0$. The triangle
inequality for the oscillatory integral, added to the real integral,
cancels the upper estimate contributed by $g_-$. Thus uniformly in $t$,
\[
 |A_k(\sigma+it)|+\Re A_k(\sigma)
 \le {2\max(C_\theta,0)\over(\sigma-\theta)^{k+1}}
 +O\!\left(\int_{u\notin[U,3U/2],\ u\ge0}
          {u^k e^{-(\sigma-2)u}\over k!}\,du\right).
\]
Here the first term follows by extending the integral of the positive
majorant to the whole half-line. To bound the tails put $a=\sigma-2$.
After multiplication by $a^{k+1}$ the tail integral is the probability
that a gamma random variable $V$ of shape $k+1$ and rate $a$ lies
outside the indicated interval. For all sufficiently large $\sigma$
and then $U$, one has $aU/(k+1)<0.9$ and $3aU/(2(k+1))>1.1$.
For $v=0.9,1.1$ the corresponding lower or upper tail of $aV/(k+1)$
is at most $(v e^{1-v})^{k+1}$: apply Markov's inequality to
$\exp((1-1/v)aV)$, whose expectation is $v^{k+1}$.
Both bases are strictly below one. Hence $\sigma$ can be fixed so large
that both bases times $(\sigma-\theta)/a$ are below one. Along the
resulting unbounded sequence of $k$ the entire right side is
$O((\sigma-\theta)^{-k})$.

Suppose now there were a pole $\rho$ with $\Re\rho>\theta$. Write
$D$ for the distance from $\sigma$ to the nearest pole in $\Re s>0$,
truncated at $\sigma-\theta$. There is a centre $z=\sigma+it$ with a
unique nearest pole at distance $d<D$. Indeed, if $D=\sigma-\theta$,
take first $t=\Im\rho$; otherwise move the real centre slightly
toward the ordinate of a nearest pole, which is nonreal. In both cases
the strict inequality persists on an open interval of ordinates.
Ties can be avoided there: two distinct poles to the left of the
vertical line are equidistant at at most one ordinate, and there are
only countably many pairs. Choose $r$ strictly between $d$ and $D$
so the closed disk about $z$ of radius $r$ contains only that pole.
Subtracting its simple principal part and applying Cauchy's estimate
gives $|A_k(z)|\ge c d^{-k}-O(r^{-k})$ with $c>0$.
The real centre is regular throughout its closed disk of radius $r$,
so $|A_k(\sigma)|=O(r^{-k})$. These estimates contradict the window
bound as $k\to\infty$, since $d<r<\sigma-\theta$.
This proves the implication. Its arithmetic hypothesis still requires
a construction of the indicated intervals for every $\theta>1/2$
sufficiently close to $1/2$.


For the application to $\zeta$, the passage from a character family
to the scalar function on these intervals is exact. Indeed, for a
primitive character $\chi$, form the set $S_\chi$ of primitive
characters inducing the powers $\chi^k$, $k\in\mathbb Z$. The level
sum consists of $G(x)$ plus a weighted sum over
$S_\chi\setminus\{\mathbf 1\}$. For the character of modulus one,
every power is trivial and its conductor is one, so the latter set is
empty. Thus an upper bound for this level sum on a window is already
an upper bound for $G$ on the entire same window, with the same
constant and exponent. The remaining requirement is the existence of
arbitrarily large such windows with uniform exponents.

For the trivial Dirichlet character the auxiliary fields in one
cyclotomic construction of these intervals admit an elementary
description. Let $q\equiv3\pmod8$ and $d\equiv7\pmod{16}$ be primes,
$e\ge0$, and $m=2^{e+2}$. In the cyclotomic field of conductor $mqd$,
consider the subgroup of units $a$ satisfying
\[
 a\equiv\pm1\pmod m,\qquad
 \left(\frac{a}{q}\right)=\left(\frac{a}{d}\right)=1.
\]
Its fixed field is
\[
 E=\mathbb Q(\zeta_m+\zeta_m^{-1},\sqrt{-q},\sqrt{-d}),
 \qquad [E:\mathbb Q]=2^{e+2}.
\]
Indeed the Chinese remainder theorem splits the unit group into three
factors. The first subgroup fixes the maximal real subfield of
$\mathbb Q(\zeta_m)$. For $p=q,d$ the quadratic Gauss sum
$\sum_{a\bmod p}(a/p)\zeta_p^a$ transforms by the Legendre character
and its square is $-p$, since $p\equiv3\pmod4$. Thus the stabilizer of
all three generators is precisely the displayed subgroup. Its index is
$(\varphi(m)/2)\cdot2\cdot2=2^{e+2}$. Passing to any cyclotomic
conductor divisible by $mqd$ takes the inverse image of the same
subgroup and leaves its fixed field unchanged. In particular, taking
$e=\lceil(2+8\varepsilon)\log_2 x\rceil$ for $x\ge1$ gives
$4x^{2+8\varepsilon}\le [E:\mathbb Q]<8x^{2+8\varepsilon}$.
This determines the fields and their degrees; a uniform arithmetic
estimate on these varying fields is still needed in the construction.

The parameters controlling the twists can be constructed for a
matrix on a free $\mathbb Z_\ell$-module. Let $S$ have size $r$ and
let $u\in\mathbb Z_\ell^\times$ satisfy $v_\ell(u-1)=c\ge1$,
with $c\ge2$ for $\ell=2$. Fix $j\ge0$ and put $v=u^{\ell^j}$.
For $a,b\in\mathbb Z_\ell$, $a\ne b$, one has
\[
 |v^a-v^b|_\ell=\ell^{-(c+j)}|a-b|_\ell .
\]
For integral exponents this follows by binomial expansion: raising
a principal unit to a power prime to $\ell$ preserves the valuation
of its difference from one, and raising it to $\ell$ increases that
valuation by one under the stated condition on $c$. Extension by
continuity gives the displayed formula and defines the powers on
$\mathbb Z_\ell$.

For each eigenvalue $z$ of $S$ in a splitting field, counted with
multiplicity, choose $s_z\in\mathbb Z_\ell$ minimizing
$|v^s z-1|_\ell$; compactness ensures existence. The eigenvalues
are integral since the characteristic polynomial is monic. If $z$
is a unit, the ultrametric inequality and minimality imply
$|v^a-v^{s_z}|_\ell\le |v^a z-1|_\ell$. If $|z|_\ell<1$,
the right side is one and the inequality still holds. Multiplying
over the eigenvalues gives
\[
 v_\ell(\det(v^a S-I))\le r(c+j)+\sum_z v_\ell(a-s_z),
\]
with $v_\ell(0)=+\infty$. Thus when every distance is at least
$\ell^{-B}$ the determinant is nonzero and its valuation is at most
$r(c+j+B)$. By Smith normal form this bounds the length of the
cokernel on the free module. This also gives a single list of centres for all layers of a
finitely generated module, including its torsion. Let $M$ have at most
$R\ge1$ generators and let $g$ act unipotently on $M/\ell M$.
Choose generators lifting a basis of $M/\ell M$ and write
$\pi:\mathbb Z_\ell^r\twoheadrightarrow M$, $r\le R$, for the
resulting presentation (Nakayama's lemma). Lift the images of the
generators under $g$ to obtain a matrix $A$ with $\pi A=g\pi$.
Its reduction modulo $\ell$ is unipotent since $\pi$ induces an
isomorphism modulo $\ell$. Put $j_0=\lfloor\log_\ell R\rfloor+3$.
Apply the preceding determinant estimate just once, with $j=j_0$ and
$S=A^{\ell^{j_0}}$. Write $H=u^a A$ for an integral twist $a\ge1$.
For $m\le j_0$, $H^{\ell^m}-I$ divides $H^{\ell^{j_0}}-I$ in
the polynomial ring evaluated at $H$, so the earlier cokernel is a
quotient of the later one.

For $m\ge j_0$ put $U=H^{\ell^m}$. Then $U\equiv I\pmod{\ell^2}$:
modulo $\ell$, $(H-I)^r=0$, so $H^{\ell^k}\equiv I\pmod\ell$
once $\ell^k\ge r$, and one further $\ell$th power gives the stated
congruence. Consequently $I+U+\cdots+U^{\ell-1}$ is $\ell$ times
a matrix congruent to $I$ modulo $\ell$. The determinant valuation
of $U^\ell-I$ therefore increases by exactly $r$ whenever
$\det(U-I)\ne0$. The presentation induces a surjection from each
free cokernel onto the corresponding cokernel on $M$. It follows
that outside the balls of radius $\ell^{-B}$ around the same $r$
centres the length on $M$ at every layer $m$ is at most
$r(\max(m,j_0)+c+B)$. The zero module satisfies this directly.
Identification of these module cokernels with arithmetic groups is
a separate descent assertion.

For the upper estimate the direction of the descent map matters:
an injection into a group of bounded length suffices. Here is the
injectivity argument for a nonzero cyclotomic twist. Let $P$ be the
Galois group at a finite layer, $H$ the kernel of its cyclotomic
character, and suppose $P/H$ is topologically generated by $\gamma$
whose character value $u$ has infinite order. Give
$D=\mathbb Q_\ell/\mathbb Z_\ell$ the action through the character
to the power $-e$, where $e\ne0$. Restriction of continuous cocycles
induces an injection
\[
 H^1(P,D)\longrightarrow \operatorname{Hom}_{\rm cont}(H,D).
\]
Indeed coboundaries restrict to zero since $H$ acts trivially. A
cocycle restricting to zero factors through $P/H$. Multiplication
by $u^{-e}-1$ is surjective on $D$, since this scalar is a nonzero
$\ell$-adic integer. Subtracting a coboundary therefore kills its
value on $\gamma$, and hence on all integral powers of $\gamma$.
Continuity and density kill the cocycle everywhere. This injection
restricts to the subgroup satisfying the local conditions. Let $X$ be the maximal abelian pro-$\ell$ quotient of $H$ on
which all inertia subgroups intersected with $H$, and all
decomposition subgroups above $\ell$ intersected with $H$, vanish.
For a class locally trivial at the primes above $\ell$ and
unramified elsewhere, its restriction $f$ factors through $X$:
a local coboundary evaluates to zero on $H$, where the coefficient
action is trivial. Continuity makes the factorization valid through
the closed subgroup generated by the local relations. Write $A$ for
conjugation by a lift of the generator on $X$. The cocycle identity
applied to $gh=(ghg^{-1})g$ gives
$f(ghg^{-1})=\chi(g)^{-e}f(h)$ for $h\in H$, where $\chi$ denotes
the cyclotomic character. Thus $f$ kills $(u^e A-I)X$. At the layer
generated in the cyclotomic quotient by $\gamma^{\ell^m}$ the
same identity gives the operator $(u^e A)^{\ell^m}-I$, using the same quotient $X$. Indeed, in a common algebraic
closure the subgroup $P_m$ of the layer contains $H$, so the kernel
of the restricted cyclotomic character is still $H$. For any
prolongation of a finite place the decomposition subgroup in $P_m$
is the intersection of $P_m$ with its stabilizer in $P$; inertia
is obtained by imposing trivial action on the residue field.
Intersecting either subgroup with $H$ thus gives the original
intersection. Taking all prolongations, with the same residue
characteristic condition above $\ell$, gives identical local relations.
For a finite cokernel its character group into
$\mathbb Q_\ell/\mathbb Z_\ell$ has the same length, by decomposition
into cyclic summands. Hence the injection suffices for the upper
estimate once the uniform generator bound for $X$ is established.

One source of exceptional twists is proximity to a finite list
of $\ell$-adic parameters. Suppose a tower estimate applies to a
twist $j$ whenever $|j-s_i|_\ell\ge\ell^{-B}$ for each of $r$
parameters in a normed extension of $\mathbb Q_\ell$, with
$|\ell|_\ell=\ell^{-1}$. For an integer cutoff $Z\ge1$ take
$B=\lfloor\log_\ell Z\rfloor$. Each open ball of radius $\ell^{-B}$
contains at most one integer in $[1,Z]$. Indeed two distinct such
integers would have difference divisible by $\ell^{B+1}$, by the
ultrametric inequality, while their ordinary distance is less than
$Z<\ell^{B+1}$. Hence at most $r$ twists must be removed, regardless
of the positions of the parameters or the cutoff. In applying this
observation to varying fields, it is the number of parameters in the
tower estimate that must be bounded uniformly.

A finite exceptional set can be avoided on a whole square window
with explicit dependence on its cardinality. Suppose $0<\eta\le A$,
$X>1$, and $S$ is a set of at most $R$ exceptional twist indices.
Put $D=2A/\eta$ and $Y_i=X^{1/(2D^i)}$ for $0\le i\le R$.
For the window $[Y_i,Y_i^2]$ all the indices that could lie in
$(x^\eta,x^A]$ at some point of the window belong to
$(Y_i^\eta,Y_i^{2A}]$. These last intervals are pairwise disjoint:
the upper endpoint for $i+1$ equals the excluded lower endpoint for
$i$. At least one misses $S$. Its left endpoint $Y=Y_i$ satisfies
\[
 X^{1/(2D^R)}\le Y,\qquad Y^2\le X,\qquad
 S\cap(x^\eta,x^A]=\varnothing\quad(Y\le x\le Y^2).
\]
Several arithmetic estimates can use the same window by taking the
union of their exceptional sets before choosing $i$. Their sets may
vary with $X$ and with the auxiliary primes; a bound on the cardinality
of the union independent of these choices keeps $2D^R$ fixed.
Applying this argument requires proving that outside the excluded
indices the arithmetic estimates hold, as well as bounding their number.

For the trivial character the number of primes above $q,d$ in
$\mathbb Q(\mu_{2^\infty})$ is uniformly bounded. In fact it is
exactly two above $q$ and four above $d$. At conductor $2^n$,
$n\ge4$, an odd prime $p$ has residue degree
$\operatorname{ord}_{2^n}(p)$ since Frobenius raises roots of
unity to the $p$th power. Put $b=v_2(p+1)$, equal to $2$ for
$q$ and $3$ for $d$. As $v_2(p-1)=1$, successive squaring gives
$v_2(p^{2^k}-1)=b+k$ for $k\ge1$: after the first square the
factor $p^{2^k}+1$ has valuation one. The order is therefore
$2^{n-b}$, and its index in the unit group is $2^{b-1}$.
The prime counts stabilize, giving the same counts in the union.
Thus allowing ramification at these two rational primes introduces
only six places in this tower, independently of their sizes.

Suppose the space of everywhere unramified quadratic characters
of this fixed infinite field has dimension at most $g$ over
$\mathbb F_2$. Allowing ramification above $q,d$ increases the
dimension by at most six: restriction to the six inertia groups
has unramified kernel, and each restriction contributes at most
one dimension since wild inertia has odd residue characteristic
and the tame inertia has procyclic pro-$2$ quotient.

Adjoining each of $\sqrt{-q},\sqrt{-d}$ at most doubles this
new dimension bound. Indeed take the maximal pro-$2$ Galois group
$P$ unramified outside these allowed places. Each quadratic step is
unramified outside them (the quadratic discriminants are $-q,-d$).
The maximal such extension upstairs lies in the same tower: the
normal closure downstairs of a finite Galois $2$-extension upstairs
has $2$-power degree, since over the quadratic field it is a compositum
of two conjugate $2$-extensions. Unramifiedness is preserved in the
compositum. If $\dim\operatorname{Hom}_{\rm cont}(P,\mathbb F_2)=h$,
then $P$ has $h$ topological generators: lifts of a basis modulo the
closed subgroup of squares and commutators generate, since every
proper closed subgroup is contained in an open subgroup of index two.
For an open subgroup of index two choose one generator outside it,
and multiply the other generators outside it by its inverse.
Rewriting words using the two coset representatives then gives at
most $1+2(h-1)$ generators upstairs (the square of that generator,
the remaining generators and their conjugates). A trivial field step
preserves the bound. After both steps the dimension is at most
$4(g+6)$, uniformly in $q,d$. Imposing extra local vanishing can
only reduce it. This calculation requires a finite $g$ for the fixed
base field.

The passage from square windows to the intervals used in the pole
argument allows a fixed normalization of the degree. Fix $c\ge1$, put
$h=2+8\varepsilon$ and $k=c^{1/h}$. Suppose the arithmetic bounds
hold throughout $[W,W^2]$ for the dyadic exponents $e(u)$ with
$u^h\le2^{e(u)}\le2u^h$, up to twist $\lfloor u^{401}\rfloor$.
Set $Y=kW$. Once $W\ge\max(k^3,k^{400})$, each
$x\in[Y,Y^{3/2}]$ and $u=x/k$ both belong to $[W,W^2]$.
Indeed $(kW)^{3/2}\le W^2$ and $W\le u\le x$. Then
\[
 x^h\le c\,2^{e(u)}\le2x^h,\qquad
 x^{400}=k^{400}u^{400}\le u^{401}.
\]
Thus the same field at $u$ supplies the required degree at $x$,
and any bound on a sum of nonnegative lengths up to the original
cutoff applies up to $\lfloor x^{400}\rfloor$. Its allowed upper
bound $\exp(u^\varepsilon/2)$ can be enlarged to
$\exp(x^\varepsilon/2)$. If $W\ge X^{1/J}$ for arbitrarily large
$X$ with fixed $J\ge1$, then $Y$ is arbitrarily large as well.
The arithmetic input here is uniformity on the original square
window for a single pair of auxiliary primes.

The elementary size comparison in that estimate can be made uniform
over the auxiliary primes. Fix $\varepsilon>0$. Let $2\le Z\le x$ and
suppose
\[
 \log q\le Z^{\varepsilon/5},\qquad
 \log d\le Z^{\varepsilon/5}(\log Z)^2 .
\]
Write $N=mqd$ and $t=\lfloor x^{401}\rfloor$. For each fixed $A,B\ge0$
and all sufficiently large $x$, uniformly over these choices,
\[
 (Nt)^A+B\exp(x^{4\varepsilon/5})
       \le \exp(x^\varepsilon/2).
\]
Indeed $m<8x^{2+8\varepsilon}$, whence
$\log(Nt)\le\log8+(403+8\varepsilon)\log x+
x^{\varepsilon/5}(1+(\log x)^2)=o(x^\varepsilon)$.
Also $\log(1+B)+x^{4\varepsilon/5}=o(x^\varepsilon)$.
Eventually each summand is at most $\exp(x^\varepsilon/4)$,
which gives the asserted inequality. Thus a polynomial estimate of
this shape for the weighted sum of cohomological lengths would give
the requisite numerical bound. Its uniformity in the primes and
the finiteness of the lengths are separate arithmetic assertions.

There is a further simplification in separating the odd primes in
the cohomological sum. For the trivial character the auxiliary prime
$p_*$, the least prime at least $7$ congruent to $3$ modulo $4$ and
coprime to the conductor, is $7$. The auxiliary quadratic field is
$\mathbb Q(\sqrt{-7})$ and has class number one: its discriminant is
$-7$, and Minkowski's ideal bound $(2/\pi)\sqrt7<2$ gives an integral
ideal of norm one in each class. The second auxiliary prime is the
least prime coprime to $2p_*$ and congruent to one modulo the odd part
of the character order. That order is one, so this prime is $3$.
Consequently the exceptional set in the odd-prime count, defined by
adjoining $3$ to the odd prime divisors of the product of the character
order, conductor, $p_*$, class number and second auxiliary prime,
is exactly $\{3,7\}$. Bounds proved away from this set therefore
require two exceptional odd-prime contributions; the prime $2$ is
treated separately.

The same field description removes one branch of the character
count altogether. Write $j$ for the twist index, $1\le j<t$, and
suppose $q,d>2t+1$. Every auxiliary character $\eta$ factors through
$\mathrm{Gal}(E/\mathbb Q)$, of order $2^{e+2}$. Its order is
therefore a power of two. In particular there is no nonprincipal
character of odd order.

More precisely, consider an odd prime $\ell$ with either
$\ell\in\{3,7\}$ or $\ell-1\mid 2j$, taking $q,d>7$ as well.
In the divisibility case $\ell\le2j+1<q,d$; in the exceptional case
$\ell\le7<q,d$. Thus $\ell\nmid N$. Passing from $\eta$ to its
prime-to-$\ell$ order part leaves it unchanged, since its order is
a power of two. Removing its local component at $\ell$ also leaves it
unchanged, since $\ell\nmid N$. Hence the odd-order condition on this
modified character singles out the principal character exactly.
The parity condition $\eta(-1)=(-1)^j$ then forces $j$ even.
For all nonprincipal characters in these ranges the modified character
has even order. The length estimates for those characters, and the
estimate for $\ell\notin\{3,7\}$ with $\ell-1\nmid2j$, still require
arithmetic input.

For $j=1$ the conductor separation at the exceptional primes can be
made explicit. Write $2^c$ for the conductor of the $2$-component
of $\eta$. On removing this component one is left with a product of
quadratic characters modulo $q,d$, of order at most two. The least
common multiple of this order and four is therefore four. The
multiplicative order of $\ell$ modulo four is two for both exceptional
primes. Thus the parameter $v_2(\ell^2-1)$ is $3$ at $\ell=3$ and
$4$ at $\ell=7$. The conductor test
$2^c<16\ell^2 4^{v_2(\ell^2-1)}$ becomes respectively
$c\le13$ and $c\le17$, since the two right sides are $9216$ and
$200704$.

In this range the field cut out by $\eta$ has degree at most $2048$
at $3$ and $32768$ at $7$. Indeed the $2$-component is even and, for
$c\ge3$, factors through the cyclic group
$(\mathbb Z/2^c\mathbb Z)^\times/\{\pm1\}$ of order $2^{c-2}$.
For $c<3$ it is trivial. Adjoining the two quadratic character
components can raise the order of the product only to the least common
multiple of this bound and two. The field cut out by the product has
degree equal to its order. These bounds concern field degrees;
estimating the cohomological lengths in those fields remains necessary.

For these fields the local correction at $\ell$ can be evaluated.
Let $L=\mathbb Q(\eta,\mu_\ell)$, where $\mathbb Q(\eta)$ denotes the
field fixed by the kernel of $\eta$. For $v\mid\ell$, the completion
$L_v$ is an unramified extension of $\mathbb Q_\ell(\mu_\ell)$:
the conductor of $\eta$ divides $N$, which is prime to $\ell$.
The extensions $\mathbb Q_\ell(\mu_{\ell^a})/
\mathbb Q_\ell(\mu_\ell)$ are totally ramified, as follows from the
Eisenstein polynomial $\Phi_{\ell^a}(1+X)$ over $\mathbb Q_\ell$.
Consequently the cyclotomic character on the absolute Galois group
of $L_v$ has image $1+\ell\mathbb Z_\ell$. For every $j\ge1$,
\[
 \operatorname{length}_{\mathbb Z_\ell}
  H^0(L_v,\mathbb Q_\ell/\mathbb Z_\ell(-j))=1+v_\ell(j).
\]
Indeed $1+\ell$ topologically generates this image for odd $\ell$.
The fixed points are the kernel of multiplication by $(1+\ell)^j-1$,
up to a unit. Raising to an exponent prime to $\ell$ preserves the
valuation of $(1+\ell)-1$, and each subsequent $\ell$th power
increases it by one, by the binomial expansion. This proves the formula.

In particular for $j=1$ the sum of local lengths is exactly the number
of places above $\ell$, at most $[L:\mathbb Q]/(\ell-1)$ by the local
degree formula. Since $[L:\mathbb Q]\le(\ell-1)[\mathbb Q(\eta):\mathbb Q]$,
this contribution is bounded by $2048$ at $3$ and $32768$ at $7$
in the small-conductor range. Passing from these local lengths to the
global second cohomology still requires control of the kernel of
localization.

The finite class group calculation arising from cyclotomic descent
has the following concrete form. Fix a base field $L$ containing
$\mu_{\ell^c}$ but not $\mu_{\ell^{c+1}}$, with $\ell$ odd and $c\ge1$.
At level $\nu\ge c$, put $F=L(\mu_{\ell^\nu})$, let $\gamma$ act
on the roots by the exponent $1+\ell^c$, and set $a=(1+\ell^c)^j$.
Such a generator exists since the cyclotomic extension above level $c$
is cyclic of prime-power degree and every nontrivial intermediate
extension contains its first degree-$\ell$ layer. Write $A$ for the
ideal class group of $F$ in additive notation. The number of functions
$f$ from its finite primes to $\mathbb Z/\ell^\nu$ satisfying
\[
 \sum_P v_P(u)f(P)=0\quad(u\in F^\times),\qquad
 f(\gamma^{-1}P)=a f(P)
\]
is exactly
\[
 \left|A\big/\big(\ell^\nu A+(\gamma^{-1}-a)A\big)\right|.
\]
Indeed $f$ extends uniquely to a homomorphism on fractional ideals.
The first equation makes it factor through $A$, and the second makes
it annihilate $(\gamma^{-1}-a)A$. Since the target is killed by
$\ell^\nu$, it factors through the displayed quotient. Conversely
every homomorphism from that quotient satisfies both equations.
Each cyclic summand of the quotient has order $\ell^h$ with
$0\le h\le\nu$, and admits exactly $\ell^h$ homomorphisms to
$\mathbb Z/\ell^\nu$, proving the count. In the fields just considered
$c=1$ by the ramification calculation above. Finiteness of the class
group gives a finite count at each level. To control the global kernel
by this descent one needs a bound independent of $\nu$.

One factor in the class group comparison can be evaluated in every
layer. Put $n=[F:L]=\ell^{\nu-c}$, $r=r_2(L)$, and
$U_F=\mathcal O_F^\times/\mu(F)$. For $a=(1+\ell^c)^j$ as above,
\[
 [U_F:(\gamma-a)U_F]=\frac{(a^n-1)^r}{a-1},
 \qquad
 v_\ell([U_F:(\gamma-a)U_F])
 =r(\nu+v_\ell(j))-(c+v_\ell(j)).
\]
Here the free unit group is written additively. To prove the first
identity, note that $L$ is totally imaginary and each of its $r$
complex places splits into $n$ complex places on which $\gamma$ acts
as an $n$-cycle. By the logarithmic unit theorem the real span of
$U_F$ is the hyperplane of sum zero in this permutation representation.
On the full space the characteristic polynomial is $(X^n-1)^r$;
on the hyperplane it is $(X^n-1)^r/(X-1)$. Evaluating at $a>1$
gives the absolute determinant, hence the index on the integral
lattice. Finally for odd $\ell$ the binomial expansion gives
$v_\ell((1+\ell^c)^h-1)=c+v_\ell(h)$ for $h\ge1$.
Apply this with $h=jn$ and $h=j$ to obtain the second identity.

Thus the $\ell$-part of this unit index grows by a factor
$\ell^r$ whenever $\nu$ increases by one. A uniform estimate for the
class group quotient needs a matching upper estimate on the arithmetic
objects to which the units are compared; the determinant alone
does not give it.

One ingredient in the upper estimate is detection of powers by
reduction. Let $m=\ell^\nu$, assume $\mu_m\subset F$, and let
$\beta\in F^\times$. If $\beta$ is an $m$th power in the residue
field at every prime of $F$ of residue degree one over $L$, apart
from finitely many exceptions, then $\beta\in F^{\times m}$.

Indeed, otherwise $K=F(\beta^{1/m})$ is a nontrivial Galois extension:
all the roots differ by elements of $\mu_m$. A nonidentity automorphism
fixes no root, since fixing any one root fixes $K$. By the Chebotarev
density theorem a positive-density set of unramified primes of $F$
has Frobenius in the conjugacy class of such an automorphism. Exclude
the primes dividing $m$ or the divisor of $\beta$. At every remaining
prime the roots reduce distinctly and Frobenius has no fixed root,
so the reduction of $\beta$ is not an $m$th power. The primes of $F$
of relative residue degree at least two have density zero: their
Dirichlet series for real $s>1$ is bounded by
$[F:L]\sum_{\mathfrak p\subset\mathcal O_L}
(\mathrm N\mathfrak p)^{-2s}$, summed over nonzero prime ideals,
which stays bounded as $s\downarrow1$. Thus the positive-density set
contains a relative degree-one prime outside the exceptional set,
a contradiction. Applied to a quotient of two elements, this proves
that distinct classes modulo $F^{\times m}$ are distinguished by
some such reduction. The argument allows the detecting prime to
depend on the pair of classes and on $\nu$.

The rank factor in the count can be understood from block sums.
Write $D_\ell$ for the homology map induced by $g\mapsto
\operatorname{diag}(g,\ldots,g)$ with $\ell$ blocks. Stabilization by
identity blocks identifies the levels in the following calculation.
For a degree-three class $z$, the two successive differences have
stable image
\[
 P(D_\ell)z=(D_\ell-\ell^3)(D_\ell-\ell^2)z.
\]
In fact in every odd degree $m\ge3$ the image of the $m-1$
successive differences has rational rank at most $r_2(L)$ for $L$
totally imaginary. To see this, start at level $n\ge 2m+4$. The
target level $\ell^{m-1}n$ is at least $4m+5$, so Charney's stability
theorem for Dedekind domains identifies the target homology with
stable homology. Only injectivity at this final level is needed.
In stable rational homology the operator is
$\prod_{h=2}^m(D_\ell-\ell^h)$.

With direct sum as multiplication, stable rational homology is the
free graded-commutative Hopf algebra on its primitives, whose degree
$i$ part is $K_i(\mathcal O_L)\otimes\mathbb Q$ by the rational
Hurewicz theorem for the plus construction. The block map is diagonal
followed by $\ell$-fold direct sum. It therefore multiplies a
primitive by $\ell$, and a product of $h$ primitives by $\ell^h$.
Every monomial in positive degree $m$ has between one and $m$
factors. The displayed polynomial kills all monomials of length at
least two. Its image is contained in the degree-$m$ primitives,
whose dimension is $r_2(L)$ by Borel's rank theorem for a totally
imaginary field (see \cite{CharneyStab,WeibelK}, especially
\cite[IV, \S1]{WeibelK}). This proves the asserted rank bound at
the finite target level.

For a finitely generated abelian group $M$ of rank at most $r$, any
image of $M$ in a group killed by $q=\ell^\nu$ has at most
$q^r|M_{\rm tors}|$ elements, by factoring through $M/qM$ and applying
the elementary divisor theorem. This explains the power of $q$ in
the homological count once the rank estimate is available. The integral torsion estimates needed for the arithmetic count remain to be supplied.

For the integral estimate, one ingredient has a bound independent of
the discriminant. Put $d=[L:\mathbb Q]$ and let $S$ be a finite
subgroup of $\mathrm{GL}_N(\mathcal O_L)$. Then
\[
 |S|\le 3^{dN^2}.
\]
Indeed reduction modulo $3\mathcal O_L$ is injective on $S$.
Otherwise its kernel contains an element $g$ of prime order $p$.
Write $g=1+3^bA$ with $b\ge1$ maximal, using the freeness of the
matrix ring over $\mathbb Z$. If $p\ne3$, the first term of $g^p-1$
is $p3^b A$ and the remaining terms vanish modulo $3^{b+1}$,
a contradiction. If $p=3$, division by $3^{b+1}$ in the cubic
expansion gives $A+3^bA^2+3^{2b-1}A^3$, nonzero modulo three.
Finally the matrix ring modulo three has $3^{dN^2}$ elements.

In particular this bound applies to the stabilizer of an ordered
simplex of positive definite Hermitian forms on the Minkowski space
of $\mathcal O_L^N$. Fixing one vertex already forces the images of a
lattice basis to have prescribed lengths in a positive definite real
norm. Each image therefore has finitely many choices. This bounds
the isotropy contribution in a cell count; the number of cell orbits
and the resulting integral boundary maps must also be controlled to
obtain a discriminant bound for homology torsion.

More explicitly, suppose an augmented acyclic ordered simplicial
chain complex with a group action has at most $A$ simplex orbits in
each degree up to $m+1$, and all stabilizers have order at most $s$.
Here the chains are freely generated by the ordered simplices and
the action preserves their face maps. Then
\[
 \log|H_m(G,\mathbb Z)_{\rm tors}|
 \le A\log(m+2)+A\log(s)\sum_{q=1}^m s^q .
\]
To prove this, tensor the chain complex with a free right bar
resolution over $\mathbb ZG$. Exactness of the augmentation gives
$H_*(G,\mathbb Z)$ as the total homology. Filtering by simplex degree
gives in bidegree $(p,q)$ a direct sum of $H_q(S,\mathbb Z)$ over
orbit representatives with stabilizers $S$, by the induced-module
formula. For $q>0$ transfer kills each summand by $|S|$; the bar
complex gives its cycles at most $|S|^q$ generators, since they form
a subgroup of a free group of that rank. Its order is thus at most
$s^{s^q}$. All subsequent terms in this bidegree are subquotients.

In row zero the first differential is the signed face map on orbit
generators. Every column from degree $m+1$ has Euclidean norm at most
$m+2$, including when some faces coincide. If its rank is $b>0$,
Smith normal form bounds the torsion order of its cokernel by the
absolute value of any nonzero $b$-minor, at most $(m+2)^A$ by
Hadamard's inequality; at rank zero the cokernel is free.
Row-zero homology embeds in this cokernel. No later differential
enters bidegree $(m,0)$, so its final torsion satisfies the same
bound. Multiplying the bounds along total degree $m$ proves the
inequality by the filtration of the abutment. Thus for the arithmetic
group it remains in this step to exhibit such a complex with a
sufficiently small orbit count.

Here is an orbit estimate for the faces that have a common frame
of minimal vectors. Put $b=dN$, $D=|\operatorname{disc}L|$ and
$I=D^N b^b$. Write the form as $Q(v)=\sum_w v_w^*H_wv_w$, summing
over the complex places. Normalize its minimum on nonzero integral
vectors to one. In the Euclidean metric given by $Q$ the lattice has
covolume $2^{-b/2}D^{N/2}\prod_w\det H_w$, by the discriminant formula
and the real Jacobian of the complex linear map $H_w^{1/2}$.
Balls of radius $1/2$ have disjoint interiors modulo this lattice;
each contains a cube of side $1/\sqrt b$. Consequently
$\prod_w\det H_w\ge I^{-1}$.

Suppose the vertices of an ordered $p$-simplex share $N$ minimal
vectors forming an invertible matrix $A$ over $L$. For any one of its
vertices let $G_w=A_w^*H_w A_w$. All diagonal entries are at most one,
so Hadamard's inequality gives $\det G_w\le1$. Moreover
$\prod_w\det G_w=k\prod_w\det H_w$, where
$k=|\mathrm N_{L/\mathbb Q}\det A|$ is the index of the column lattice.
It follows that $k\le I$ and each $\det G_w\ge I^{-1}$.
Two frames give the same left $\mathrm{GL}_N(\mathcal O_L)$-orbit
if their lattices $A^{-1}\mathcal O_L^N$ agree, since then
$A'A^{-1}$ preserves $\mathcal O_L^N$ in both directions.
Such a lattice contains $\mathcal O_L^N$ with index $k$, and modulo
$\mathcal O_L^N$ is a subgroup of $(\mathbb Z/k)^b$. It has at most
$b$ generators by Smith normal form, hence at most $k^{b^2}$
possibilities for fixed $k$. There are at most $I^{1+b^2}$ frame orbits.

For any further minimal vector $v$, put $u=A^{-1}v$. Each $G_w$
has trace at most $N$, so its smallest eigenvalue is at least
$I^{-1}N^{1-N}$. Since $u_w^*G_w u_w\le1$ and $ku$ is integral,
every conjugate of every coordinate of $ku$ has modulus at most
$R=I\sqrt{IN^{N-1}}$. There are at most $(2R+1)^d$ algebraic integers
in this box: distinct integers have distance at least one in the
maximum norm on the complex places, by the field norm, and disjoint
open balls of radius $1/2$ fit in the ball of radius $R+1/2$.
Thus at most $(2R+1)^b$ vectors can be minimal with this frame.
A perfect vertex is determined by $h=(d/2)N^2$ linearly independent
equations $Q(v)=1$, since $h$ is the real dimension of the form
space. The orbit count for ordered tuples of such vertices with a common
frame is therefore at most $I^{1+b^2}(2R+1)^{bh(p+1)}$.

For completeness construct the complex on the dual cone. Let
$\mathcal R$ be the set of positive definite tuples $H$ satisfying
$Q_H(v)\ge1$ for every nonzero integral $v$. It is nonempty by scaling
and convex. For each positive definite dual tuple $X$, minimize
$\sum_w\operatorname{tr}(X_wH_w)$ over $\mathcal R$. Bounded
sublevels bound the traces of all $H_w$ from above, uniformly for
$X$ in compact subsets of the open dual cone. The determinant bound
just proved consequently bounds their least eigenvalues away from
zero. These sublevels are compact and only finitely many lattice
vectors can give active constraints there. The minimizing face is
therefore a compact polytope, whose vertices are perfect forms:
at an extreme point the active linear equations must span the dual,
since otherwise a small displacement in both directions preserves
all constraints. Conversely such a spanning set determines a vertex.

The regions of dual tuples minimized by the various vertices form a
locally finite polyhedral cover of the open dual cone. Indeed on a
compact neighborhood the common trace bound applies by comparison
with a single fixed feasible form. Only finitely many vertices can
occur there by the finite active set, and their regions are cut out
locally by the linear comparisons of their values. All nonempty
intersections are convex. Their nerve is contractible by the
polyhedral nerve theorem (subdivide the locally finite polyhedral
decomposition inside the open cone to obtain a cover by subcomplexes).
For vertices in a simplex of this nerve take their average $H$ and
a common minimizing $X$. The active constraints at $H$ are precisely
their common minimal vectors. If those vectors fail to span $L^N$,
take a nonzero $L$-linear functional vanishing on them and let $P$
be the positive semidefinite tuple of its squared absolute values.
For small $s>0$, $H-sP$ remains feasible by local finiteness of the
constraints. But $\sum_w\operatorname{tr}(X_wP_w)>0$, contradicting
minimality. Thus every nerve simplex has the common frame used above.

Use all ordered tuples supported on nerve simplices, allowing repeats,
with the alternating deletion boundary. This augmented complex has
the homology of the nerve: project a tuple to its oriented simplex
(zero on repetitions) and choose a vertex ordering for a chain section.
The composite is chain homotopic to the identity by induction on
degree inside each support simplex, where prepending any fixed
vertex contracts the augmented tuple complex. The acyclicity so
obtained holds integrally; the section need not respect the group
action. The tuple complex itself carries the required action. This
supplies both the acyclicity and the orbit bound used above.

Two integral constructions enter when passing from divisors to
homology classes. First, let $I_j,J_j$ be nonzero fractional ideals
of $F$, indexed by $1\le j\le n$, and suppose
$\prod J_j=\beta\prod I_j$ as fractional ideals. There is an
$F$-linear map carrying $\bigoplus I_j$ isomorphically to
$\bigoplus J_j$ with determinant exactly $\beta$. Indeed Steinitz
classification gives an isomorphism of the two projective modules,
which extends to $F^n$. Its determinant $b$ satisfies
$b\prod I_j=\prod J_j$ by taking top exterior powers. Hence
$\beta/b$ is a unit, and multiplication by this unit on the first
target summand corrects the determinant. In rank zero the assertion
requires $\beta=1$.

Second, one can isolate a prescribed slope over the integers before
reducing the cycles. Put $q=\ell^\nu$, $d_0=\ell^c$,
$\kappa=1+d_0$, $c\le\nu$, and fix $0\le t_0\le M<d_0$. Write
$P(Z)=\prod_{j=0}^{M-1}(Z-\kappa^j)=\sum_i p_i Z^i$.
In the integral group ring of $(\mathbb Z/q)^2$ there is an element
whose image under $(x,y)\mapsto x+ty$ vanishes for
$0\le t\le M$, $t\ne t_0$, and at $t_0$ equals
$F(X)=\sum_i p_iX^{\kappa^i}$, with $X^q=1$.

To prove this, set $u_i=(\kappa^i-1)/d_0$ and
$H(Z)=\sum_i p_iZ^{u_i}$. For $0\le h<M$ the coefficient of $W^h$
in $H(1+W)$ is $\sum_i p_i\binom{u_i}{h}=0$: the binomial is a
polynomial of degree $h$ in $\kappa^i$, and
$\sum_i p_i(\kappa^i)^j=P(\kappa^j)=0$ for $j<M$.
Thus $(Z-1)^M$ divides $H$ in $\mathbb Z[Z]$, and
$(X^{d_0}-1)^M$ divides $F(X)=XH(X^{d_0})$. For each
$t\ne t_0$ in the range, $\gcd(t_0-t,q)$ divides $d_0$.
Consequently $X^{t_0-t}-1$ divides $X^{d_0}-1$ in the group ring:
choose an integer $n\ge0$ with $n(t_0-t)\equiv d_0\pmod q$
and use the geometric sum. Hence $F$ is a multiple of
$\prod_{t\ne t_0}(X^{t_0}-X^t)$ there, say with multiplier $B(X)$.
Using the two generators $X,Y$ of the larger group ring, the required
element is $B(X)\prod_{t\ne t_0}(Y-X^t)$. Substitution $Y=X^t$
proves both assertions, without division in the residue ring.
For odd $\ell$ the scalar in the corresponding separated weight
has a loss independent of the layer. Indeed for $0\le u<M$ its
integer value and valuation are
\[
 C_u=\prod_{0\le j<M,\ j\ne u}(\kappa^u-\kappa^j),\qquad
 v_\ell(C_u)=(M-1)c+v_\ell(u!)+v_\ell((M-1-u)!).
\]
Each factor has valuation $c+v_\ell(|u-j|)$ by the binomial
valuation formula above. In degree three ($M=3,u=1$) this gives
$C_1=-\kappa\ell^{2c}$. Multiplication by $C_1$ on
$\mathbb Z/\ell^\nu$ therefore has kernel of order
$\ell^{\min(2c,\nu)}$, uniformly bounded with $c$ fixed.
This isolates the determinant component in the coefficient construction;
realizing the resulting data as a cycle with the specified reduction
at every admissible residue field remains to be shown.

The chain that realizes a slope relation has coefficients in
$R=\mathbb Z/q$, $q=\ell^\nu$. Put $C=C_q$, $D=C\times C$, and
$d_t(x)=(x,x^t)$. For any $R$-cycle $z$ of degree $h=2e+2$ in
$C$, there exist coefficients $n_t\in R$, $0\le t\le h+1$, at
least one a unit, and a chain $b$ in $D$ satisfying
\[
 \partial b=\sum_t n_t(d_t)_*z.
\]
Indeed the cyclic integral free resolution has alternating
differentials $g-1$ and $1+g+\cdots+g^{q-1}$. Tensoring two copies
over $\mathbb Z$ resolves the trivial module for $D$: both augmented
complexes split over $\mathbb Z$. After taking coinvariants with
coefficients $R$ both differentials vanish. Hence $H_h(D;R)$ is free
of rank $h+1$, one summand for each bidegree of total degree $h$.
Lift the $h+2$ coordinate columns of the indicated classes to
$\mathbb Z^{h+1}$. They have an integral relation with relatively
prime coefficients, obtained by clearing denominators in a nonzero
rational relation and dividing by their gcd. Reduction gives a
relation with a unit coefficient, and its vanishing in homology
gives the asserted $b$ in the bar complex.

For homomorphisms $\rho,\sigma:D\to H$ suppose that
$\rho d_t=k_t^{-1}(\sigma d_t)k_t$ with $k_t\in H$.
The conjugation prism $P_t$ on bar chains satisfies
$\partial P_t+P_t\partial=(\rho d_t)_*-(\sigma d_t)_*$:
it is the simplicial homotopy of the two one-object group functors
given by conjugation, with the orientation chosen as displayed.
Then the degree-$h+1$ chain
\[
 y=\sigma_*b-\rho_*b+\sum_t n_t P_tz
\]
is a cycle. Applying $\partial$ cancels the two differences term by
term since $\partial z=0$. For $e=0$ this gives a degree-three cycle;
the argument neither requires nor produces a lift to a cycle over
$\mathbb Z$. Application to the divisor construction requires the two
representations in the same integral lattice stabilizer and the
prescribed values on the reduced prisms.

In degree three the scalar evaluation of a reduced prism can be
computed explicitly. Let $k$ be a finite field, $M=|k^\times|$,
$q\mid M$, and $u$ a generator of $k^\times$. For $a\in k^\times$
let $A(a)\in[0,M-1]$ be its discrete logarithm to base $u$. The
integer carry
\[
 B(a,b)=(A(a)+A(b)-A(ab))/M
\]
is a symmetric two-cocycle by associativity. Write the chosen
primitive $q$th root as $\zeta=u^{(M/q)s}$, so $s$ is a unit modulo
$q$. For any homomorphism $v:k^\times\to R$ (written additively),
the cup product $s^{-1}v\smile B$ is a three-cocycle over $R$.
Take $z=\sum_{i=0}^{q-1}[g^i|g]$ for a generator $g$ of $C_q$;
its boundary is $q[g]=0$ over $R$. On inserting $b\in k^\times$
into the image of $z$ under $g\mapsto w=\zeta^x$, the last two
insertion terms cancel by symmetry of $B$. The remaining value is
\[
 s^{-1}v(b)\sum_{i=0}^{q-1} B(w^i,w)
   =s^{-1}v(b)\,q A(w)/M=xv(b)\quad\text{in }R.
\]
The middle equality telescopes since $w^q=1$. In degree $2j+1$,
$j\ge1$, use $s^{-j}v\smile B^{\smile j}$ and the degree-$2j$
cycle obtained by summing
$[g^{i_1}|g|\cdots|g^{i_j}|g]$ over all residues $i_1,\dots,i_j$.
Its first face has coefficient $q=0$; the adjacent internal faces
cancel in pairs by shifting the summation indices, and the last
internal face cancels the last face by the same shift.
On inserting $b$, all terms except the first cancel in pairs:
the two insertions in each carry slot interchange its arguments.
The first term factors into $j$ copies of the telescoping sum.
Its value is therefore $x^jv(b)$.

Here is the extension to matrices. Let $\rho:G\to\mathrm{GL}(V)$
be a representation over $k$. Choose a nonzero vector on each line
in $\mathbb P(V)$. Regard $G$ as a category with one object. A bar
simplex lifts to its action groupoid on the lines once the initial
line is specified; sum over all initial lines. This operation commutes
with the boundary: composing consecutive arrows commutes with lifting,
and deleting the first arrow permutes the starting lines. Send each
lifted arrow to the scalar by which it carries the chosen vector on
its source to the chosen vector on its target. Scalars multiply under
composition, giving a chain map to the bar complex of $k^\times$.
Composing with the preceding cocycle defines an evaluation on cycles.
Rescaling the chosen vectors gives naturally isomorphic functors,
whose chain maps are homotopic by the bar prism. The evaluation is
therefore independent of the choices and invariant under isomorphism.
For a one-dimensional representation it is the original evaluation.

For a block sum $V\oplus W$, send each line with nonzero first
projection to its image in $V$ and the remaining lines to $\mathbb P(W)$.
Over a line in $V$ with chosen vector $v$, the lifts are represented
uniquely by $(v,w)$ with $w\in W$. With these representatives the
scalar for a lifted arrow is precisely that on the first projection.
Each line of $W$ has one lift and its usual scalar. Thus the transfer
of the block sum is the sum of $|k|^{\dim W}$ copies of the first
transfer and one of the second. Since $q\mid |k|-1$, the first
multiplicity equals one in $R$. This proves additivity on cycles,
including the contributions of mixed lines. In positive degree a
trivial representation contributes zero, so stabilization preserves
the evaluation and repeating a block $d$ times multiplies it by $d$.

For a nonselected slope the lattice transport can also be made
explicit. Let $J=(\mathbb Z/q)^2$ and write $\eta_z$ for the integral
coefficient at $z$ of the selector above. For a fractional ideal
$\mathfrak a$ form the lattices
$M_0=\bigoplus_{z\in J}\mathcal O_F$ and
$M_1=\bigoplus_{z\in J}\mathfrak a^{\eta_z}$ in $F^J$.
The group $D$ acts diagonally by
$\rho(u,v)e_z=\zeta^{uz_1+vz_2}e_z$ in additive coordinates and
preserves both lattices. At $t\ne t_0$, every fiber of
$z\mapsto z_1+t z_2$ has coefficient sum zero. The product of its
target ideals is therefore $\mathcal O_F$. The determinant construction
above gives on that fiber a map of determinant one carrying the source
to the target lattice. Their direct sum $g_t$ carries $M_0$ onto
$M_1$ and commutes with $\rho d_t$, which acts by a scalar on
each fiber.

At the selected slope write $x=z_1+t_0z_2$ and let $S$ act on
divisors by $S\mathfrak d=\gamma^{-1}\mathfrak d$. The divisor
condition on the input is
$\operatorname{div}\alpha=(S-\kappa^j)\mathfrak d$, with
$\mathfrak d=\operatorname{div}(\mathfrak a)$ here denoting the
divisor of the fractional ideal. Put
$Q(Z)=P(Z)/(Z-\kappa^j)=\sum_i b_i Z^i$ for $0\le j<M$, and set
\[
 \beta=\prod_i(\gamma^{-i}\alpha)^{b_i}.
\]
Equivariance of principal divisors gives
$\operatorname{div}\beta=Q(S)(S-\kappa^j)\mathfrak d
=P(S)\mathfrak d$.

Change coordinates over $L$ as follows. On a coordinate of weight
$x=\kappa^i\bmod q$ apply $\gamma^{-i}$, and on all other coordinates
apply the identity. This is well-defined since $\gamma$ acts faithfully
on $\zeta$ with exponent $\kappa$. Call the resulting map $U$.
In the new coordinates the selected slope acts by $\zeta$ on the whole
orbit of weight one, and by the original eigenvalues elsewhere.
The source lattice is unchanged. In the weight-one block the product
of the transformed target ideals has divisor
$\sum_i p_i S^i\mathfrak d=P(S)\mathfrak d$ by the defining pushforward
of $\eta$. This identity also handles collisions among the
$\kappa^i\bmod q$, since their powers of $S$ then coincide. In each
other block the divisor sum is zero. The weight-one block is nonempty
(it contains $z=(1,0)$). The prescribed-determinant construction
therefore gives a block diagonal map $g'$ from $M_0$ to $UM_1$
with determinant $\beta$ on that block and one on the others.
Now $g_0=U^{-1}g'U$ transports $M_0$ to $M_1$ over $L$ and commutes
with the selected slope. With $H=\operatorname{Aut}_{\mathcal O_L}(M_0)$
and $\sigma=g_0^{-1}\rho g_0$, both representations preserve $M_0$.
Also $k_t=g_0^{-1}g_t\in H$ satisfies
$k_t^{-1}(\sigma d_t)k_t=\rho d_t$, as required in the prism construction.

The scalar that the determinant must recover can be checked before
taking homology. Let $f$ be reduction followed by $v$ at a prime
where all conjugates of $\alpha$ are units. For such units put
$W(b)=\sum_{h=0}^{n-1}\kappa^{hj} f(\gamma^h b)$ in $R$, where
$n=[F:L]$. Since $\kappa^n=1$ in $R$, reindexing gives
$W(\gamma^{-i}\alpha)=\kappa^{ij}W(\alpha)$. Thus
\[
 W(\beta)=Q(\kappa^j)W(\alpha)
 =C_j f\left(\prod_{h=0}^{n-1}
                (\gamma^h\alpha)^{\kappa^{hj}}\right).
\]
To evaluate the cycle, exclude the finitely many primes dividing
$q$, the relative discriminant, or the denominators of the lattices,
maps and their inverses. At a remaining relative degree-one prime
the Galois extension splits completely over the underlying prime.
The Chinese remainder theorem identifies restriction of scalars
after reduction with the direct sum of the $n$ conjugate blocks,
with $\gamma^h$ acting on the entries in the $h$th evaluation.
If $M_0$ is projective over $\mathcal O_L$, choose a projective
complement to make it free and extend its automorphisms by the
identity. Such a complement exists since a torsion-free finite
module over a Dedekind domain is projective. The added trivial
representation contributes zero.

In the localized group, write $P_g$ for the homotopy from a
representation to its conjugate by $g^{-1}$. On cycles the prisms
for two successive conjugations add modulo boundaries: the two-edge
path and its composite bound the natural-transformation triangle
(the usual triangulation of the bar simplex times a two-simplex).
Applying this with $g_t=g_0 k_t$ shows that the image of $y$ is
homologous to $\sum_t n_t P_{g_t}(\rho d_t)_*z$.
Indeed the terms $P_{g_0}(\rho d_t)_*z$ add to
$P_{g_0}\rho_*\partial b$, homologous to
$\sigma_*b-\rho_*b$ by the prism identity. All these maps can now
be reduced at the chosen prime.

For a scalar slope acting by $w$ on a block, evaluation of its prism
is additive as a function of the commuting matrix $g$ by the same
triangle. It factors through the determinant: over the residue field
of size greater than two the elementary matrices are commutators
(conjugate by a diagonal matrix with distinct entries in the two
relevant positions), and generate the special linear group by
elimination. On $g=\operatorname{diag}(b,1,\ldots,1)$ block
additivity and the scalar calculation give $x^jv(b)$ for $w=\zeta^x$;
the identity prism is null in homology. Thus the nonselected slopes
give zero because their block determinants are one. On the selected
slope use $U$ to read the reduced blocks in the coordinates of $g'$.
In the $h$th conjugate the surviving block has weight $\kappa^h$
and determinant the reduction of $\gamma^h\beta$. The evaluation
of $y$ is therefore $n_{t_0}W(\beta)$. Scaling the cycle by the
inverse of the unit $n_{t_0}$ gives the claimed determinant evaluation
$C_jW(\alpha)$ at every allowed degree-one prime.

We now count several divisor classes at once. Fix $m=2j+1$,
$j\ge1$, and put $N_*=4m+5$. Let $T_i$ be the order of the torsion
in $H_i(\mathrm{GL}_{N_*}(\mathcal O_L),\mathbb Z)$, for $i=m,m-1$.
These two numbers are independent of $q$. Consider a finite set of
inputs $\alpha$ satisfying the divisor condition. Write
$u_\alpha=\prod_{h=0}^{n-1}(\gamma^h\alpha)^{\kappa^{hj}}$, and put
$b=m-1+v_\ell(C_j)$. If the elements $u_\alpha^{\ell^b}$ are distinct
in $F^\times/F^{\times q}$, the number of inputs is at most
\[
                     q^{r_2(L)}T_mT_{m-1}.
\]

Indeed exclude the union of the exceptional primes for this finite
set. The power-detection argument supplies, for each distinct pair,
a relative degree-one prime where their quotient after raising to
$\ell^b$ is not a $q$th power. The residue group is cyclic and
contains $\mu_q$, so there exists a homomorphism $v$ to $R$ detecting
this quotient. Collect these finitely many evaluations in a single
map $E$ with target a product of copies of $R$. The evaluations
respect stabilization on positive-degree cycles. They thus define
a homomorphism on stable homology with $R$ coefficients, since bar
chains use finitely many matrices. By the determinant calculation,
the vectors $\ell^{m-1}E(y_\alpha)$ are pairwise distinct: dividing
$C_j$ by $\ell^{v_\ell(C_j)}$ gives a unit modulo $q$.

It remains to bound the image of $\ell^{m-1}E$. In the coefficient
exact sequence for stable group homology, the image of integral
homology modulo $q$ has quotient $H_{m-1}(\mathbb Z)[q]$, of order
at most $T_{m-1}$ by stability at $N_*$. On integral lifts use the
endomorphism $P=\prod_{h=2}^m(D_\ell-\ell^h)$ from the rank calculation.
Additivity under block sums gives
$E(Pz)=\prod_{h=2}^m(\ell-\ell^h)E(z)$ on reduction modulo $q$.
The coefficient on the right is $\ell^{m-1}$ times a unit modulo $q$.
The integral image of $P$ has rank at most $r_2(L)$, and its torsion
embeds in the stable torsion of order $T_m$. Its image in the target
of $E$ therefore has order at most $q^{r_2(L)}T_m$. Each coset
of the integral lifts contributes a translate of this image after
the same unit rescaling. Multiplying by the number of cosets proves
the estimate. In particular the loss in the power tested on the
arithmetic side is fixed before the cyclotomic layer $\nu$ is chosen.

To relate this count to the class group, write $F^\times$ additively
in the following calculation. Put $S=\gamma^{-1}$, $a=\kappa^j$ and
$\delta=S-a$. Form the group of pairs $(\alpha,\mathfrak d)$ with
$\operatorname{div}\alpha=\delta\mathfrak d$, modulo pairs of the form
$(\delta v,\operatorname{div}v)$, and denote it by $W_0$. There is
an exact sequence
\[
 0\longrightarrow \mathcal O_F^\times/\delta\mathcal O_F^\times
 \longrightarrow W_0\longrightarrow \ker(\delta:\mathrm{Cl}(F)
 \to\mathrm{Cl}(F))\longrightarrow0 .
\]
The last map takes the class of $\mathfrak d$. Its surjectivity is
the definition of a principal divisor; if $\mathfrak d=\operatorname{div}v$,
subtracting the relation leaves a unit in the first coordinate.
Two such units differ by a relation precisely when that relation
comes from a unit. Since the class group is finite, the cardinalities
of the kernel and cokernel of $\delta$ on it agree.

Let $T=\sum_{h=0}^{n-1}a^h\gamma^h$. Then
$T\delta=\gamma^{-1}(1-a^n)$, and $q\mid1-a^n$.
Consequently $(\alpha,\mathfrak d)\mapsto\ell^bT\alpha\bmod q F^\times$
is well defined on $W_0$. The preceding homological count bounds its
image. For the kernel estimate divide $F^\times$ by its roots of unity,
writing $Y=F^\times/\mu(F)$ additively, and use the same definition
of pairs with $Y$ in the first coordinate. Denote this quotient of
$W_0$ by $W$. The exact sequence above now has free units modulo
$\delta$ on the left. The target of the counting map becomes $Y/qY$,
so the preceding image bound still applies. Both $Y$ and the divisor
group are torsion-free, and $\delta$ is injective on each: $S^n=1$
and $a^n>1$. In particular $W$ identifies, via the first coordinate,
with a subgroup of $Y/\delta Y$.

If a class in the counting kernel is represented by $y$, then
$\ell^bTy=qz$ in $Y$. Applying $\delta$ and cancelling $q$ in $Y$
gives $\ell^b(1-a^n)Sy/q=\delta z$. Thus the $\ell$-primary kernel
is killed by $\ell^{b+v_\ell(j)}$, since
$v_\ell(a^n-1)=\nu+v_\ell(j)$ and $S$ is invertible.

Its number of generators also admits a bound independent of $\nu$.
If a pair $(y,D')$ represents an element of $W[\ell]$, then
$\ell y=\delta v$, $\ell D'=\operatorname{div}v$. The class of $v$
in $Y/\ell Y$ is well-defined by injectivity of $\delta$, and it
determines the pair class: if $v=\ell z$, the pair is the relation
$(\delta z,\operatorname{div}z)$. This homomorphism embeds $W[\ell]$
in the $S$-invariant classes represented by elements whose valuations
are all divisible by $\ell$, since $a\equiv1\pmod\ell$.

Here is a uniform count of the latter classes. In multiplicative
notation choose a generator $\rho$ of $\mu(F)$ and represent an
invariant class by $v$ with
$S(v)/v=w^\ell\rho^i$, $0\le i<\ell$, absorbing any remaining
power of $\rho^\ell$ into $w$. The pair $(i,\mathrm N_{F/L}w)$ has
at most $\ell^2$ values since $(\mathrm N w)^\ell=(\mathrm N\rho)^{-i}$.
For two representatives with the same pair, Hilbert's theorem 90
gives $w/w'=S(z)/z$. Consequently $v/(v'z^\ell)$ lies in $L^\times$.
Its valuations over $F$ are divisible by $\ell$; away from places
above $\ell$ this also holds over $L$, since $F/L$ is unramified
there. The number of such base classes modulo $L^{\times\ell}$ is
at most $\ell^{r_2(L)+s_L}|\mathrm{Cl}(L)[\ell]|$, where $s_L$ is
the number of primes above $\ell$. Indeed division of the divisor
away from these primes by $\ell$ maps into the $\ell$-torsion of the
localized class group, with kernel the localized units modulo
$\ell$th powers. The unit theorem gives the stated unit factor;
the localized class group is a quotient of the ordinary class group,
so its $\ell$-rank is no larger. As $s_L\le[L:\mathbb Q]$, we may
take $g=2[L:\mathbb Q]+2+\lfloor\log_2|\mathrm{Cl}(L)|\rfloor$ and
obtain $|W[\ell]|\le\ell^g$.

The elementary divisor theorem now bounds the $\ell$-part of $|W|$
by $q^{r_2(L)}T_mT_{m-1}\ell^{(b+v_\ell(j))g}$.
Divide by the $\ell$-part of the free unit index computed above
(the determinants of $S-a$ and $\gamma-a$ agree in absolute value).
Writing $A=\mathrm{Cl}(F)$, this yields
\[
 |A/(qA+\delta A)|\le
 T_mT_{m-1}\ell^{(b+v_\ell(j))g+c+v_\ell(j)} .
\]
Indeed the quotient on the left has order at most the $\ell$-part
of the cokernel of $\delta$ on $A$. The factor $q^{r_2(L)}$ has
cancelled. This bound is independent of the layer for fixed base,
$j$ and $c$, under the slope-separation hypothesis $2j+1<\ell^c$.

The uniformity in $\nu$ permits descent to the localization kernel.
Write $M=\mathbb Q_\ell/\mathbb Z_\ell(-j)$, $j\ge1$, and
$L_\infty=L(\mu_{\ell^\infty})$. Use continuous cohomology of the
maximal extension unramified outside $\ell$. A class in the
localization kernel restricts to zero at every place above $\ell$.
Restriction to $L_\infty$ is injective even on the whole $H^1$:
by inflation--restriction its kernel comes from
$H^1(\Gamma,M)$ with $\Gamma=\operatorname{Gal}(L_\infty/L)$.
The value on a topological generator determines a continuous
cocycle. Multiplication by $\kappa^{-j}-1$ is surjective on the
divisible group $M$, since it is a nonzero $\ell$-adic integer.
Thus a coboundary can be subtracted to kill this value, after which
continuity kills the cocycle on all of $\Gamma$.

Take any finite set of distinct classes in the localization kernel
and choose continuous cocycles for them. Their images are finite by
compactness. At each place above $\ell$ their restrictions are
coboundaries; choose an element of $M$ giving each such coboundary.
Take $\nu$ large enough that all values and all these finitely many
elements are killed by $\ell^\nu$. Over $F=L(\mu_{\ell^\nu})$ the
cocycles are now homomorphisms to $M[\ell^\nu]$, since the action
on their values is trivial. They vanish on all decomposition groups
above $\ell$: this holds at one prolongation of each base place by
the chosen coboundaries, and at their conjugates by the cocycle law
and triviality of the action on its finite image. They are unramified
elsewhere. By unramified abelian class field theory each therefore
factors through $A=\mathrm{Cl}(F)$. The resulting characters are
distinct since their restrictions to $L_\infty$ are already distinct.

Conjugating by a lift of $\gamma^{-1}$ multiplies these characters
by $\kappa^j$ on their values; the extra coboundary term vanishes
because the subgroup over $F$ acts trivially on the cocycle values.
Equivariance of the Artin map therefore makes them factor through
$A/(qA+\delta A)$. The order of this quotient bounds the size of
the original finite set, by the elementary divisor theorem. Since
the bound is independent of $\nu$, it bounds the cardinality of
the entire localization kernel.

One can first enlarge the base to attain $\ell^c>2j+1$. Indeed
restriction from $L$ to any intermediate cyclotomic layer is
injective by the same injectivity at $L_\infty$ and preserves
local triviality. If the original top root has exponent $c_0$,
take $c\ge c_0$ least with $\ell^c>2j+1$. The new relative
degree is $\ell^{c-c_0}$; when $c>c_0$ it is at most $2j+1$,
since $c_0\ge1$ and $\ell^{c-1}\le2j+1$.
In particular for $j=1$ and $\ell=3$, adjoining $\mu_9$ to a base
with $c_0=1$ suffices and has degree three.

The dependence on the base in this descent is quantitative.
Let $n_0=[L:\mathbb Q]$, $D_0=|\operatorname{disc}L|$, and write
$s$ for the degree of the enlargement. The polynomial
$X^s-\zeta_{\ell^{c_0}}$ generates it on choosing compatible primitive
roots; its degree is the extension degree. Its discriminant is a
unit times $s^s$, so the enlarged field has degree $d=sn_0$ and
absolute discriminant $D\le D_0^s s^{sn_0}$. For any such field
Minkowski's bound gives a representative of each ideal class of norm
at most $2^d\sqrt D$. Counting all ideals up to this bound by their
Dirichlet series at two gives
$|\mathrm{Cl}|\le4^dD\zeta(2)^d$, using the Euler product.
Also $\ell^{c-1}(\ell-1)\le d$, by the presence of the primitive
root. Thus $g\ll d+\log(2D)$ and
$b=2j+2jc+v_\ell(j!)+v_\ell(j!)$ at $M=2j+1$, $u=j$.

In particular there is an absolute $C$ such that for odd $\ell$ the
logarithm of the size of the localization kernel just bounded is at
most $(2D_0)^{(2jn_0)^C}$. Here are the size comparisons. For
$N_*=4(2j+1)+5$ and $p\le2j+2$ the explicit orbit estimate above
gives $\log A\ll (2jd)^{C_1}\log(2D)$ for an absolute $C_1$,
by taking logarithms of $I$ and $R$. The stabilizer bound has
$\log |S|\le dN_*^2\log3$. Inserting these in the torsion
inequality gives $\log T_i\le (2D)^{(2jd)^{C_2}}$ for both degrees.
The additional term $((b+v_\ell(j))g+c+v_\ell(j))\log\ell$
satisfies the same bound by the preceding class number and root
estimates. Finally $s\le2j+1$ and the discriminant inequality
absorb the change of base by increasing the absolute exponent $C$.

This also bounds the global second cohomology. Write $S$ for the
places above $\ell$ together with the archimedean places and
$G_{L,S}$ for the Galois group unramified outside $S$.
Poitou--Tate duality identifies the kernel of localization on
$H^2(G_{L,S},\mathbb Z_\ell(j+1))$ with the Pontryagin dual of the
localization kernel just bounded in degree one. The image embeds in
the sum of the local groups. At $w\mid\ell$, local Tate duality identifies
the latter group's dual with the invariants of
$\mathbb Q_\ell/\mathbb Z_\ell(-j)$. Since $L$ contains $\mu_\ell$
and $\ell$ is odd, the local cyclotomic image is $1+\ell^{c_w}\mathbb Z_\ell$,
where $c_w$ is maximal with $\mu_{\ell^{c_w}}\subset L_w$.
Indeed the closed subgroups of $1+\ell\mathbb Z_\ell$ are given by
the $\ell$-adic logarithm, and the local cyclotomic image is open.
Thus the local length is $c_w+v_\ell(j)$ and
$\ell^{c_w-1}(\ell-1)\le[L_w:\mathbb Q_\ell]\le n_0$.
There are at most $n_0$ such places, and the complex places contribute
zero. In particular $H^2$ is finite and its logarithmic cardinality
obeys the same bound $(2D_0)^{(2jn_0)^C}$ after increasing $C$.

For the small-conductor branch at $j=1$ one can pass directly to the
twisted degree-one kernel. Let $E_\ell$ be a finite coefficient
extension of $\mathbb Q_\ell$ containing the values of $\eta$, with
integer ring $O$ and ramification index $e_O$. Take
$L=\mathbb Q(\eta,\mu_\ell)$. Its degree divides
$\operatorname{ord}(\eta)(\ell-1)$ and is prime to $\ell$.
Restriction on degree-one cohomology with coefficients
$(E_\ell/O)(-1)\otimes\eta^{-1}$ is injective since corestriction
composed with restriction is multiplication by $[L:\mathbb Q]$.
Over $L$ the character twist disappears. Classes which vanish locally
at the primes dividing $\ell N$ and on inertia elsewhere restrict
to the untwisted localization kernel over $L$. Indeed they remain
locally trivial above $\ell$ and are unramified elsewhere, where the
untwisted coefficient module itself is unramified. Tensoring the
untwisted discrete module with $O$ commutes with these kernel
conditions by finite freeness over $\mathbb Z_\ell$. Consequently the
$O$-length of the twisted kernel divided by $e_O$ is at most the
$\mathbb Z_\ell$-length of the untwisted kernel: a cyclic summand of
length $h$ extends to $O/\ell^h O$, of length $e_O h$.

For $\ell=3,7$ the degrees of $L$ in the small-conductor branch are
bounded absolutely as shown above. Its conductor divides $\ell N$,
so the conductor-discriminant formula gives $D_0\le(\ell N)^{n_0}$.
The preceding bound therefore gives a normalized length, multiplied
by $\log\ell$, of at most $(\ell N)^{C'}$ with one absolute $C'$
for both primes, uniformly in the auxiliary primes and in the
coefficient extension. This supplies the exceptional small-conductor
estimate at twist one.

\begin{theorem}\label{thm:all-support}
Assume \eqref{eq:dirichlet-input}. For every $T>0$ and every $F\in\mathcal D_T$ one has
\[
 W_T(F)=\sum_{\rho}|\Phi_F(\rho)|^2\ge0 .
\]
The series converges, with multiplicities. In particular
$0\le\lambda^{\rm W}_T\le\lambda_{\rm A}(T)$, and
$I_\infty(\Gamma_F,\Gamma_F)\le0$ for every smooth primitive radial test.
\end{theorem}
\begin{proof}
Let $\rho$ be a nontrivial zero. Since $\rho\ne1$,
\eqref{eq:dirichlet-input} implies $\Re\rho\le1/2$. If $\Re\rho<1/2$, the factors relating $\xi$ to
$\zeta$ are finite and nonzero in the open critical strip. Hence
$\xi(1-\rho)=\xi(\rho)=0$ gives a zero of $\zeta$ with real part
$1-\Re\rho>1/2$, already excluded. Thus $\rho=1-\bar\rho$ for every
nontrivial zero, with no restriction on its ordinate.

Fix $T>0$. For real smooth tests the asserted identity follows from
\eqref{comparison} as in the preceding theorem. It also holds for complex
smooth tests $F=U+iV$: conjugate ordinates give opposite cross terms in
$|\Phi_U(\rho)+i\Phi_V(\rho)|^2$, and $W_T(F)=W_T(U)+W_T(V)$.
The rapid decay of both transforms justifies the pairing.

For $F\in\mathcal D_T$, choose smooth tests $F_n$ converging in the form
norm, using the core construction in the proof of Theorem~\ref{thm:KEP2}.
The vectors $(\Phi_{F_n}(\rho))_\rho$ are Cauchy in $\ell^2$, since the
squared distance between the $n$th and $m$th vectors is $W_T(F_n-F_m)$,
which tends to zero by form continuity. Their limit has coordinates
$\Phi_F(\rho)$: on the critical line Cauchy--Schwarz gives
$|\Phi_{F_n-F}(\rho)|\le\sqrt{2T}\,\|F_n-F\|_2$. Taking squared norms
gives the stated equality. The polar term vanishes on $\mathcal P_T$,
so the infima over complete and primitive unit vectors are ordered as
claimed. On smooth primitive tests \eqref{comparison} gives the
intersection inequality with the opposite sign.
\end{proof}

For example, take $R=5$ or $R=10$. Under \eqref{eq:dirichlet-input},
Theorem~\ref{thm:all-support} gives, for every $F\in\mathcal D_R$,
\[
 W_R(F)=\sum_{\gamma}
   \left|\int_{-R}^{R}F(t)e^{i\gamma t}\,dt\right|^2 .
\]
Here $R$ bounds the support of the test; the sum runs over all ordinates
with multiplicity at either radius. There is also an unconditional
consequence of the bounded-window certificate at both radii. Put
$r=2.8022$, take $G\in C_c^\infty(-r,r)$ and set $F(t)=G(t-h)$ with
$|h|\le R-r$. Then $F$ is supported in $(-R,R)$ and
$\Phi_F(s)=e^{(s-1/2)h}\Phi_G(s)$. In each term of the comparison
formula the two translation factors have product
$e^{(\rho-1/2)h}e^{(1/2-\rho)h}=1$. Thus
$W_R(F)=W_r(G)$, and \eqref{eq:power-window} applies with the same
norm since translation preserves $\|G\|_2$. For a sum of tests with
different translations the mixed terms carry different exponential
factors, so this argument gives no bound for them.

\begin{proposition}[the matrices in every degree]\label{prop:all-degree-gram}
With the arithmetic input \eqref{eq:dirichlet-input}, let $a_j>1$
and $u_j=a_j-\tfrac12$, $1\le j\le d$. Then, for $c_j\in\mathbb C$,
\[
 \sum_{j,k=1}^d c_j\overline{c_k}K(a_j,a_k)
 =\sum_{\gamma}\left|\sum_{j=1}^d\frac{c_j}{u_j-i\gamma}\right|^2 .
\]
Here $\tfrac12+i\gamma$ ranges over the nontrivial zeros with multiplicity.
In particular the hypothesis of Proposition~\ref{ladder} holds in every degree.
\end{proposition}
\begin{proof}
Write $X(z)=\xi(\tfrac12+z)$, where
$\xi(s)=\tfrac12 s(s-1)\pi^{-s/2}\Gamma(s/2)\zeta(s)$.
The logarithmic derivative $X'(u)/X(u)$ is exactly $L(\tfrac12+u)$.
By the specialization in Theorem~\ref{thm:all-support} all zeros of $X$ are imaginary.
Since $X$ is even, of order one and nonzero at the origin, its Hadamard
product is $X(0)\prod_{\gamma>0}(1+z^2/\gamma^2)$.
Indeed the exponential factors for $i\gamma$ and $-i\gamma$ cancel,
and evenness makes the remaining linear exponent constant. The zero count
$N(H)=O(H\log H)$ gives $\sum_{\gamma>0}\gamma^{-2}<\infty$ and justifies
logarithmic differentiation on compact sets away from the divisor.
Thus $L(\tfrac12+u)$ is the paired sum of $(u-i\gamma)^{-1}$ over signed
ordinates. For $v>0$ we may use $(v+i\gamma)^{-1}$ for the second paired
sum, by symmetry. The identity
\[
 \frac{(u-i\gamma)^{-1}+(v+i\gamma)^{-1}}{u+v}
 =\frac1{(u-i\gamma)(v+i\gamma)}
\]
now gives the asserted Gram matrix. The series of products converges
absolutely. Taking $a_j=32+j$ for $j=0,\ldots,m$ covers every $m$, so
Proposition~\ref{ladder} applies to every compact support radius.
\end{proof}

The involution in \cite{P43} supplies the adjoint in the radial
self-intersection. The Frobenius contacts in \cite{P44} fix its prime-power
coefficients, and the comparison theorem \cite[Theorem~C]{Paper45} identifies
the primitive self-intersection with the negative of the zero sum.
Together with the theorem above, these results complete the equivalence
between the intersection inequality for every compact support and the
classical Riemann hypothesis. The certificates of \S\ref{sec:log5} give explicit margins on
bounded windows. Using \eqref{eq:dirichlet-input},
Theorem~\ref{thm:all-support} gives the spectral identity on the full
domain for every radius, and Proposition~\ref{prop:all-degree-gram}
gives its matrix form in every degree.

\begin{thebibliography}{99}
\bibitem{BoasKac1945} R.~P.~Boas and M.~Kac, \emph{Inequalities for Fourier transforms of positive
functions}, Duke Math.~J. \textbf{12} (1945), 189--206.

\bibitem{Bombieri2000} E.~Bombieri, \emph{Remarks on Weil's quadratic functional in the theory of
prime numbers, I}, Atti Accad.~Naz.~Lincei Cl.~Sci.~Fis.~Mat.~Natur.~Rend.~Lincei (9) Mat.~Appl.
\textbf{11} (2000), no.~3, 183--233.

\bibitem{ConnesConsani2021} A.~Connes and C.~Consani, \emph{Weil positivity and trace formula, the archimedean place}, Selecta Math.\ (N.S.) \textbf{27} (2021), no.~4, Paper No.~77; arXiv:2006.13771.

\bibitem{CvS2025} A.~Connes and W.~D.~van Suijlekom, \emph{Quadratic forms, real zeros and echoes of the spectral action}, arXiv:2511.23257 (2025).

\bibitem{DLMF} NIST Digital Library of Mathematical Functions, \url{https://dlmf.nist.gov}, \S5.9.

\bibitem{Groskin2026} A.~Groskin, \emph{A finite Guinand--Weil dictionary and archimedean tail
order for the truncated Weil quadratic form}, arXiv:2607.02828 (2026).

\bibitem{HLP} G.~H.~Hardy, J.~E.~Littlewood and G.~P\'olya, \emph{Inequalities}, 2nd ed., Cambridge Univ.~Press, 1952.

\bibitem{Haynsworth1968} E.~V.~Haynsworth, \emph{Determination of the inertia of a partitioned Hermitian
matrix}, Linear Algebra Appl.\ \textbf{1} (1968), 73--81.

\bibitem{IK} H.~Iwaniec and E.~Kowalski, \emph{Analytic Number Theory}, Amer.~Math.~Soc.
Colloq.~Publ.~53, 2004.

\bibitem{Johansson2017} F.~Johansson, \emph{Arb: efficient arbitrary-precision midpoint-radius interval
arithmetic}, IEEE Trans.\ Comput.\ \textbf{66} (2017), 1281--1292.

\bibitem{Katznelson} Y.~Katznelson, \emph{An Introduction to Harmonic Analysis}, 3rd ed.,
Cambridge Univ.~Press, 2004.

\bibitem{Kim2026} T.~Kim, Y.~Hong, M.~Kim, S.~Choi, J.~Jang and M.~Kim, \emph{A numerical realization of Suzuki's Weil-quadratic-form operator: the archimedean spectral law, its universality, and an operator form of Weil's positivity criterion}, arXiv:2607.24830 (2026).

\bibitem{Krein1940} M.~G.~Kre\u{\i}n, \emph{Sur le probl\`eme du prolongement des fonctions
hermitiennes positives et continues}, Dokl.~Akad.~Nauk SSSR \textbf{26} (1940), 17--22.

\bibitem{Liu2026} V.~Liu, \emph{Certified Weil positivity beyond the unit window: source-exact
block-Schur and tail-compensation bounds for the Riemann zeta function}, manuscript dated 14 Sep 2026,
alphaXiv; under review at Math.\ Comp.\ (unrefereed; AI tools declared by the author).

\bibitem{MoriUpper} R.~Mori, \emph{Improved upper bounds for windowed Weil forms with exponential factor
$e^{-4\pi\mu}$}, Zenodo (2026), v3 doi:\nolinkurl{10.5281/zenodo.23363767}; concept doi:\nolinkurl{10.5281/zenodo.23175060}
(unrefereed).

\bibitem{P43} D.~A.~Trejo Pizzo,
\emph{A correspondence algebra with a canonical trace on the idele class
group, its Tate symbol and its numerical form}, preprint.

\bibitem{P44} D.~A.~Trejo Pizzo,
\emph{An arithmetic Frobenius family and its integral shadow on a
correspondence algebra of the idele class group}, preprint.

\bibitem{Paper45} D.~A.~Trejo Pizzo, \emph{An explicit formula and a comparison theorem for the canonical trace on a correspondence
algebra of the idele class group}, preprint, 2026.

\bibitem{PlattTrudgian2021} D.~J.~Platt and T.~S.~Trudgian, \emph{The Riemann hypothesis is true up to
$3\cdot10^{12}$}, Bull.\ London Math.\ Soc.\ \textbf{53} (2021), 792--797.

\bibitem{Rockafellar} R.~T.~Rockafellar, \emph{Convex Analysis}, Princeton Univ.~Press, 1970.

\bibitem{Rosser1941} J.~B.~Rosser, \emph{Explicit bounds for some functions of prime numbers}, Amer.\ J.\
Math.\ \textbf{63} (1941), 211--232.

\bibitem{RosserYoheSchoenfeld1968} J.~B.~Rosser, J.~M.~Yohe and L.~Schoenfeld, \emph{Rigorous computation and
the zeros of the Riemann zeta-function}, Proc.\ IFIP Congress 1968, 70--76.

\bibitem{SSV2012} R.~L.~Schilling, R.~Song and Z.~Vondra\v{c}ek, \emph{Bernstein Functions:
Theory and Applications}, 2nd ed., de Gruyter, 2012.

\bibitem{Sasvari1994} Z.~Sasv\'ari, \emph{Positive Definite and Definitizable Functions},
Akademie Verlag, Berlin, 1994.

\bibitem{SongVondracek2003} R.~Song and Z.~Vondra\v{c}ek, \emph{Potential theory of subordinate
killed Brownian motion in a domain}, Probab.~Theory Related Fields \textbf{125} (2003), 578--592.

\bibitem{Trefethen2013} L.~N.~Trefethen, \emph{Approximation Theory and Approximation Practice},
SIAM, Philadelphia, 2013.

\bibitem{Weil52} A.~Weil, \emph{Sur les ``formules explicites'' de la th\'eorie des nombres
premiers}, Comm.~S\'em.~Math.~Univ.~Lund (1952), 252--265.

\bibitem{Yoshida1992} H.~Yoshida, \emph{On Hermitian forms attached to zeta functions}, in:
\emph{Zeta Functions in Geometry}, Adv.~Stud.~Pure Math.~\textbf{21}, Kinokuniya, Tokyo, 1992, 281--325.

\bibitem{Zhu2026} X.~Zhu, \emph{Weil positivity in compact windows: a finite reduction, certified
two-sided bounds, and a Landau--Widom decay law}, arXiv:2608.24827 (2026).
\bibitem{CharneyStab} R. Charney, Homology stability for $GL_n$ of a Dedekind domain, Invent. Math. 56 (1980), 1--17.
\bibitem{WeibelK} C. Weibel, The $K$-book, Graduate Studies in Mathematics 145, AMS, 2013.
\end{thebibliography}
\end{document}
